% STRIKE 방법 상세 정리 (논문 작성용). XeLaTeX로 컴파일한다.
%   cd results/strike_paper && xelatex strike_methods_ko.tex && xelatex strike_methods_ko.tex
% 표와 그림은 scripts/strike_paper.py 와 scripts/strike_episode_viz.py 의 출력
% (docker_batch/run_strike_paper.sh -> outputs/strike_final/paper)을 읽는다. 다른 폴더를 쓰려면
%   xelatex "\def\paperdir{/경로/paper}% STRIKE 방법 상세 정리 (논문 작성용). XeLaTeX로 컴파일한다.
%   cd results/strike_paper && xelatex strike_methods_ko.tex && xelatex strike_methods_ko.tex
% 표와 그림은 scripts/strike_paper.py 와 scripts/strike_episode_viz.py 의 출력
% (docker_batch/run_strike_paper.sh -> outputs/strike_final/paper)을 읽는다. 다른 폴더를 쓰려면
%   xelatex "\def\paperdir{/경로/paper}% STRIKE 방법 상세 정리 (논문 작성용). XeLaTeX로 컴파일한다.
%   cd results/strike_paper && xelatex strike_methods_ko.tex && xelatex strike_methods_ko.tex
% 표와 그림은 scripts/strike_paper.py 와 scripts/strike_episode_viz.py 의 출력
% (docker_batch/run_strike_paper.sh -> outputs/strike_final/paper)을 읽는다. 다른 폴더를 쓰려면
%   xelatex "\def\paperdir{/경로/paper}% STRIKE 방법 상세 정리 (논문 작성용). XeLaTeX로 컴파일한다.
%   cd results/strike_paper && xelatex strike_methods_ko.tex && xelatex strike_methods_ko.tex
% 표와 그림은 scripts/strike_paper.py 와 scripts/strike_episode_viz.py 의 출력
% (docker_batch/run_strike_paper.sh -> outputs/strike_final/paper)을 읽는다. 다른 폴더를 쓰려면
%   xelatex "\def\paperdir{/경로/paper}\input{strike_methods_ko.tex}"
\providecommand{\paperdir}{../strike_reduced/final/paper}   % the report's results in this branch
\providecommand{\episodedir}{\paperdir/episodes}
\providecommand{\snapdir}{\paperdir/traj4/snapshots}   % filter snapshots
\documentclass[10pt,a4paper]{article}
\usepackage[margin=2.1cm]{geometry}
\usepackage{fontspec}
\setmainfont{Noto Serif CJK KR}
\setsansfont{Noto Sans CJK KR}
\setmonofont{Noto Sans Mono CJK KR}
\usepackage{amsmath,amssymb,bm,mathtools}
\usepackage{graphicx,float}
\usepackage{booktabs,array,tabularx,multirow}
\usepackage[dvipsnames]{xcolor}
\usepackage{enumitem}
\usepackage{caption,subcaption}
\usepackage{titlesec}
\usepackage{fancyhdr}
\usepackage[most]{tcolorbox}
\usepackage[ruled,vlined,linesnumbered]{algorithm2e}
\usepackage{fvextra}
\usepackage{tikz}
\usetikzlibrary{arrows.meta,positioning,fit,calc}
\usepackage[colorlinks=true,linkcolor=NavyBlue,urlcolor=NavyBlue,citecolor=NavyBlue]{hyperref}

\linespread{1.18}
\setcounter{tocdepth}{1}
\floatplacement{table}{H}
\setlength{\parindent}{1em}
\setlength{\parskip}{2pt}
\emergencystretch=2em
\setlist{nosep,leftmargin=1.5em}
\captionsetup{font=small,labelfont=bf,labelsep=period}
\renewcommand{\figurename}{그림}
\renewcommand{\tablename}{표}
\renewcommand{\contentsname}{목차}
\renewcommand{\refname}{참고 문헌}
\renewcommand{\abstractname}{요약}
\renewcommand{\appendixname}{부록}
\SetAlgorithmName{알고리즘}{알고리즘}{알고리즘 목록}
\SetKwInput{KwIn}{입력}
\SetKwInput{KwOut}{출력}
\SetKwInput{KwMem}{기억}
\SetKw{KwRet}{반환}
\SetKwIF{If}{ElseIf}{Else}{만약}{이면}{그렇지 않고 만약}{그렇지 않으면}{끝}
\SetKwFor{For}{각}{에 대해}{끝}
\titleformat{\section}{\large\bfseries\sffamily\color{NavyBlue}}{\thesection}{0.6em}{}
\titleformat{\subsection}{\normalsize\bfseries\sffamily}{\thesubsection}{0.6em}{}
\titleformat{\paragraph}[runin]{\normalsize\bfseries}{}{}{}[.]
\pagestyle{fancy}\fancyhf{}
\lhead{\small STRIKE: 방법과 결과 상세 정리}
\rhead{\small 2026.\ 10.\ 23}
\cfoot{\small\thepage}
\renewcommand{\headrulewidth}{0.3pt}

\newtcolorbox{keybox}[1]{colback=NavyBlue!4,colframe=NavyBlue!55,fonttitle=\bfseries\sffamily\small,title={#1},
  left=4pt,right=4pt,top=2pt,bottom=2pt,fontupper=\small}
\newtcolorbox{warnbox}[1]{colback=Orange!6,colframe=Orange!70,fonttitle=\bfseries\sffamily\small,title={#1},
  left=4pt,right=4pt,top=2pt,bottom=2pt,fontupper=\small}

\newcommand{\R}{\mathbb{R}}
\newcommand{\sat}{\operatorname{sat}}
\newcommand{\clip}{\operatorname{clip}}
\newcommand{\sdf}{\operatorname{sdf}}
\newcommand{\nrm}{\operatorname{normalize}}
\newcommand{\Safe}{\mathcal{S}}
\newcommand{\Base}{\bar{\mathcal{B}}}
\newcommand{\BaseP}{\mathcal{B}_{\mathrm{I}}}
\newcommand{\Term}{\mathcal{T}}
\newcommand{\unom}{u_{\mathrm{nom}}}
\newcommand{\upur}{u_{\mathrm{pur}}}
\newcommand{\uesc}{u_{\mathrm{esc}}}
\newcommand{\ubase}{u_{\mathrm{base}}}
\newcommand{\ustar}{u^{\star}}
\newcommand{\file}[1]{\texttt{\small #1}}
\newcommand{\figpath}[1]{\paperdir/figures/#1}
\newcommand{\tabin}[1]{\IfFileExists{\paperdir/tables/#1.tex}{\input{\paperdir/tables/#1.tex}}{\begin{center}\fbox{표 파일 없음: #1.tex}\end{center}}}
\newcommand{\figin}[2][\linewidth]{\IfFileExists{#2}{\includegraphics[width=#1]{#2}}{\fbox{그림 파일 없음}}}

\title{\vspace{-1.0cm}\textbf{\sffamily STRIKE: 다중 회복 정책 안전 필터와\\ 필터 포함 임무 학습(PS2)을 이용한 쿼드로터 저고도 표적 접근}\\[4pt]
\large 방법, 수식, 알고리즘, 평가 방법과 결과의 상세 정리 (논문 작성용)}
\author{김중원 \quad (정리 보조: Claude)}
\date{\small 2026년 10월 23일. 대상: 희소 경기장, 원래 추격자(반경 0.4 m, 2 m/s, 2 m/s$^2$)로 학습한 최종 구성}

\begin{document}
\maketitle
\thispagestyle{fancy}

\begin{abstract}
\noindent
장애물과 이동 추격자가 함께 있는 저고도 경기장에서 쿼드로터가 이동 표적의 0.6\,m 이내로 접근하는 문제(STRIKE)를 다룬다.
방법은 두 층으로 이루어진다.
\emph{안전 필터}는 임무 정책의 명목 입력을 일곱 개 회복 정책과 다섯 단계로 혼합한 35개 후보 각각을 한 스텝 실행한 뒤,
그 다음 상태에서 회복 정책을 3\,s 동안 가상으로 전개하여 바닥·장애물·벽·천장 제약과 상대 접근 속도를 반영한 추격자 여유도 $h_c$를 확인하고,
종단 상태가 호버하거나 제동할 수 있는 상태인지 검사한다. 통과한 후보 가운데 명목 입력을 가장 많이 유지하는 입력을 실행하고,
통과 후보가 없으면 직전에 인증한 계획을 그 추격자 예측 그대로 이어서 실행한다.
\emph{임무 정책}은 6\,m/s 기준 접근 제어에 잔차를 더하는 SAC 정책이며, 실제 필터를 포함한 환경에서 필터의 수정량을 벌점으로 받으며 재학습한다(PS2).
1000개 평가 장면과 15\,s 제한에서 PS2 정책과 최종 필터의 조합은 성공률 97.6\,\%, 접촉 2000회 중 1회(0.05\,\%, 95\,\% 상한 0.28\,\%)를 보였고,
같은 필터를 쓰는 잔차 정책(94.4\,\%, 접촉 0/2000)보다 성공률이 높았다(같은 장면 쌍 비교, 정확 McNemar 검정 $p<10^{-6}$).
수치는 현재 업로드된 평가(시드 0 전체, 시드 1 일부)에 기반하며, 세 시드 평가가 끝나면 같은 스크립트로 표와 그림이 갱신된다.
\end{abstract}

\tableofcontents
\vspace{4pt}

%======================================================================
\section{이 문서의 범위와 읽는 법}\label{sec:scope}

이 문서는 학술대회 원고(\emph{STRIKE 다중 회복 정책과 안전 필터 학습을 이용한 쿼드로터의 저고도 표적 접근}) 작성을 위해,
최종 구성의 수식과 알고리즘을 구현(\file{ps2rl/})과 대조하여 정리한 것이다.
모든 수식과 상수는 저장소의 코드와 학습에 사용된 설정 파일(\file{configs.json})에서 확인했다.
원고 초안과 다른 부분은 부록~\ref{app:draft}에, 주요 수식의 한글(HWP) 수식 편집기 입력은 부록~\ref{app:hwp}에 정리했다.

\begin{keybox}{이 버전의 구성 (원고의 ``최종 구성'')}
\begin{itemize}
\item \textbf{경기장}: 40\,m$\times$20\,m, 높은 장애물 3개, 이동 추격자 1개(반경 $r_c=0.4$\,m, 속력 $\le v_c=2$\,m/s, 가속도 $\le a_c=2$\,m/s$^2$), 이동 표적, 15\,s 제한.
\item \textbf{안전 필터}: 이산 다중 회복 필터. 회복 정책 7개(학습 회복 정책 또는 호버 LQR과 회피 반사의 결합), 혼합률 5단계, 3\,s 예측, 신규 계획 여유, 확장 종단 집합, 추격자 예측을 재생하는 계획 유지.
\item \textbf{임무 정책}: 6\,m/s 기준 접근 제어 + SAC 잔차(관측 135차원). PS2: 배치하는 필터를 그대로 포함한 환경에서 300만 스텝 재학습.
\item \textbf{평가}: 세 시드(전체 파이프라인을 시드별로 학습), 1000개 평가 장면(시드 9001), 15\,s 에피소드와 10\,s 절단 결과, Clopper--Pearson 구간과 같은 장면 쌍 비교.
\end{itemize}
\end{keybox}

\begin{figure}[H]\centering
\begin{tikzpicture}[node distance=5mm and 7mm, every node/.style={font=\small},
  box/.style={draw, rounded corners, align=center, minimum height=10mm, inner sep=4pt, fill=#1!8}, box/.default=gray,
  >={Stealth[length=2mm]}]
\node[box=blue] (p1) {Phase I\\ 학습 회복 정책 $\pi_\theta$\\ (안전 도착, TD3)};
\node[box=orange, right=of p1] (fam) {회복 정책군 $\{\pi_b^{(j)}\}_{j=0}^{6}$\\ 기본 회복 + 회피 반사\\ (속력비, 회피 각)};
\node[box=red, right=of fam] (filt) {이산 안전 필터 $F$\\ 35개 후보, 3\,s 전개,\\ 종단 검사, 계획 유지};
\node[box=gray, right=of filt] (plant) {쿼드로터,\\ 장애물, 추격자,\\ 표적};
\node[box=green, below=8mm of filt] (str) {임무 정책\\ $\unom=\sat(\upur+\Delta u_\xi)$};
\draw[->] (p1) -- (fam);
\draw[->] (fam) -- (filt);
\draw[->] (str) -- node[right]{$\unom$} (filt);
\draw[->] (filt) -- node[above]{$u$} (plant);
\draw[->] (plant) |- node[pos=0.75,below]{$o$ (상태, 장면)} (str);
\node[draw, dashed, rounded corners, fit=(filt)(str)(plant), inner sep=5pt,
  label={[font=\scriptsize]below:PS2: 필터를 포함한 환경에서 임무 정책을 학습($r^{\mathrm{PS2}}=r-w_f\|\eta\|^2$)}] {};
\end{tikzpicture}
\caption{전체 구성. 안전 필터는 회복 가능성을, 임무 정책은 표적 접근 성능을 담당한다.}\label{fig:overview}
\end{figure}

\paragraph{재현 경로} 표·그림·본문 수치는 \file{docker\_batch/run\_strike\_paper.sh}가 만든다.
\file{STEPS=tables}는 몇 초 안에 표(\file{.md/.csv/.tsv/.tex}), 그림 7종, 본문 인용 수치(\file{numbers.md})를 만든다.
나머지 단계는 대표 에피소드를 다시 실행해 2차원·3차원 궤적 그림, 상세 신호 그림, 2차원·3차원 GIF를 만든다(\S\ref{sec:results-episodes}).

%======================================================================
\section{문제 설정}\label{sec:problem}

\subsection{기체 모델}
기체 상태는 위치 $p\in\R^3$, 관성좌표 속도 $v\in\R^3$, 자세 쿼터니언 $q\in S^3$로 이루어지고,
입력은 질량으로 나눈 총추력(추력 가속도) $a_T$와 기체 각속도 명령 $\omega$이다.
\begin{equation}
x=(p,v,q),\qquad u=(a_T,\omega),\qquad 0\le a_T\le 4g=39.24\ \mathrm{m/s^2},\qquad |\omega_i|\le 18\ \mathrm{rad/s}.
\label{eq:state}
\end{equation}
시간 간격 $\Delta t=0.02$\,s(50\,Hz)의 명시적 Euler 적분으로 상태를 갱신한다. $R(q)$는 기체 좌표에서 관성좌표로의 회전행렬, $e_3=(0,0,1)^\top$, $\otimes$는 쿼터니언 곱이다.
\begin{align}
p_{k+1}&=p_k+\Delta t\,v_k, \label{eq:dyn-p}\\
v_{k+1}&=v_k+\Delta t\,\big(a_{T,k}\,R(q_k)e_3-g\,e_3\big), \label{eq:dyn-v}\\
\tilde q_{k+1}&=q_k+\tfrac{\Delta t}{2}\,q_k\otimes(0,\omega_k),\qquad q_{k+1}=\tilde q_{k+1}/\|\tilde q_{k+1}\|. \label{eq:dyn-q}
\end{align}
이를 $x_{k+1}=f_d(x_k,u_k)$로 쓴다. 입력은 항상 상자 $\mathcal U$로 포화시킨다($\sat$).
입력의 정규화 척도는 $u_{\mathrm{scale}}=(a_{\max},\omega_{\max},\omega_{\max},\omega_{\max})=(39.24,18,18,18)$이다.
이 모델은 PX4 오프보드 제어가 받는 추력--각속도 인터페이스에 맞춘 운동 모델이며, 모터 동특성, 지연, 공력, 추정 오차는 포함하지 않는다.

\subsection{경기장과 장면 생성}\label{sec:arena}
\begin{itemize}
\item \textbf{경기장}: $x\in[0,40]$\,m, $|y|\le10$\,m, 기준 바닥 고도 $z_f=1$\,m, 상한 고도 $z_{\max}=10$\,m. 기체 반경 $r_d=0.3$\,m.
\item \textbf{장애물}: 높은 기둥 3개. 확률 0.5로 원기둥(반경 $\mathcal U[0.8,2.0]$\,m), 아니면 직사각 기둥(반폭 $\mathcal U[1.5,3.0]\times\mathcal U[1.0,2.5]$\,m, 회전 $\pm30^\circ$). 중심 $x\in[7,31]$\,m.
  장애물 간 간격 3.5\,m 이상, 출발점과 표적에서 4\,m 이상, 옆벽에서 1\,m 이상 떨어지게 하고, 장애물과 옆벽에서 1.5\,m 넘게 떨어진 0.5\,m 격자점만으로 출발점과 표적을 이을 수 있는 장면만 사용한다.
\item \textbf{출발}: $x=2$\,m, $y\in[-3,3]$\,m, $z\in[2,4]$\,m, 속력 $\mathcal U[4,6]$\,m/s(표적 방향 $\pm15^\circ$), 기울기 $10^\circ$ 이하.
\item \textbf{표적}: $x\in[32,37]$\,m, $|y|\le5$\,m, $z\in[1.6,2.5]$\,m에서 시작해 수평으로 $\mathcal U[0.2,0.8]$\,m/s의 일정 속도로 떠돌며, $x\in[32,39]$, $|y|\le9$의 영역 경계에서 반사된다.
\item \textbf{추격자}: 출발점--표적 직선의 $35$--$85\,\%$ 지점에서 옆으로 $5$--$9$\,m 떨어진 곳(높이 $2$--$4$\,m)에서 시작한다.
\end{itemize}

\paragraph{추격자 운동(환경)} 추격자 $j$의 위치 $c_j$, 속도 $\nu_j$는 장애물을 밀어내는 순수 추적으로 움직인다.
\begin{align}
w&=\frac{p-c_j}{\|p-c_j\|}+2\varpi_c\begin{bmatrix}\nabla\sdf_O(c_{j,xy})\\0\end{bmatrix},\qquad
\varpi_c=\clip\!\Big(\frac{d_{\mathrm{rep}}-\sdf_O(c_{j,xy})}{d_{\mathrm{rep}}},0,1\Big),\quad d_{\mathrm{rep}}=2\ \mathrm m,\\
\nu_{j,k+1}&=\Pi_{v_c}\!\Big(\nu_{j,k}+\Delta t\,\Pi_{a_c}\big(\tfrac{v_c w/\|w\|-\nu_{j,k}}{\Delta t}\big)\Big),\qquad c_{j,k+1}=c_{j,k}+\Delta t\,\nu_{j,k+1},
\end{align}
여기서 $\Pi_r(\cdot)$는 노름을 $r$ 이하로 줄이는 연산이며, 높이는 $[z_f+r_c,z_{\max}]$로 제한한다.
필터는 이 운동을 알지 못하고, 거리 감소에 가장 불리한 \emph{가속도 제한 없는 순수 추적}(식~\eqref{eq:pursue})으로 예측한다.

\subsection{장애물의 부호 거리}
장애물 $i$는 중심 $c_i$, 반폭 $b_i=(b_x,b_y)$, 회전 $\psi_i$, 모서리 반경 $\rho_i$로 표현한다(원기둥은 $b_i=0$, $\rho_i$=반경, 직사각 기둥은 $\rho_i=0$).
$\ell=R(\psi_i)^\top(p_{xy}-c_i)$로 두면
\begin{equation}
\sdf_i(p_{xy})=\big\|\max(|\ell|-b_i,0)\big\|_2+\min\!\big(\max(|\ell_x|-b_x,\ |\ell_y|-b_y),\,0\big)-\rho_i,\qquad
\sdf_O=\min_i\sdf_i ,
\label{eq:sdf}
\end{equation}
이며 $|\cdot|$와 $\max$는 성분별로 적용한다. 밖에서 양수, 표면에서 0, 안에서 음수이다.

\subsection{종료 조건과 성과 지표}\label{sec:outcomes}
매 스텝 다음 사건을 검사한다(같은 스텝에서 접촉이 우선).
\begin{itemize}
\item \textbf{성공}(strike): $\|p-p_t\|<0.6$\,m.
\item \textbf{접촉}(contact): 장애물 $\sdf_O(p_{xy})<r_d$, 바닥 $z<z_f$, 추격자 $\|p-c_j\|<r_c+r_d=0.7$\,m 중 하나.
\item \textbf{영역 이탈}(out): $x\notin[-1,41]$, $|y|>11$ 또는 $z>z_{\max}$. 임무 실패이지만 접촉으로 세지 않는다.
\item \textbf{시간초과}(time-up): 15\,s(750스텝) 동안 위 사건이 없음. 10\,s 결과는 같은 에피소드를 10\,s에서 절단한 것이다.
\end{itemize}

%======================================================================
\section{안전 집합}\label{sec:safe}

\subsection{기하 제약}
\begin{equation}
h_f=z-z_f,\qquad h_o=\sdf_O(p_{xy})-r_d,\qquad h_w=d_{\mathrm{wall}}(p_{xy})-r_d,\qquad h_z=z_{\max}-z-r_d,
\label{eq:geo}
\end{equation}
여기서 $d_{\mathrm{wall}}$은 네 외벽($x=0$, $x=40$, $y=\pm10$)까지의 최소 거리다.
바닥 조건은 과제의 고도 정의($z\ge z_f$)를 그대로 따르므로 $r_d$를 빼지 않는다.

\subsection{추격자 여유도 $h_c$의 유도}\label{sec:hc}
추격자 $j$에 대해 거리와 시선(추격자에서 기체로 향하는 단위벡터)을
\begin{equation}
R_j=\|p-c_j\|,\qquad \hat e_j=\frac{p-c_j}{R_j}
\end{equation}
로 두면 기체가 거리를 벌리는 속도는 $\hat e_j^\top v$이다. 추격자는 거리를 최대 $v_c$로 줄이므로(기체를 정면으로 향할 때) 최악의 접근 속도는
\begin{equation}
s_j=v_c-\hat e_j^\top v,\qquad s_j^+=\max(s_j,0)
\label{eq:closing}
\end{equation}
이다. 기체가 반응 지연 $\delta$ 뒤에 유효 이탈 가속도 $a_e$로 거리를 벌리는 속도를 높인다고 하자.
지연 동안에는 거리를 최대 $s_j^+\delta$ 잃는다. 그 뒤 접근 속도가 $s(t)=s_j^+-a_e t$로 줄어 0이 될 때까지 잃는 거리는
\begin{equation}
\int_0^{s_j^+/a_e}(s_j^+-a_e t)\,dt=\frac{(s_j^+)^2}{2a_e}
\qquad\Big(\text{등가로}\ 0=(s_j^+)^2-2a_e d_{\mathrm{brake}}\Big)
\end{equation}
이다. 충돌 반경과 두 손실을 현재 거리에서 빼면 추격자 여유도를 얻는다.
\begin{equation}
\boxed{\;h_{c,j}=R_j-(r_c+r_d)-s_j^+\,\delta-\frac{(s_j^+)^2}{2a_e}\;},\qquad \delta=0.2\ \mathrm s,\quad a_e=5\ \mathrm{m/s^2}.
\label{eq:hc}
\end{equation}
$a_e$는 기체의 최대 가속도보다 훨씬 작게 잡은 값이다. 요구 거리 $d_{\mathrm{req}}=r_c+r_d+s^+\delta+(s^+)^2/(2a_e)$의 예는 다음과 같다.

\begin{center}\small
\begin{tabular}{lccl}\toprule
상황 & $s^+$ [m/s] & $d_{\mathrm{req}}$ [m] & 구성 \\\midrule
추격자로부터 $v_c$ 이상으로 멀어짐 & 0 & 0.7 & 충돌 반경만 \\
정지(호버) & 2 & 1.5 & $0.7+0.4+0.4$ \\
추격자 쪽으로 6\,m/s 정면 접근 & 8 & 8.7 & $0.7+1.6+6.4$ \\\bottomrule
\end{tabular}
\end{center}

\noindent
같은 거리 4\,m라도 정지 중이면 $h_c=2.5$\,m이지만 6\,m/s로 정면 접근 중이면 $h_c=-4.7$\,m이다.
즉 아직 충돌하지 않은 상태도 안전 집합 밖일 수 있고, 이것이 거리만 보는 조건과의 차이다.

\paragraph{열린 공간에서의 불변성(근사)} 지연 뒤 기체가 $\hat e_j$ 방향으로 $a_e$의 가속을 유지하고 추격자가 정면으로 $v_c$로 다가오면 $\dot R_j=-s_j$, $\dot s_j=-a_e$이므로 $s_j>0$에서
\begin{equation}
\dot h_{c,j}=\dot R_j-\dot s_j\,\delta-\frac{s_j\dot s_j}{a_e}=-s_j+a_e\delta+s_j=a_e\delta\ \ge0
\end{equation}
이다. 따라서 막힌 곳이 없으면 $\{h_c\ge0\}$은 ``도피''에 대해 유지된다. 시선 회전에 의한 항(크기 $\lesssim v_c^2/(4R)$)은 $a_e$를 실제 최대보다 작게 잡아 흡수한다고 본다.
장애물과 벽이 도피 방향을 막을 수 있다는 점이 이 식만으로는 부족한 이유이고, 필터가 회복 궤적을 직접 검사하는 이유이다.

\subsection{안전 집합}
장면 정보 $\zeta$(장애물과 추격자 상태)에 대해
\begin{equation}
\Safe(\zeta)=\big\{x:\ h_f\ge0,\ h_o\ge0,\ h_w\ge0,\ h_z\ge0,\ h_{c,j}\ge0\ \ \forall j\big\}.
\label{eq:safe}
\end{equation}

\subsection{비교: 구형 근사(초기 방식)}\label{sec:ball}
초기 방식은 추격자를 현재 위치에 고정된 반경 $R_{\mathrm{ball}}=r_c+r_d+v_cT=3.7$\,m($T=1.5$\,s)의 구로 보고 $\|p-c_j\|\ge R_{\mathrm{ball}}$을 요구한다.
$T$초 안에 추격자가 닿을 수 있는 영역을 덮지만, 그 뒤 추격자의 위치와 기체의 접근 속도를 반영하지 못한다.
Phase I 학습(\S\ref{sec:phase1})은 이 구형 근사를 사용하며, 구형 근사 필터는 비교 실험(표~\ref{tab:table1_15s}, \ref{tab:table3_ablation})에 쓰인다.

%======================================================================
\section{기준 집합: 호버 LQR}\label{sec:base}

\subsection{선형화와 LQR}
호버 목표점 $a$ 주변에서 오차 상태 $e=(p-a,\ v,\ \phi)\in\R^9$를 쓴다. $\phi=2\,\mathrm{sgn}(q_w)\,q_{xyz}$는 작은 회전벡터(자세 좌표)이다.
호버 입력 $\ustar=(g,0,0,0)$ 주변의 연속시간 선형 모델은
\begin{equation}
\dot p=v,\qquad \dot v_x=g\,\phi_y,\qquad \dot v_y=-g\,\phi_x,\qquad \dot v_z=\Delta a_T,\qquad \dot\phi=\omega
\end{equation}
이고, Euler 이산화 $A_d=I+\Delta tA$, $B_d=\Delta tB$에 대해 이산 LQR을 설계한다.
\begin{equation}
e_{k+1}=A_de_k+B_d\Delta u_k,\qquad \Delta u_k=-Ke_k,\qquad K=(R_u+B_d^\top PB_d)^{-1}B_d^\top PA_d,\qquad V(e)=e^\top Pe.
\label{eq:lqr}
\end{equation}
$P$는 이산 Riccati 방정식의 해이다. 가중치는 $Q=\mathrm{diag}(2,2,8,\ 0.01,0.01,0.5,\ 8,8,2)$, $R_u=\mathrm{diag}(0.06,0.15,0.15,0.12)$이다.

\subsection{정지점의 재설정}
속도가 있는 기체는 제자리에 멈출 수 없으므로, $V$를 최소로 만드는 목표점을 해석적으로 구한다.
$r=(v,\phi)$로 두고 $P$를 위치와 나머지로 분할하면
\begin{equation}
V=(p-a)^\top P_{pp}(p-a)+2(p-a)^\top P_{pr}r+r^\top P_{rr}r .
\end{equation}
$\partial V/\partial a=-2P_{pp}(p-a)-2P_{pr}r=0$에서
\begin{equation}
a^\star(x)=p+Dr,\quad D=P_{pp}^{-1}P_{pr};\qquad
V_{\min}(r)=\min_aV=r^\top Sr,\quad S=P_{rr}-P_{rp}P_{pp}^{-1}P_{pr}
\label{eq:anchor}
\end{equation}
이다($S$는 Schur 여행렬). 이때 $e=(-Dr,r)$이므로 재설정 LQR은 $u=\ustar-\tilde Kr$, $\tilde K=K_r-K_pD$로, 어디서든 목표점 없이 정지하는 제어기가 된다.
설계값에서 $D$의 속도 성분은 약 0.385\,s로, 수평 속도 1\,m/s마다 정지점이 0.385\,m 앞에 놓인다.

\subsection{호버 타원체와 기준 집합}
$\{e:e^\top Pe\le c\}$에서 위치 성분의 방향 $n$으로의 최대 변위는 Cauchy--Schwarz 부등식으로 $\sqrt{c\,n^\top[P^{-1}]_{pp}n}$이다. 따라서
\begin{equation}
\rho_{xy}=\sqrt{c_B\,\lambda_{\max}\big([P^{-1}]_{xy}\big)}=1.22\ \mathrm m,\qquad
z_{\mathrm{ext}}=\sqrt{c_B\,[P^{-1}]_{zz}}=0.63\ \mathrm m,\qquad c_B=50 .
\label{eq:extent}
\end{equation}
필터가 쓰는 기준 집합은 정적 환경의 호버 회복 가능 상태로, 추격자 조건을 포함하지 않는다(추격자는 회피 반사와 종단 조건으로 다룬다, \S\ref{sec:family}, \ref{sec:terminal}).
\begin{equation}
\Base=\left\{x\ \middle|\
\begin{aligned}
&h_f,h_o,h_w,h_z\ge0,\qquad r^\top Sr\le c_B,\\
&\sdf_O(a^\star_{xy})\ge\rho_{xy}+r_d,\qquad d_{\mathrm{wall}}(a^\star_{xy})\ge\rho_{xy}+r_d,\\
&a^\star_z-z_{\mathrm{ext}}\ge z_f,\qquad z_{\max}-a^\star_z\ge\rho_{xy}+r_d
\end{aligned}\right\}.
\label{eq:base}
\end{equation}
즉 충분히 느리고($V_{\min}\le c_B$), LQR이 정착하는 동안 움직일 수 있는 타원체 전체가 장애물·바닥·벽에서 떨어져 있다.
천장 조건은 구현대로 $\rho_{xy}$를 쓰며, $\rho_{xy}\ge z_{\mathrm{ext}}$이므로 보수적이다.
Phase I의 기준 집합 $\BaseP$는 벽과 천장 대신 추격자 구 조건 $\|a^\star-c_j\|\ge R_{\mathrm{ball}}+\rho$($\rho=1.22$\,m, 3차원 범위)를 쓴다.

\subsection{$c_B$의 인증}
목표점을 $a_0$에 고정한 LQR $u=\ustar-K(p-a_0,r)$이 $\{V\le c_B\}$를 유지하려면 $c_B$가 다음 세 한계보다 작아야 한다.
\begin{itemize}
\item 입력 한계: 타원체 위에서 $\max|K_ie|=\sqrt{c\,K_iP^{-1}K_i^\top}\le m_i$이므로 $c_U=\min_i m_i^2/(K_iP^{-1}K_i^\top)=56.1$ ($m=(g,18,18,18)$).
\item 자세 좌표: $|\phi|<2$이면 좌표가 일대일이므로 $c_{\mathrm{chart}}=4/\lambda_{\max}([P^{-1}]_{\phi\phi})=102.1$.
\item 비선형 감소: 비선형 Euler 모델에서 $\{V=c\}$ 위 4000개 표본의 한 스텝 감소 $V^+/V<1$을 확인한다. $c=\min(c_U,c_{\mathrm{chart}})=56.1$에서 이미 모든 표본이 감소하므로(최악 방향 0.3\,m/s$^2$ 바람 포함) $c_{\mathrm{Lyap}}\ge56.1$이다. 감소하지 않으면 이분법으로 줄인다.
\end{itemize}
$c_B=50$에서 표본 최대 $V^+/V$는 0.980(바람 포함 0.986)이다. $\Base$ 안에서 수평 속력은 2.67\,m/s, 수직 속력은 2.12\,m/s, 기울기는 $52.7^\circ$ 이하이다.
표본 검사이므로 연속 상태공간 전체의 엄밀한 증명은 아니다.

\subsection{호버 latch}
계획이 $\Base$에 들어가고 회피 반사가 꺼져 있으면($\sigma=0$, \S\ref{sec:family}) 그때의 목표점 $a_0$을 고정하고 고정 목표점 LQR로 호버한다.
인증이 다루는 제어기가 바로 이 고정 목표점 LQR이기 때문이다. 매 스텝 재설정하면 목표점이 밀리미터 단위로 흔들려, 경계에서 들어온 상태가 다시 $\Base$ 밖으로 나가는 일이 생긴다.
고정된 호버는 $V(p-a_0,r)\le c_B$인 동안 종단 조건으로 인정된다.

%======================================================================
\section{Phase I: 안전 도착 회복 정책}\label{sec:phase1}

\paragraph{정책} LQR 입력에 잔차를 더한다.
\begin{equation}
\pi_\theta(x,\zeta)=\sat\big(u_{\mathrm{LQR}}(x)+\Delta u_\theta(o)\big),\qquad u_{\mathrm{LQR}}(x)=\ustar-\tilde Kr .
\label{eq:phase1-policy}
\end{equation}
관측 $o\in\R^{98}$은 $z-z_f$, $v$, $q$(8개), 가까운 순서의 장애물 슬롯 8개(최근접점 벡터 2, 중심 오프셋 2, 반폭 2, $\cos2\psi$, $\sin2\psi$, 모서리 반경, 존재 여부: 각 10개), 추격자 구 슬롯 2개(중심 오프셋 3, 반경 $R_{\mathrm{ball}}$, 존재 여부: 각 5개)로 구성된다.
슬롯 순서는 회복 전개를 시작할 때의 순서로 고정하여 전개 중 입력이 불연속으로 바뀌지 않게 한다. $\Safe$와 $\BaseP$의 소속은 $o$의 정확한 함수이다.

\paragraph{목적함수와 Bellman 관계} 에피소드는 표본 (상태, 장면)에서 시작한 회복 전개 하나이며, $\BaseP$ 도달(성공), $\Safe$ 이탈(실패), 112스텝(2.24\,s) 중 먼저 일어나는 사건에서 끝난다.
\begin{equation}
J(\theta)=\mathbb E\big[\beta^{\tau_B}\,\mathbf 1\{\tau_B<\tau_{\mathrm{out}}\}\big],\qquad \beta=0.99 .
\label{eq:sa}
\end{equation}
장면을 포함한 상태 $y=(x,\zeta)$와 결합 전이 $y^+=G(y,u)$에 대해 최적값은
\begin{equation}
W^\star(y)=1\ (y\in\BaseP\cap\Safe),\qquad W^\star(y)=0\ (y\notin\Safe),\qquad W^\star(y)=\beta\max_{u\in\mathcal U}W^\star(G(y,u))\ \ (\text{그 밖})
\end{equation}
를 만족하고, $W^\star(y)=\beta^{d(y)}$($d$는 안전하게 도착하는 최소 스텝 수)이다. TD3의 critic 목표는 이 관계를 그대로 쓴다.
\begin{equation}
y_{\mathrm{TD}}=b(o)+\beta\,c(o)\,\min_{i=1,2}Q'_i\big(o',\mu'(o')\big),\qquad b=\mathbf 1\{\BaseP\},\ c=\mathbf 1\{\Safe\setminus\BaseP\},
\end{equation}
여기서 $Q'(o',\cdot)$는 $o'\in\BaseP$이면 1, $o'\notin\Safe$이면 0으로 둔다. 엔트로피 항은 critic 목표에 넣지 않는다(도착 순서 ``빨리 $>$ 늦게 $>$ 실패''를 보존하기 위함).
actor 손실은 $-\mathbb E[\min_iQ_i(o,\Delta u_\theta(o))]+0.2\,\mathbb E\big[\tfrac14\|\Delta u/u_{\mathrm{scale}}\|^2\big]$이고, 둘째 항이 잔차를 LQR 쪽으로 붙잡는다. actor는 LQR 모방($\Delta u\approx0$)으로 초기화한다.

\paragraph{학습 분포} 일반(속력 $\le10$\,m/s, 기울기 $\le45^\circ$, 장애물 0--6개, 추격자 구 0--2개), 정지 경로 막힘(속력 $\ge3$\,m/s, LQR 정지점 부근에 장애물), 추격자 근접(속력 $\ge3$\,m/s, LQR 정지점 부근에 추격자 구), 호버 근방(느리고 정지점이 비어 있음) 영역을 가중치 $(0.4,0.2,0.2,0.2)$로 섞는다.
평가에는 저속 상태 근처에 추격자 구가 있는 저속 탈출 영역을 추가로 쓴다(가중치 0). 영역마다 5만 개 시작점을 미리 뽑고, 검증 1000개와 시험 2000개는 따로 뽑는다.

\paragraph{평가 지표와 결과} 시험 상태가 $C_N$에 속하는 비율 $\hat m$을 쓴다. $C_N$은 $N=T/\Delta t=75$스텝(1.5\,s) 안에 $\Safe$를 벗어나지 않고 $\BaseP$에 처음 도달하는 상태의 집합이다.
학습 회복 정책은 가중 평균 $70.1\pm0.3\,\%$(세 시드)로 호버 LQR의 $43.8\,\%$보다 높다(표~\ref{tab:table4_phase1}, 그림~\ref{fig:phase1}).
이 평가는 추격자를 정지한 구로 근사한 Phase I 조건이며, 최종 필터의 이동 추격자 성능과는 구분해야 한다.

%======================================================================
\section{회복 정책군}\label{sec:family}

\paragraph{기본 회복 입력} 기준 집합 안이거나 호버가 고정되어 있으면 고정 목표점 LQR, 밖이면 학습 회복 정책(또는 호버 LQR 변형에서는 재설정 LQR)을 쓴다.
\begin{equation}
\ubase^{(j)}(x)=\begin{cases}\ustar-K(p-a_0,r) & x\in\Base\ \text{또는 호버 고정},\\ \pi_\theta(x,\zeta_{\mathrm{view}}) & \text{그 밖(학습 변형)},\\ u_{\mathrm{LQR}}(x) & \text{그 밖(LQR 변형)}.\end{cases}
\end{equation}
호버가 고정되기 전에는 $a_0=a^\star(x)$이므로 첫 줄은 재설정 LQR과 같다. 학습 회복 정책의 장면 입력 $\zeta_{\mathrm{view}}$는 계획을 만든 시점의 장면(가까운 장애물 8개와 추격자 구)으로 고정한다. Phase I에서 학습한 범위(장애물 10\,m, 추격자 12\,m 이내)만 넣는다.

\paragraph{회피 반사} 최소 추격자 여유도 $h_{\min}=\min_jh_{c,j}$에 따라 회피 입력을 섞는다.
\begin{equation}
\sigma(x)=\clip\!\Big(\frac{h_{\mathrm{off}}-h_{\min}(x)}{h_{\mathrm{off}}-h_{\mathrm{on}}},0,1\Big),\qquad
\pi_b^{(j)}(x)=\sat\big((1-\sigma)\,\ubase^{(j)}(x)+\sigma\,\uesc^{(j)}(x)\big),\qquad h_{\mathrm{on}}=0.5,\ h_{\mathrm{off}}=2.0\ \mathrm m.
\label{eq:reflex}
\end{equation}
$h_{\min}\ge2$\,m이면 기본 회복 입력만, $h_{\min}=1.25$\,m이면 절반씩, $h_{\min}\le0.5$\,m이면 회피 입력만 쓴다. 반사가 작동하면($\sigma>0$) 호버 고정을 푼다.

\paragraph{회피 입력} 가장 위험한 추격자 $j^\star=\arg\min_jh_{c,j}$의 반대 방향 $\hat e_{xy}$를 변형 각 $\vartheta_j$만큼 돌리고, 정지점 $a^\star$ 주변의 장애물과 벽의 바깥 법선으로 민다.
\begin{align}
d&=\nrm\Big(R(\vartheta_j)\,\hat e_{xy}+2\varpi_o\,n_o+2\textstyle\sum_{i=1}^4\varpi_{w,i}\,n_{w,i}\Big),\qquad
\varpi=\clip\!\Big(\frac{d_{\mathrm{push}}-(\mathrm{dist}(a^\star_{xy})-\rho_{xy}-r_d)}{d_{\mathrm{push}}},0,1\Big),\\
v_{\mathrm{des}}&=\big(v_{\mathrm{esc},j}\,d,\ 1.5\,(z_{\mathrm{ref}}-z)\big),\qquad z_{\mathrm{ref}}=\clip(z,\ z_f+1,\ z_{\max}-3),
\end{align}
여기서 $n_o=\nabla\sdf_O(a^\star_{xy})/\|\nabla\sdf_O\|$, $n_{w,i}$는 벽의 안쪽 법선, $\mathrm{dist}$는 각 장애물·벽까지의 거리, $d_{\mathrm{push}}=1.5$\,m이다.
속도 루프는 $a=k_v(v_{\mathrm{des}}-v)$(수평 성분은 $\|a_{xy}\|\le g\tan45^\circ$로 제한), $a_{\mathrm{des}}=(a_{xy},a_z+g)$에서
$a_T=a_{\mathrm{des}}^\top z_b$, $\omega=k_{\mathrm{att}}R^\top(z_b\times\hat a_{\mathrm{des}})-k_\psi\psi\,e_3$로 입력을 만들고 포화시킨다($k_v,k_{\mathrm{att}},k_\psi)=(2,10,3)$, $z_b=R(q)e_3$, $\psi$는 요각).
추격자 여유가 충분할 때는 학습 회복 정책이 장애물을 처리하고, $h_c$가 줄어들수록 회복 정책은 ``$v_{\mathrm{esc}}$로 도피''로 바뀐다. 이는 열린 공간에서 $h_c$를 줄이지 않는 행동이다(\S\ref{sec:hc}).

\paragraph{일곱 변형} 회피 속력비 $v_{\mathrm{esc}}/v_c$와 회피 각 $\vartheta$의 조합으로 회복 정책을 늘린다(표~\ref{tab:variants}).
필터는 어느 하나라도 인증하는 후보를 받아들이므로, 회복 가능 집합은 각 변형 회복 가능 집합의 합집합이 된다.

\begin{table}[h]\centering\small
\caption{필터의 회복 정책 일곱 변형. 번호는 동률일 때의 우선순위이다.}\label{tab:variants}
\begin{tabular}{cccc}\toprule
번호 $j$ & 기본 회복 입력(기준 집합 밖) & $v_{\mathrm{esc}}/v_c$ & $\vartheta$ \\\midrule
0 & 학습 회복 정책 $\pi_\theta$ & 2 & $0^\circ$ \\
1 & 재설정 호버 LQR & 2 & $0^\circ$ \\
2, 3 & 학습 회복 정책 & 2 & $+60^\circ$, $-60^\circ$ \\
4 & 학습 회복 정책 & 3.5 & $0^\circ$ \\
5, 6 & 학습 회복 정책 & 3.5 & $+60^\circ$, $-60^\circ$ \\\bottomrule
\end{tabular}
\end{table}

%======================================================================
\section{이산 다중 회복 안전 필터}\label{sec:filter}



\subsection{후보 입력}
임무 정책의 명목 입력 $\unom$을 각 회복 정책의 현재 입력과 섞어 후보를 만든다.
\begin{equation}
u^{(j)}_\lambda=\sat\big(\lambda\,\unom+(1-\lambda)\,\pi_b^{(j)}(x)\big),\qquad \lambda\in\Lambda=\{1,0.75,0.5,0.25,0\},\quad j=0,\dots,6 .
\label{eq:cand}
\end{equation}
$7\times5=35$개 후보를 검사한다($\lambda=1$에서는 일곱 개가 같은 입력이다).
이는 입력 집합 전체를 최적화하는 것이 아니라 명목 입력과 각 회복 입력을 잇는 선분 위의 유한한 점을 검사하는 것이다.

\subsection{회복 궤적의 전개와 추격자 예측}
후보를 한 스텝 실행한 다음 상태에서 시작해 그 변형의 회복 정책을 $N=150$스텝(3\,s) 전개한다. 추격자는 예측 기체를 향해 $v_c$로 직진한다고 가정한다(거리 감소에 가장 불리한 운동).
\begin{align}
x_1&=f_d\big(x,u^{(j)}_\lambda\big),\qquad x_{i+1}=f_d\big(x_i,\ \pi_b^{(j)}(x_i;c_i)\big),\qquad i=1,\dots,N, \label{eq:rollout}\\
c_{\ell,i+1}&=c_{\ell,i}+\min\!\big(v_c\Delta t,\ \|p_i-c_{\ell,i}\|\big)\,\frac{p_i-c_{\ell,i}}{\|p_i-c_{\ell,i}\|},\qquad c_{\ell,0}=c_\ell(\text{현재}),\qquad c_{\ell,i,z}\ge z_f+r_c . \label{eq:pursue}
\end{align}
$\ell$은 추격자 번호이다. 회피 반사는 예측 추격자 위치 $c_i$를 보고, 학습 회복 정책의 장면 입력은 계획 시점의 것으로 고정된다. 종단 첨자는 $H=N+1$이다.

\subsection{궤적 조건(신규 계획 여유)}
새 계획은 모든 예측 노드 $i=1,\dots,H$에서 여유를 더한 안전 집합에 있어야 한다.
\begin{equation}
\min\{h_f,h_o,h_w,h_z\}(x_i)\ge m_o=0.15\ \mathrm m,\qquad \min_\ell h_{c,\ell}(x_i;c_{\ell,i})\ge m_c=0.2\ \mathrm m .
\label{eq:fresh}
\end{equation}
여유는 반사 방향의 재조준과 실제 추격자가 예측과 다르게 움직이는 것에 대비한 것이다. 계획 유지(\S\ref{sec:commit})의 재검사는 여유 없이($m_o=m_c=0$) 안전 집합만 요구한다.

\subsection{종단 집합}\label{sec:terminal}
종단 상태 $x_H$는 호버하거나 직선 제동으로 멈출 수 있어야 한다.
\begin{align}
\Term_{\mathrm{hover}}&=\big\{x_H:\ \text{호버 고정이면 } V(p_H-a_0,r_H)\le c_B,\ \text{아니면 } x_H\in\Base\big\},\\
\Term_{\mathrm{flee}}&=\big\{x_H:\ \sdf_O(p_{\mathrm{brake}})\ge r_d,\ d_{\mathrm{wall}}(p_{\mathrm{brake}})\ge r_d\big\},\qquad
p_{\mathrm{brake}}=p_{H,xy}+\frac{v_{H,xy}\|v_{H,xy}\|}{2a_e},\\
\Term&=\big(\Term_{\mathrm{hover}}\cup\Term_{\mathrm{flee}}\big)\cap\big\{\min_\ell h_{c,\ell}(x_H;c_{\ell,H})\ge m_c\big\}.
\label{eq:terminal}
\end{align}
$p_{\mathrm{brake}}$는 수평 속도 방향으로 $a_e$로 제동할 때의 정지점이다. 따라서 종단이 항상 $\Base$인 것은 아니며, 도피 중인 상태도 제동 공간이 남아 있으면 허용한다.
회복 정책 $j$의 유한시간 회복 가능 집합은
\begin{equation}
C^{(j)}=\big\{(x,\zeta):\ x_i\in\Safe^{+}\ (i=1,\dots,H),\ x_H\in\Term\big\},\qquad C=\bigcup_{j=0}^{6}C^{(j)}
\end{equation}
이며($\Safe^+$는 식~\eqref{eq:fresh}의 여유를 더한 안전 집합), 후보 실행 직후 상태가 $C^{(j)}$에 속하는지를 검사한다.
합집합은 각 집합을 포함하므로 같은 검사 조건에서 한 회복 정책만 쓰는 경우보다 선택지가 줄지 않는다. 다만 모든 회피 가능 상태를 찾는다는 뜻은 아니다.

\subsection{선택 규칙}
인증된 조합의 집합을 $\mathcal A$라 할 때 명목 입력의 비중이 가장 큰 후보를 실행한다.
\begin{equation}
(\lambda^\star,j^\star)\in\arg\max_{(\lambda,j)\in\mathcal A}\lambda,\qquad u=u^{(j^\star)}_{\lambda^\star}\quad(\text{동률이면 작은 }j).
\label{eq:select}
\end{equation}
변형마다 회복 입력이 다르므로 이는 $\|u-\unom\|$의 최소화와 일반적으로 같지 않다. ``명목 입력이 회복 가능성을 유지하면 그대로 통과시키고, 필요할 때만 회복 쪽으로 움직인다''는 최소 수정의 이산 형태이다.
그림~\ref{fig:snapshot}은 실제 에피소드에서 필터가 한 스텝에 검사한 내용을 다시 계산한 것이다. 명목 입력($\lambda=1$)은 앞의 장애물 때문에 모든 회복 정책에서 거부되고, 회복 정책 0과 $\lambda=0.75$의 조합이 인증되어 실행되었다(평가 기록과 일치).
\IfFileExists{\snapdir/snapshot_ep845_step091.png}{%
\begin{figure}[H]\centering
\includegraphics[width=\linewidth]{\snapdir/snapshot_ep845_step091.png}
\caption{필터의 한 스텝 결정(시드 0, 에피소드 845, $t=1.8$\,s). 왼쪽: 일곱 회복 정책의 3\,s 계획($\lambda=0$)과 끝나는 방식(점선 원: 호버 타원체, $\times$: 제동 정지점), 선택한 후보의 계획(굵은 파랑)과 추격자 예측(초록 점선), 추격자의 $h_c=0$ 거리(옅은 원). 오른쪽: 후보 35개의 인증 결과(초록: 인증, 빨강: 궤적 조건 위반, 주황: 종단 조건 위반)와 선택한 칸.}\label{fig:snapshot}
\end{figure}}{}

\subsection{계획 유지와 추격자 예측 재생}\label{sec:commit}
학습 회복 정책의 입력(가까운 장애물 순서, 범위 안의 장애물)은 스텝마다 바뀔 수 있어 직전에 인증한 계획이 다음 스텝에 달라질 수 있다.
그래서 필터는 기억 $m=(\zeta_{\mathrm{view}},a_0,\text{호버 고정 여부},j_c,\text{추격자 슬롯 순서},\hat c)$를 저장한다. $\hat c$는 인증한 계획이 예측한 다음 스텝의 추격자 위치이다.
$\mathcal A=\varnothing$이면 저장한 변형 $j_c$의 회복 입력을 \emph{저장한 장면 입력과 예측 추격자 위치} $\hat c$로 계산해 실행한다.
모델이 정확하면 이 입력은 인증된 회복 궤적의 다음 입력과 정확히 같다(committed continuation). 이 입력도 실제 추격자 위치에 대해 여유 없이 다시 검사하고, 실패하면 그대로 실행하되 ``인증 없음''으로 기록한다.
예측 추격자 위치는 식~\eqref{eq:pursue}로 한 스텝씩 진행한다.

\subsection{알고리즘}
\begin{algorithm}[H]
\caption{이산 다중 회복 안전 필터 $F$ (한 제어 스텝)}\label{alg:filter}
\KwIn{상태 $x$, 장면(장애물, 추격자 $c$), 명목 입력 $\unom$}
\KwMem{$m=(\zeta_{\mathrm{view}},a_0,\text{고정 여부},j_c,\text{슬롯 순서},\hat c)$}
\KwOut{실행 입력 $u$, 갱신된 기억 $m'$, 정보 $(\lambda^\star,\text{인증 여부})$}
현재 장면에서 새 장면 입력 $\zeta_{\mathrm{new}}$를 만들고, 추격자 예측 $c_1\gets$ 식~\eqref{eq:pursue}를 한 스텝 적용\;
\For{$j=0,\dots,6$}{
  $\pi_b^{(j)}(x)$ 계산(식~\eqref{eq:reflex})\;
  \For{$\lambda\in\{1,0.75,0.5,0.25,0\}$}{
    $u^{(j)}_\lambda\gets\sat(\lambda\unom+(1-\lambda)\pi_b^{(j)}(x))$, \ $x_1\gets f_d(x,u^{(j)}_\lambda)$\;
    회복 정책 $j$를 $N=150$스텝 전개(식~\eqref{eq:rollout}--\eqref{eq:pursue})\;
    모든 노드가 식~\eqref{eq:fresh}를 만족하고 $x_H\in\Term$이면 $(\lambda,j)$를 $\mathcal A$에 추가\;
  }
}
\eIf{$\mathcal A\neq\varnothing$}{
  $(\lambda^\star,j^\star)\gets\arg\max_{\mathcal A}\lambda$ (동률이면 작은 $j$), \ $u\gets u^{(j^\star)}_{\lambda^\star}$\;
  $m'\gets(\zeta_{\mathrm{new}},a_0',\text{고정 여부}',j^\star,\text{새 슬롯 순서},c_1)$\;
}{
  $u\gets\pi_b^{(j_c)}(x;\zeta_{\mathrm{view}},\hat c)$ \tcp*{직전 인증 계획의 다음 입력}
  여유 없이 실제 추격자에 대해 재검사하고, 실패하면 ``인증 없음''으로 기록\;
  $m'\gets(\zeta_{\mathrm{view}},a_0',\text{고정 여부}',j_c,\text{슬롯 순서},\Pi(\hat c))$\;
}
\KwRet{$u,\ m'$}
\end{algorithm}

\subsection{보장되는 것과 보장되지 않는 것}\label{sec:guarantee}
\begin{keybox}{명제 (가정 아래의 진술)}
가정: (A1) 기체 모델이 정확하고 외란이 없다. (A2) 각 추격자의 속력은 $v_c$ 이하이다. (A3) 3\,s 예측 구간에서 속력 $v_c$의 순수 추적이 $h_c$에 대해 최악의 경우이다(거리에 대해서는 성립, 시선 방향에 대해서는 근사).

이때 실행되는 모든 입력은 인증된 후보이거나 직전 인증 계획의 이어짐이다. 인증된 계획은 3\,s 동안 상태를 여유를 둔 $\Safe$ 안에 두고 $\Term$에서 끝난다.
추격자 예측 재생 덕분에 이어짐은 인증된 회복 궤적의 꼬리를 정확히 재현하며, 실제 추격자 위치에 대해 다시 검사된다.
\end{keybox}
남는 간극은 $\Term$이 추격자에 대해 불변 집합이 아니라는 점이다. 호버는 결국 추격자에게 따라잡히고, 도피는 열린 공간이 남아 있는 동안만 안전하다.
또한 이어짐을 재검사할 때 예측 구간이 한 스텝 늘어나므로, 그 마지막 상태가 $\Term$을 만족하지 못해 인증을 잃을 수 있다.
따라서 이 필터는 3\,s 앞을 인증하고 20\,ms마다 다시 계획하는 장치이며, 이동 추격자에 대한 무한시간 불변성의 증명은 아니다.
실제 측정에서는 인증 상실이 PS2 정책 에피소드의 1.5\,\%에서 나타났다(표~\ref{tab:table2_final}).

\subsection{계산량}
한 제어 스텝에서 후보 35개와 이어짐 7개, 즉 42개의 150스텝 전개(제어 스텝당 약 6300회의 회복 입력 계산)를 수행하며, 학습 변형은 그때마다 98차원 입력의 신경망을 평가한다.
GPU(RTX 4050 노트북)에서 1000개 에피소드를 묶어 평가하면 15\,s 에피소드 전체가 약 350\,s 걸린다(묶음 기준 약 0.46\,ms/기체-스텝).
단일 기체의 실시간 최악 실행 시간은 측정하지 않았다.

%======================================================================
\section{비교 필터: bCBF-QP}\label{sec:bcbf}

PS2-RL의 backup CBF QP를 같은 집합에 적용한 비교 구성이다. 기체와 추격자를 포함한 상태를 $y=(p,v,q,c_1,\dots,c_C)$, 결합 동역학을 $\dot y=f(y)+g(y)u$,
고정된 회복 정책(반사 포함)의 흐름을 $\varphi(\tau,y)$, 그 벡터장을 $f_\pi$라 하자.
민감도 $Q_i=\partial\varphi(\tau_i,y)/\partial y$는 Euler 전개의 정확한 Jacobian으로 전파한다.
절대 예측 시각 $t_i$를 고정하면 $\tau_i=t_i-t$가 시간에 따라 줄어들므로, 매끄러운 구간에서
\begin{equation}
\frac{d}{dt}h\big(\varphi(t_i-t,y(t))\big)=\nabla h(\varphi_i)\big[Q_i\big(f(y)+g(y)u\big)-f_\pi(\varphi_i)\big]
\end{equation}
이다. 둘째 항이 남은 예측 시간이 줄어드는 데서 나온다. 이를 $-\alpha h(\varphi_i)$ 이상으로 제한한 선형 제약을 노드 $\tau_i$(1.5\,s 구간, 3스텝 간격)와 모든 안전 함수 $h\in\{h_o,h_f,h_w,h_z,h_c\}$에 대해 둔다($\alpha=4$).
종단 노드에서는 기준 집합 함수(반사가 작동 중이면 제동점 함수)와 $h_c$에 같은 형태의 제약을 $\alpha_B=2$로 둔다.
\begin{equation}
\min_{u,\ s\ge0}\ \tfrac12\sum_kw_k(u_k-u_{\mathrm{nom},k})^2+\tfrac12w_ss^2\quad\text{s.t. 완화된 CBF 제약, }u\in\mathcal U,\qquad w=(1,0.1,0.1,1),\ w_s=10^6 .
\end{equation}
roll/pitch 각속도의 가중치가 작아 추력보다 기울기로 경로를 바꾸는 해를 선호한다.
회복 입력 $u=\pi_b(x)$가 제약을 $10^{-2}$보다 크게 위반하면(정규화된 행 기준) QP를 건너뛰고, QP 출력은 다시 이산 필터(최종 회복 정책군)를 통과시킨다.
따라서 이 구성의 접촉 억제는 QP 단독의 결과가 아니며, 선형화·이산화·slack 때문에 표준 연속시간 CBF 불변성 정리를 그대로 적용할 수 없다.

%======================================================================
\section{임무 정책}\label{sec:striker}

\paragraph{기준 접근 제어} 표적 위치 $p_t$와 속도 $v_t$에 대해
\begin{equation}
v_{\mathrm{des}}=v_{\mathrm{cr}}\frac{p_t-p}{\|p_t-p\|}+v_t,\qquad v_{\mathrm{cr}}=6\ \mathrm{m/s}
\end{equation}
를 회피 반사와 같은 속도·자세 루프($k_v=2.5$, 기울기 $50^\circ$, $k_{\mathrm{att}}=10$, $k_\psi=3$)로 추종해 $\upur$를 만든다. 장애물과 추격자를 고려하지 않는다.

\paragraph{잔차 정책}
\begin{equation}
\unom=\sat\big(\upur(x)+\Delta u_\xi(o)\big),\qquad |\Delta a_T|\le0.1\,a_{\max}=3.92\ \mathrm{m/s^2},\quad |\Delta\omega_i|\le0.25\,\omega_{\max}=4.5\ \mathrm{rad/s}.
\label{eq:residual}
\end{equation}
관측은 135차원이다: 회복 정책 형식의 장면 관측 98개, 표적 오프셋/10, 표적 속도, 추격자 속도, 경기장 내 위치(14개),
유도 특징 19개(시선 방향 3, $d/10$, $\log(d+0.1)$, 상대 속도/5 3, 접근 속도/5, 시선각 변화율/2 3, 요 좌표계의 상대 위치/10 3과 상대 속도/5 3, 표적 높이/2),
그리고 보정 대상인 기준 입력 $(\upur-\ustar)/u_{\mathrm{scale}}$(4개).

\paragraph{보상} 표적 거리 $d_k=\|p_k-p_{t,k}\|$에 대해
\begin{equation}
r_k=w_p(d_{k-1}-d_k)-w_t-\frac{w_a}{4}\Big\|\frac{\Delta u_k}{u_{\mathrm{scale}}}\Big\|^2+100\cdot\mathbf 1_{\text{성공}}-150\cdot\mathbf 1_{\text{접촉 또는 이탈}},\qquad w_p=2,\ w_t=w_a=0.05 .
\label{eq:reward}
\end{equation}
거리 항은 $w_p(d_0-d_{\mathrm{end}})$로 합쳐지고, 시간 비용은 빠른 성공을 장려한다. 시간초과는 critic 목표에서 종료로 보지 않는다(부트스트랩 유지).

\paragraph{학습} SAC, $\gamma=0.995$, 보상 척도 0.1, 병렬 환경 64개, 배치 256, 재생 버퍼 $10^6$, 학습률 $3\times10^{-4}$, $\tau=0.005$, 온도 초기값 0.1(하한 $10^{-3}$, 목표 엔트로피 $-4$), 은닉층 256$\times$2(ELU).
처음 20만 스텝은 actor를 고정하고 critic만 학습한다(critic warm-up). actor 손실에 고정점 벌점 $0.2\,\mathbb E\big[\tfrac14\|\Delta u/b\|^2\big]$($b$는 잔차 한계)을 더한다.
1000만 스텝 학습하며, 25만 스텝마다 128개 에피소드로 평가해 반환값이 가장 큰 체크포인트를 쓴다. 처음부터 학습한 정책은 실패했고(성공 0, 이후 0.41--0.49), 기준 접근 제어의 잔차로 두어야 안정적으로 학습되었다.

%======================================================================
\section{PS2: 필터를 포함한 임무 학습}\label{sec:ps2}

필터를 환경의 일부로 두고 임무 정책을 다시 학습한다.
\begin{align}
u_{\mathrm{nom},k}&=\sat\big(\upur(x_k)+\Delta u_\xi(o_k)\big),\qquad (u_k,m_{k+1})=F\big(x_k,\zeta_k,m_k,u_{\mathrm{nom},k}\big),\qquad x_{k+1}=f_d(x_k,u_k), \label{eq:ps2}\\
r^{\mathrm{PS2}}_k&=r_k-w_f\|\eta_k\|^2,\qquad \eta_k=\frac{u_k-u_{\mathrm{nom},k}}{u_{\mathrm{scale}}},\qquad w_f=0.5 . \label{eq:ps2-reward}
\end{align}
정책은 자신이 제안한 입력이 그대로 실행된다고 가정하지 않고, 필터 개입 이후의 상태와 보상으로 학습한다. 필터가 자주 막는 행동보다 실제로 진전이 생기고 수정 비용이 작은 행동이 유리해진다.
필터 기억 $m$까지 포함한 전이 $(x,\zeta,m)\mapsto(x^+,\zeta^+,m^+)$는 정책에서 보면 환경의 일부이므로, 이산 필터를 미분하지 않고도 off-policy SAC로 학습할 수 있다.
다만 $m$은 관측되지 않으므로 관측이 완전한 상태라고 주장하지 않는다. 필터 출력을 복사하는 모방학습이 아니라, 수정된 전이에서 같은 강화학습 목표를 다시 최적화하는 것이다.

\paragraph{설정} 같은 시드의 잔차 정책 최고 체크포인트에서 시작하여, 처음 10만 스텝은 actor를 고정하고 critic을 재적응시키며, 온도는 시작 값으로 고정한다(학습률 0).
배치할 필터(최종 회복 정책군, 추격자 예측 재생)와 같은 필터로 300만 스텝 학습하고(학습 = 배치), 25만 스텝마다 64개 에피소드로 평가해 반환값이 가장 큰 체크포인트를 쓴다.
PS2 정책은 필터와 함께 쓰도록 학습한 정책이며, 필터 없이 실행하면 성공률이 잔차 정책보다 낮다(62.1\,\% 대 74.0\,\%). 필터가 장애물에서 비켜 주는 것에 의존하기 때문이다.

\paragraph{학습 시간} RTX 4050 노트북 GPU에서 시드당 Phase I 약 10분(1000만 스텝), 잔차 정책 28--34분(1000만 스텝), PS2 72--76분(300만 스텝)이 걸렸다.

%======================================================================
\section{평가 방법}\label{sec:protocol}

\paragraph{절차} 시드 $s\in\{0,1,2\}$마다 전체 파이프라인(Phase I $\to$ 잔차 정책 $\to$ PS2)을 그 시드로 학습한다.
모든 구성을 같은 1000개 평가 장면(장면 시드 9001, 학습 장면 풀과 별도)으로 평가하며, 장면마다 출발 교란과 표적·추격자 초기 조건까지 같다.
에피소드는 15\,s이고, 10\,s 결과는 같은 에피소드를 10\,s에서 절단한 것이다(10\,s까지 끝나지 않은 에피소드는 시간초과로 센다).
각 시드는 그 시드의 임무 정책을 학습한 추격자 설정(원래 추격자)으로 평가한다.

\paragraph{지표} 성공률, 접촉률(장애물·바닥·추격자), 시간초과율, 영역 이탈률, 성공 에피소드의 평균 도달 시간,
평균 명목 유지율 $\bar\lambda$(에피소드 스텝 평균 $\lambda^\star$, 새 인증이 없는 스텝은 0),
인증 상실률(한 스텝 이상 유효한 인증 계획이 없었던 에피소드의 비율). 인증 상실은 접촉과 다른 지표이다.

\paragraph{통계} 시드별 비율의 평균 $\pm$ 표준편차를 보고하고, 시드를 합친 횟수 $k/n$에 대해 양측 95\,\% Clopper--Pearson 구간을 쓴다.
\begin{equation}
p_L=\mathrm{B}^{-1}\!\big(\tfrac{\alpha}{2};\,k,\,n-k+1\big),\qquad p_U=\mathrm{B}^{-1}\!\big(1-\tfrac{\alpha}{2};\,k+1,\,n-k\big),\qquad \alpha=0.05,
\end{equation}
여기서 $\mathrm B^{-1}$는 베타 분포의 분위수이다. $k=0$이면 $p_U=1-(\alpha/2)^{1/n}$이고, 단측 95\,\% 상한은 $1-0.05^{1/n}\approx3/n$이다.
$n=1000$에서 0.30\,\%(단측)와 0.37\,\%(양측), $n=3000$에서 0.10\,\%와 0.12\,\%이다.
모든 구성이 같은 장면을 쓰므로 두 구성의 비교는 \emph{쌍 비교}이다. 구성 A만 성공한 장면 수 $b$와 B만 성공한 장면 수 $c$로 정확 McNemar 검정을 한다.
\begin{equation}
p_{\mathrm{McNemar}}=\min\Big\{1,\ 2\sum_{i=0}^{\min(b,c)}\binom{b+c}{i}2^{-(b+c)}\Big\}.
\end{equation}
구성끼리 비교하는 표(표~\ref{tab:table2_final}, \ref{tab:table3_ablation})는 모든 행이 평가된 공통 시드만 쓴다.

%======================================================================
\section{결과}\label{sec:results}

\begin{warnbox}{수치의 기준}
아래 표와 그림은 업로드된 평가 결과(시드 0 전체, 시드 1 일부, 시드 2 미평가)에서 만든 것이다. 각 행의 ``시드 $\times$ 장면'' 열이 사용한 에피소드 수를 나타낸다.
세 시드 평가를 마친 뒤 \file{STEPS=tables bash docker\_batch/run\_strike\_paper.sh}로 다시 만들고 이 문서를 다시 컴파일하면 표와 그림이 갱신된다.
\end{warnbox}

\subsection{주요 비교}
\tabin{table1_10s}
\tabin{table1_15s}

\begin{figure}[H]\centering
\figin{\figpath{fig_outcomes.pdf}}
\caption{구성별 결과 비율(왼쪽 10\,s, 오른쪽 15\,s 제한). 막대 안의 수는 성공률(\%)이다.}\label{fig:outcomes}
\end{figure}

\paragraph{필터의 효과} 같은 임무 정책에서 최종 필터는 접촉을 거의 없앤다. 잔차 정책은 26.1\,\%(521/2000)에서 0/2000으로, 기준 접근 제어는 57.5\,\%에서 1/2000으로 줄었다.
대신 우회와 정지 때문에 10\,s 시간초과가 늘어난다(잔차 정책 0\,\%$\to$15.5\,\%). 15\,s 제한에서는 시간초과 대부분이 늦은 성공으로 바뀐다.

\paragraph{PS2의 효과} 같은 최종 필터를 쓸 때 PS2 정책은 잔차 정책보다 성공률이 높다(10\,s 84.5$\to$90.4\,\%, 15\,s 94.4$\to$97.6\,\%).
같은 장면 쌍 비교에서 PS2만 성공한 장면이 172개, 잔차 정책만 성공한 장면이 54개로(10\,s) 차이는 유의하다($p<10^{-6}$, 표~\ref{tab:table5_paired}).
평균 도달 시간도 6.53\,s에서 6.11\,s로 줄었다. 필터를 고정하고 임무 정책만 바꾼 비교이므로 필터 포함 재학습의 효과를 보여 주는 핵심 근거이다.

\begin{figure}[H]\centering
\begin{subfigure}{0.49\linewidth}\centering\figin{\figpath{fig_strike_time.pdf}}\caption{누적 성공률}\label{fig:strike-time}\end{subfigure}\hfill
\begin{subfigure}{0.49\linewidth}\centering\figin{\figpath{fig_safety.pdf}}\caption{추격자 최소 여유 거리의 누적 분포}\label{fig:safety}\end{subfigure}
\caption{(a) 시간에 따른 누적 성공률(음영은 시드를 합친 95\,\% 구간). 점선은 10\,s. (b) 에피소드별 추격자 최소 여유 거리 $\min_t(\|p-c\|-r_c-r_d)$의 누적 분포. 필터를 쓴 구성에서 0.2\,m의 계단은 추격자가 바짝 따라오는 동안 필터가 $h_c$를 신규 계획 여유 $m_c=0.2$\,m에 유지한 에피소드들이다(멀어지는 중에는 $h_c$가 거리 여유와 같다).}
\end{figure}

\subsection{최종 필터 선택과 구성 요소}
\tabin{table2_final}
\tabin{table3_ablation}
\tabin{table5_paired}

\begin{figure}[H]\centering
\figin{\figpath{fig_ablation.pdf}}
\caption{구성 요소 제거 실험(PS2 정책, 15\,s). 오차 막대는 95\,\% Clopper--Pearson 구간이다.}\label{fig:ablation}
\end{figure}

\begin{itemize}
\item \textbf{회피 변형}: 단일 회피 필터는 접촉 1.7\,\%(17/1000)로, 최종 필터(1/1000)와의 쌍 비교에서 차이가 유의하다($p=3.1\times10^{-5}$). 성공률도 낮다(15\,s 95.7\,\% 대 98.4\,\%).
\item \textbf{bCBF-QP}: 성공률은 비슷하지만 최종 필터가 10\,s와 15\,s 모두 높고($p=0.0015$, $0.043$), 평균 도달 시간이 짧다(5.93\,s 대 6.33\,s, 공통 시드 기준). bCBF는 명목 유지율이 더 높지만(0.97) 이후 이산 검사에서 인증을 더 자주 잃는다(2.5\,\% 대 1.5\,\%).
\item \textbf{학습 회복 정책}: 회복 정책을 호버 LQR로 바꾸면 15\,s 성공률이 82--84\,\%로 크게 떨어지고 시간초과가 15--17\,\%로 늘어난다(시드 2개와 공통 시드 기준). Phase I 회복 정책이 필터의 허용 범위를 넓히는 것이 성능의 주요 원인이다.
\item \textbf{구형 근사}: 추격자를 정지한 구로 보면 접촉이 26.3\,\%로, 상대 접근 속도를 반영한 $h_c$가 필수적이다.
\item \textbf{계획 유지의 추격자 예측 재생과 호버 latch}: 1000개 장면에서 성공·접촉의 차이는 관측되지 않았다(98.6\,\%, 98.3\,\% 대 98.4\,\%). 두 요소는 성능보다 이어짐이 인증된 계획을 정확히 재현하고 인증 대상 제어기와 일치하게 하는 일관성을 위한 것이다.
\end{itemize}

\begin{figure}[H]\centering
\figin{\figpath{fig_filter.pdf}}
\caption{필터 개입. 왼쪽: 에피소드 평균 명목 유지율 $\bar\lambda$의 누적 분포(1에서의 계단은 필터가 한 번도 수정하지 않은 에피소드). 오른쪽: 인증 상실 에피소드 비율.}\label{fig:filter}
\end{figure}

\subsection{학습 과정과 Phase I}
\begin{figure}[H]\centering
\figin{\figpath{fig_training.pdf}}
\caption{학습 중 평가 성공률(시드별). 원은 평가에 사용한 체크포인트. 필터 없이 학습한 잔차 정책은 시드 1, 2에서 500--650만 스텝 이후 붕괴했으며, 평가에는 최고 반환값 체크포인트(시드 0: 850만, 시드 1: 125만, 시드 2: 175만 스텝)를 썼다. PS2 학습은 필터 안에서 안정적이다.}\label{fig:training}
\end{figure}
\tabin{table4_phase1}
\begin{figure}[H]\centering
\figin[0.55\linewidth]{\figpath{fig_phase1.pdf}}
\caption{Phase I 회복 성공률(1.5\,s 안에 안전하게 기준 집합 도달). 학습 회복 정책은 세 시드 평균과 표준편차.}\label{fig:phase1}
\end{figure}

\subsection{대표 에피소드}\label{sec:results-episodes}
\file{run\_strike\_paper.sh}는 시드 0의 평가에서 필터가 장애물 충돌을 막은 장면, 추격자 접촉을 막은 장면, 추격자 근접 회피, 유일한 접촉 장면,
기준 제어가 실패한 장면 등을 자동으로 고르고(\file{episode\_picks.json}), 각 장면을 여러 구성으로 다시 실행해 다음을 만든다.
\begin{itemize}
\item 2차원 궤적 비교(구성별 패널), 3차원 궤적(속력 색), 3차원 비교 그림;
\item 상세 신호 그림: 표적 거리, 장애물·벽·바닥·천장 여유, $h_c$와 반사 가중치 $\sigma$, 속력과 고도, $\lambda^\star$와 인증 상태(신규/이어짐/인증 없음), 사용한 회복 변형, 입력과 명목 입력;
\item 2차원 GIF와 회전하는 3차원 GIF(추격자 $h_c$ 요구 거리 구 표시);
\item 필터 스냅샷(\file{scripts/strike\_filter\_snapshot.py}): 처음 명목 입력을 줄인 스텝, 이전 계획을 이어 간 스텝, 인증을 잃은 스텝에서 필터의 검사를 다시 계산한 그림(그림~\ref{fig:snapshot}).
\end{itemize}
\IfFileExists{\episodedir/overview_main.png}{%
\begin{figure}[H]\centering
\includegraphics[width=\linewidth]{\episodedir/overview_main.png}
\caption{대표 에피소드(시드 0)의 수평 궤적. 실선은 최종 필터를 쓴 PS2 정책과 잔차 정책, 점선은 필터 없는 PS2 정책이다. 별은 성공, X는 접촉, 회색 원은 시간초과.}\label{fig:episodes-overview}
\end{figure}}{\IfFileExists{\episodedir/overview.png}{%
\begin{figure}[H]\centering
\includegraphics[width=\linewidth]{\episodedir/overview.png}
\caption{대표 에피소드(시드 0)의 수평 궤적.}\label{fig:episodes-overview}
\end{figure}}{}}

\paragraph{유일한 접촉 사례(시드 0, 에피소드 845)} 기체는 큰 원기둥과 위쪽 벽 사이의 좁은 공간($y\approx9$\,m)으로 들어간 뒤 약 6\,s부터 새 인증을 찾지 못해
이전 계획(호버 LQR + 회피 변형)을 이어 갔고, 약 7.5\,s부터는 그 계획도 재검사를 통과하지 못했다(인증 없음).
추격자가 다가오면서 회피 반사가 작동했지만 벽과 장애물이 도피 방향을 막아 9.2\,s에 접촉했다. 같은 장면에서 bCBF 구성도 같은 방식으로 실패했다.
이 사례는 \S\ref{sec:guarantee}의 남는 간극, 즉 호버 종단이 다가오는 추격자에 대해 불변 집합이 아니고 막힌 공간에서는 도피도 보장되지 않는다는 점이 실제로 나타난 경우이다.
인증을 잃은 7.5\,s의 검사를 다시 계산하면(그림~\ref{fig:snapshot-lost}) 35개 후보 모두 궤적 조건을 위반했다. 위쪽 벽($y=10$\,m) 쪽으로 움직이던 기체에서 일곱 회복 정책의 3\,s 계획이 모두 벽을 넘었기 때문이다.
Phase I 회복 정책은 벽 없이 학습되었고, 회피 반사는 추격자 여유가 $h_{\mathrm{off}}=2$\,m 이상이면 섞이지 않는다. 벽을 포함한 Phase I 학습이 남은 개선 방향이다.
\IfFileExists{\snapdir/snapshot_ep845_step374.png}{%
\begin{figure}[H]\centering
\includegraphics[width=\linewidth]{\snapdir/snapshot_ep845_step374.png}
\caption{인증을 잃은 시점의 검사(에피소드 845, $t=7.5$\,s). 일곱 회복 정책의 계획이 모두 위쪽 벽을 넘어 새 인증이 없고, 필터는 이전 계획을 이어 실행했다.}\label{fig:snapshot-lost}
\end{figure}}{}
\IfFileExists{\episodedir/ep845/compare2d.png}{%
\begin{figure}[H]\centering
\includegraphics[width=\linewidth]{\episodedir/ep845/compare2d.png}
\caption{에피소드 845를 여러 구성으로 다시 실행한 수평 궤적(CPU 재실행).}\label{fig:ep845-2d}
\end{figure}}{}
\IfFileExists{\episodedir/ep845/detail_ps2_ours.png}{%
\begin{figure}[H]\centering
\includegraphics[width=0.86\linewidth]{\episodedir/ep845/detail_ps2_ours.png}
\caption{에피소드 845, PS2 정책 + 최종 필터의 상세 신호. 위에서부터 표적 거리, 기하 여유, $h_c$와 반사 가중치 $\sigma$, 속력과 고도, $\lambda^\star$와 인증 상태(노랑: 이전 계획 유지, 빨강: 인증 없음), 사용한 회복 변형, 입력과 명목 입력.}\label{fig:ep845-detail}
\end{figure}}{}

%======================================================================
\section{한계와 적용 범위}\label{sec:limits}
\begin{enumerate}
\item \textbf{종단 집합의 불변성}: 호버 기준 집합은 정적 장애물에 대한 회복 수단이다. 추격자가 계속 다가오면 호버는 영구히 안전하지 않고, 도피 종단도 유한시간 예측의 조건이다. 정적 기준 집합에 기대는 원래 이론의 보장을 그대로 옮길 수 없다.
\item \textbf{추격자 모델}: 예측은 현재 시선 방향의 순수 추적이며, 장애물 우회와 시선각 변화를 포함한 모든 추격 전략의 최악값이라는 증명은 없다. 반응 지연과 이탈 가속도에 기반한 $h_c$도 근사이다.
\item \textbf{모델 오차}: 입력 지연, 공력, 센서 오차, 바람이 있으면 3\,s 검사와 실제 경로가 달라질 수 있다. 모든 오차를 포함하는 강건 경계는 제시하지 않았다.
\item \textbf{시간 이산 검사}: 검사는 20\,ms 노드에서 한다. 노드 사이 연속시간 충돌까지 배제하려면 이동량이나 보간 오차를 이용한 추가 조건이 필요하다.
\item \textbf{인증 상실}: 계획 유지 중에도 인증을 잃을 수 있다(PS2 에피소드의 1.5\,\%). 이 시점은 명제의 전제를 벗어난다.
\item \textbf{통계}: 1000개 장면과 2--3개 시드의 결과이다. 접촉 1/2000은 접촉확률이 0이라는 증명이 아니며 95\,\% 상한은 0.28\,\%이다.
\item \textbf{실기체}: 결과는 시뮬레이션의 성공과 접촉이다. 실시간 최악 실행 시간, PX4 추력·각속도 추종 오차, 실기체 운용 성능은 검증하지 않았다.
\end{enumerate}

%======================================================================
\appendix
\section{하이퍼파라미터}\label{app:hparams}
\begin{center}\small
\begin{tabularx}{\linewidth}{>{\raggedright\arraybackslash}p{3.0cm}>{\raggedright\arraybackslash}X}\toprule
구분 & 값 \\\midrule
모델 & $\Delta t=0.02$\,s, $a_T\in[0,4g]$, $|\omega_i|\le18$\,rad/s, $r_d=0.3$\,m, $z_f=1$\,m, $z_{\max}=10$\,m \\
경기장·장면 & 40$\times$20\,m, 장애물 3개, 추격자 1개($r_c=0.4$\,m, $v_c=2$\,m/s, $a_c=2$\,m/s$^2$, 반발 거리 2\,m), 표적 0.2--0.8\,m/s, 성공 반경 0.6\,m, 15\,s \\
$h_c$ & $\delta=0.2$\,s, $a_e=5$\,m/s$^2$ \\
호버 LQR & $Q=\mathrm{diag}(2,2,8,0.01,0.01,0.5,8,8,2)$, $R_u=\mathrm{diag}(0.06,0.15,0.15,0.12)$, $c_B=50$, $\rho_{xy}=1.22$\,m, $z_{\mathrm{ext}}=0.63$\,m \\
Phase I & TD3, 1000만 스텝, 병렬 64, 배치 256, 버퍼 $10^6$, $\beta=0.99$, 학습률 actor $10^{-4}$/critic $3\times10^{-4}$, $\tau=0.0025$, 정책 지연 2, 탐색 잡음 0.1, 잔차 벌점 0.2, 에피소드 112스텝, 구형 근사 $T=1.5$\,s, 영역 가중치 $(0.4,0.2,0.2,0.2)$ \\
회피 반사 & $h_{\mathrm{on}}=0.5$\,m, $h_{\mathrm{off}}=2.0$\,m, $(k_v,k_{\mathrm{att}},k_\psi)=(2,10,3)$, 기울기 $45^\circ$, $d_{\mathrm{push}}=1.5$\,m, $z_{\mathrm{ref}}\in[z_f+1,z_{\max}-3]$ \\
회복 정책군 & $(v_{\mathrm{esc}}/v_c,\vartheta)\in\{(2,0^\circ)\text{ 학습},(2,0^\circ)\text{ LQR},(2,\pm60^\circ),(3.5,0^\circ),(3.5,\pm60^\circ)\}$ \\
이산 필터 & $\Lambda=\{1,0.75,0.5,0.25,0\}$, 예측 3\,s(150스텝), 신규 계획 여유 $m_o=0.15$\,m·$m_c=0.2$\,m, 종단 $\Term_{\mathrm{hover}}\cup\Term_{\mathrm{flee}}$, 계획 유지(추격자 예측 재생), 호버 latch \\
bCBF-QP & $\alpha=4$, $\alpha_B=2$, 1.5\,s, 3스텝 간격, $w=(1,0.1,0.1,1)$, $w_s=10^6$, 건너뛰기 허용치 $10^{-2}$, 보조 필터 = 최종 이산 필터 \\
기준 접근 제어 & $v_{\mathrm{cr}}=6$\,m/s, $k_v=2.5$, 기울기 $50^\circ$ \\
잔차 정책 & 잔차 한계 추력 10\,\%, 각속도 25\,\%; 관측 135; SAC $\gamma=0.995$, 보상 척도 0.1, 병렬 64, 배치 256, 학습률 $3\times10^{-4}$, $\tau=0.005$, 온도 0.1(하한 $10^{-3}$), critic warm-up 20만, 고정점 벌점 0.2, 1000만 스텝 \\
보상 & $w_p=2$, $w_t=0.05$, $w_a=0.05$, 성공 $+100$, 접촉·이탈 $-150$ \\
PS2 & $w_f=0.5$, 잔차 정책에서 시작, critic 재적응 10만, 온도 고정, 300만 스텝 \\
평가 & 시드 3개, 장면 1000개(시드 9001), 15\,s와 10\,s 절단, Clopper--Pearson 95\,\%, 정확 McNemar \\\bottomrule
\end{tabularx}
\end{center}

\section{기호표}\label{app:symbols}
\begin{center}\small
\begin{tabularx}{\linewidth}{>{\raggedright\arraybackslash}p{3.4cm}X}\toprule
기호 & 의미 \\\midrule
$x=(p,v,q)$, $u=(a_T,\omega)$ & 기체 상태, 입력(추력 가속도, 기체 각속도) \\
$f_d$, $\Delta t$ & 이산 동역학, 시간 간격 \\
$\zeta$, $c_\ell$, $\nu_\ell$ & 장면 정보, 추격자 위치, 추격자 속도 \\
$r_d$, $r_c$, $v_c$, $a_c$ & 기체 반경, 추격자 반경, 추격자 속력·가속도 상한 \\
$h_f,h_o,h_w,h_z$ & 바닥, 장애물, 벽, 천장 여유 \\
$h_{c,j}$, $s_j^+$, $\delta$, $a_e$ & 추격자 여유도, 양의 접근 속도, 반응 지연, 유효 이탈 가속도 \\
$\Safe$, $\Base$, $\BaseP$ & 안전 집합, 필터의 기준 집합, Phase I 기준 집합 \\
$e$, $P$, $K$, $V$ & 호버 오차, Riccati 해, LQR 이득, $V=e^\top Pe$ \\
$a^\star$, $D$, $S$ & 최적 정지점, 정지점 행렬, Schur 여행렬 ($V_{\min}=r^\top Sr$) \\
$c_B$, $\rho_{xy}$, $z_{\mathrm{ext}}$ & 기준 집합 수준, 타원체의 수평·수직 범위 \\
$\pi_\theta$, $\Delta u_\theta$ & Phase I 학습 회복 정책과 그 잔차 \\
$\pi_b^{(j)}$, $\uesc^{(j)}$, $\sigma$ & 회복 정책 변형 $j$, 회피 입력, 반사 가중치 \\
$\Lambda$, $\lambda^\star$, $\bar\lambda$ & 혼합률 집합, 선택한 혼합률, 에피소드 평균 \\
$N$, $H$, $m_o$, $m_c$ & 예측 스텝(150), 종단 첨자($N+1$), 신규 계획 여유 \\
$\Term$, $p_{\mathrm{brake}}$ & 종단 집합, 제동 정지점 \\
$F$, $m$ & 안전 필터, 필터 기억(계획 유지) \\
$\upur$, $\Delta u_\xi$, $\unom$ & 기준 접근 입력, 임무 정책 잔차, 명목 입력 \\
$r_k$, $r^{\mathrm{PS2}}_k$, $\eta_k$, $w_f$ & 보상, PS2 보상, 정규화한 필터 수정량, 수정 비용 가중치 \\\bottomrule
\end{tabularx}
\end{center}

\section{원고 초안과 다른 점}\label{app:draft}
초안을 이 버전의 결과로 바꿀 때 확인할 항목이다.
\begin{enumerate}
\item \textbf{평가 규모}: 초안의 ``단일 시드, 128개 장면'' 수치(44.5, 82.8, 93.8, 100\,\% 등)는 이전 실험이다. 현재 표는 1000개 장면, 시드 2--3개이며 15\,s 에피소드와 10\,s 절단을 함께 보고한다.
\item \textbf{접촉 0건 표현}: PS2 정책 + 최종 필터는 접촉 1/2000(시드 0의 한 장면)이다. ``접촉이 관측되지 않았다'' 대신 ``2000회 중 1회(95\,\% 상한 0.28\,\%)''로 쓴다. 잔차 정책 + 최종 필터는 0/2000이다.
\item \textbf{15\,s 성공률}: 초안의 ``128개 장면 모두 성공(100\,\%)''은 현재 97.6\,\%이다(시간초과 2.4\,\%).
\item \textbf{PS2 정책의 학습 필터}: 초안 표 1은 단일 회피 필터로 학습한 PS2 정책을 썼다. 최종 파이프라인의 PS2 정책은 배치하는 최종 필터(회피 변형, 추격자 예측 재생)로 학습했다.
\item \textbf{bCBF slack 가중치}: 초안 식 (34) 뒤의 $w_s=10$은 $w_s=10^6$이다.
\item \textbf{신규 계획 여유}: 0.15\,m는 장애물과 벽뿐 아니라 바닥과 천장에도 적용된다.
\item \textbf{기준 집합의 추격자 조건}: 필터의 기준 집합 $\Base$에는 추격자 조건이 없고, 추격자는 종단 조건 $h_c\ge0.2$\,m와 회피 반사로 다룬다. Phase I의 기준 집합만 추격자 구 조건을 가진다.
\item \textbf{천장 조건}: 식 (20)의 $z_{\max}-a^\star_z\ge\rho_{\mathrm{top}}+r_d$에서 구현은 $\rho_{\mathrm{top}}=\rho_{xy}$이다(초안과 같음, 보수적).
\item \textbf{제동 정지점}: 식 (30)은 구현에서 수평 성분만 쓴다: $p_{\mathrm{brake}}=p_{H,xy}+v_{H,xy}\|v_{H,xy}\|/(2a_e)$. 이 점이 장애물과 벽에서 $r_d$ 이상 떨어져야 한다.
\item \textbf{추격자 운동}: 환경의 추격자는 속력 2\,m/s와 함께 가속도 2\,m/s$^2$ 상한을 가지며, 장애물 2\,m 이내에서 밀려난다(\S\ref{sec:arena}). 필터의 예측(가속도 제한 없는 순수 추적)과 다르다는 점을 밝혀 둔다.
\item \textbf{Phase I 수치}: 가중 성공률은 세 시드 평균 70.1$\pm$0.3\,\%(초안 69.9\,\%), LQR 43.8\,\%이다.
\item \textbf{이동 추격자 회복 학습 변형}: 초안 표 2의 해당 행은 최종 파이프라인 평가에 포함되지 않은 128개 장면 실험이다. 쓰려면 이전 실험임을 밝히거나 표에서 뺀다.
\item \textbf{수식 표기}: 초안 식 (2)·(6)의 $|\omega_i)$, $|\ell_i)$는 절댓값 $|\omega_i|$, $|\ell_i|$이고, 식 (4)의 괄호는 $v_{k+1}=v_k+\Delta t\,(a_{T,k}R(q_k)e_3-ge_3)$이다. 회피 각 $\vartheta$와 임무 정책 매개변수(초안 $\Delta u_\vartheta$)의 기호가 겹치므로 이 문서는 임무 정책 잔차를 $\Delta u_\xi$로 썼다.
\end{enumerate}

\section{한글(HWP) 수식 편집기 입력}\label{app:hwp}
본문의 주요 수식을 한글 수식 편집기 문법으로 옮긴 것이다(같은 내용: \file{results/strike\_paper/hwp\_equations.txt}). 붙여넣은 뒤 모양을 확인한다.
\IfFileExists{hwp_equations.txt}{\VerbatimInput[fontsize=\scriptsize,breaklines,breakanywhere,firstline=4]{hwp_equations.txt}}{\fbox{hwp\_equations.txt 없음}}

\end{document}
"
\providecommand{\paperdir}{../strike_reduced/final/paper}   % the report's results in this branch
\providecommand{\episodedir}{\paperdir/episodes}
\providecommand{\snapdir}{\paperdir/traj4/snapshots}   % filter snapshots
\documentclass[10pt,a4paper]{article}
\usepackage[margin=2.1cm]{geometry}
\usepackage{fontspec}
\setmainfont{Noto Serif CJK KR}
\setsansfont{Noto Sans CJK KR}
\setmonofont{Noto Sans Mono CJK KR}
\usepackage{amsmath,amssymb,bm,mathtools}
\usepackage{graphicx,float}
\usepackage{booktabs,array,tabularx,multirow}
\usepackage[dvipsnames]{xcolor}
\usepackage{enumitem}
\usepackage{caption,subcaption}
\usepackage{titlesec}
\usepackage{fancyhdr}
\usepackage[most]{tcolorbox}
\usepackage[ruled,vlined,linesnumbered]{algorithm2e}
\usepackage{fvextra}
\usepackage{tikz}
\usetikzlibrary{arrows.meta,positioning,fit,calc}
\usepackage[colorlinks=true,linkcolor=NavyBlue,urlcolor=NavyBlue,citecolor=NavyBlue]{hyperref}

\linespread{1.18}
\setcounter{tocdepth}{1}
\floatplacement{table}{H}
\setlength{\parindent}{1em}
\setlength{\parskip}{2pt}
\emergencystretch=2em
\setlist{nosep,leftmargin=1.5em}
\captionsetup{font=small,labelfont=bf,labelsep=period}
\renewcommand{\figurename}{그림}
\renewcommand{\tablename}{표}
\renewcommand{\contentsname}{목차}
\renewcommand{\refname}{참고 문헌}
\renewcommand{\abstractname}{요약}
\renewcommand{\appendixname}{부록}
\SetAlgorithmName{알고리즘}{알고리즘}{알고리즘 목록}
\SetKwInput{KwIn}{입력}
\SetKwInput{KwOut}{출력}
\SetKwInput{KwMem}{기억}
\SetKw{KwRet}{반환}
\SetKwIF{If}{ElseIf}{Else}{만약}{이면}{그렇지 않고 만약}{그렇지 않으면}{끝}
\SetKwFor{For}{각}{에 대해}{끝}
\titleformat{\section}{\large\bfseries\sffamily\color{NavyBlue}}{\thesection}{0.6em}{}
\titleformat{\subsection}{\normalsize\bfseries\sffamily}{\thesubsection}{0.6em}{}
\titleformat{\paragraph}[runin]{\normalsize\bfseries}{}{}{}[.]
\pagestyle{fancy}\fancyhf{}
\lhead{\small STRIKE: 방법과 결과 상세 정리}
\rhead{\small 2026.\ 10.\ 23}
\cfoot{\small\thepage}
\renewcommand{\headrulewidth}{0.3pt}

\newtcolorbox{keybox}[1]{colback=NavyBlue!4,colframe=NavyBlue!55,fonttitle=\bfseries\sffamily\small,title={#1},
  left=4pt,right=4pt,top=2pt,bottom=2pt,fontupper=\small}
\newtcolorbox{warnbox}[1]{colback=Orange!6,colframe=Orange!70,fonttitle=\bfseries\sffamily\small,title={#1},
  left=4pt,right=4pt,top=2pt,bottom=2pt,fontupper=\small}

\newcommand{\R}{\mathbb{R}}
\newcommand{\sat}{\operatorname{sat}}
\newcommand{\clip}{\operatorname{clip}}
\newcommand{\sdf}{\operatorname{sdf}}
\newcommand{\nrm}{\operatorname{normalize}}
\newcommand{\Safe}{\mathcal{S}}
\newcommand{\Base}{\bar{\mathcal{B}}}
\newcommand{\BaseP}{\mathcal{B}_{\mathrm{I}}}
\newcommand{\Term}{\mathcal{T}}
\newcommand{\unom}{u_{\mathrm{nom}}}
\newcommand{\upur}{u_{\mathrm{pur}}}
\newcommand{\uesc}{u_{\mathrm{esc}}}
\newcommand{\ubase}{u_{\mathrm{base}}}
\newcommand{\ustar}{u^{\star}}
\newcommand{\file}[1]{\texttt{\small #1}}
\newcommand{\figpath}[1]{\paperdir/figures/#1}
\newcommand{\tabin}[1]{\IfFileExists{\paperdir/tables/#1.tex}{\input{\paperdir/tables/#1.tex}}{\begin{center}\fbox{표 파일 없음: #1.tex}\end{center}}}
\newcommand{\figin}[2][\linewidth]{\IfFileExists{#2}{\includegraphics[width=#1]{#2}}{\fbox{그림 파일 없음}}}

\title{\vspace{-1.0cm}\textbf{\sffamily STRIKE: 다중 회복 정책 안전 필터와\\ 필터 포함 임무 학습(PS2)을 이용한 쿼드로터 저고도 표적 접근}\\[4pt]
\large 방법, 수식, 알고리즘, 평가 방법과 결과의 상세 정리 (논문 작성용)}
\author{김중원 \quad (정리 보조: Claude)}
\date{\small 2026년 10월 23일. 대상: 희소 경기장, 원래 추격자(반경 0.4 m, 2 m/s, 2 m/s$^2$)로 학습한 최종 구성}

\begin{document}
\maketitle
\thispagestyle{fancy}

\begin{abstract}
\noindent
장애물과 이동 추격자가 함께 있는 저고도 경기장에서 쿼드로터가 이동 표적의 0.6\,m 이내로 접근하는 문제(STRIKE)를 다룬다.
방법은 두 층으로 이루어진다.
\emph{안전 필터}는 임무 정책의 명목 입력을 일곱 개 회복 정책과 다섯 단계로 혼합한 35개 후보 각각을 한 스텝 실행한 뒤,
그 다음 상태에서 회복 정책을 3\,s 동안 가상으로 전개하여 바닥·장애물·벽·천장 제약과 상대 접근 속도를 반영한 추격자 여유도 $h_c$를 확인하고,
종단 상태가 호버하거나 제동할 수 있는 상태인지 검사한다. 통과한 후보 가운데 명목 입력을 가장 많이 유지하는 입력을 실행하고,
통과 후보가 없으면 직전에 인증한 계획을 그 추격자 예측 그대로 이어서 실행한다.
\emph{임무 정책}은 6\,m/s 기준 접근 제어에 잔차를 더하는 SAC 정책이며, 실제 필터를 포함한 환경에서 필터의 수정량을 벌점으로 받으며 재학습한다(PS2).
1000개 평가 장면과 15\,s 제한에서 PS2 정책과 최종 필터의 조합은 성공률 97.6\,\%, 접촉 2000회 중 1회(0.05\,\%, 95\,\% 상한 0.28\,\%)를 보였고,
같은 필터를 쓰는 잔차 정책(94.4\,\%, 접촉 0/2000)보다 성공률이 높았다(같은 장면 쌍 비교, 정확 McNemar 검정 $p<10^{-6}$).
수치는 현재 업로드된 평가(시드 0 전체, 시드 1 일부)에 기반하며, 세 시드 평가가 끝나면 같은 스크립트로 표와 그림이 갱신된다.
\end{abstract}

\tableofcontents
\vspace{4pt}

%======================================================================
\section{이 문서의 범위와 읽는 법}\label{sec:scope}

이 문서는 학술대회 원고(\emph{STRIKE 다중 회복 정책과 안전 필터 학습을 이용한 쿼드로터의 저고도 표적 접근}) 작성을 위해,
최종 구성의 수식과 알고리즘을 구현(\file{ps2rl/})과 대조하여 정리한 것이다.
모든 수식과 상수는 저장소의 코드와 학습에 사용된 설정 파일(\file{configs.json})에서 확인했다.
원고 초안과 다른 부분은 부록~\ref{app:draft}에, 주요 수식의 한글(HWP) 수식 편집기 입력은 부록~\ref{app:hwp}에 정리했다.

\begin{keybox}{이 버전의 구성 (원고의 ``최종 구성'')}
\begin{itemize}
\item \textbf{경기장}: 40\,m$\times$20\,m, 높은 장애물 3개, 이동 추격자 1개(반경 $r_c=0.4$\,m, 속력 $\le v_c=2$\,m/s, 가속도 $\le a_c=2$\,m/s$^2$), 이동 표적, 15\,s 제한.
\item \textbf{안전 필터}: 이산 다중 회복 필터. 회복 정책 7개(학습 회복 정책 또는 호버 LQR과 회피 반사의 결합), 혼합률 5단계, 3\,s 예측, 신규 계획 여유, 확장 종단 집합, 추격자 예측을 재생하는 계획 유지.
\item \textbf{임무 정책}: 6\,m/s 기준 접근 제어 + SAC 잔차(관측 135차원). PS2: 배치하는 필터를 그대로 포함한 환경에서 300만 스텝 재학습.
\item \textbf{평가}: 세 시드(전체 파이프라인을 시드별로 학습), 1000개 평가 장면(시드 9001), 15\,s 에피소드와 10\,s 절단 결과, Clopper--Pearson 구간과 같은 장면 쌍 비교.
\end{itemize}
\end{keybox}

\begin{figure}[H]\centering
\begin{tikzpicture}[node distance=5mm and 7mm, every node/.style={font=\small},
  box/.style={draw, rounded corners, align=center, minimum height=10mm, inner sep=4pt, fill=#1!8}, box/.default=gray,
  >={Stealth[length=2mm]}]
\node[box=blue] (p1) {Phase I\\ 학습 회복 정책 $\pi_\theta$\\ (안전 도착, TD3)};
\node[box=orange, right=of p1] (fam) {회복 정책군 $\{\pi_b^{(j)}\}_{j=0}^{6}$\\ 기본 회복 + 회피 반사\\ (속력비, 회피 각)};
\node[box=red, right=of fam] (filt) {이산 안전 필터 $F$\\ 35개 후보, 3\,s 전개,\\ 종단 검사, 계획 유지};
\node[box=gray, right=of filt] (plant) {쿼드로터,\\ 장애물, 추격자,\\ 표적};
\node[box=green, below=8mm of filt] (str) {임무 정책\\ $\unom=\sat(\upur+\Delta u_\xi)$};
\draw[->] (p1) -- (fam);
\draw[->] (fam) -- (filt);
\draw[->] (str) -- node[right]{$\unom$} (filt);
\draw[->] (filt) -- node[above]{$u$} (plant);
\draw[->] (plant) |- node[pos=0.75,below]{$o$ (상태, 장면)} (str);
\node[draw, dashed, rounded corners, fit=(filt)(str)(plant), inner sep=5pt,
  label={[font=\scriptsize]below:PS2: 필터를 포함한 환경에서 임무 정책을 학습($r^{\mathrm{PS2}}=r-w_f\|\eta\|^2$)}] {};
\end{tikzpicture}
\caption{전체 구성. 안전 필터는 회복 가능성을, 임무 정책은 표적 접근 성능을 담당한다.}\label{fig:overview}
\end{figure}

\paragraph{재현 경로} 표·그림·본문 수치는 \file{docker\_batch/run\_strike\_paper.sh}가 만든다.
\file{STEPS=tables}는 몇 초 안에 표(\file{.md/.csv/.tsv/.tex}), 그림 7종, 본문 인용 수치(\file{numbers.md})를 만든다.
나머지 단계는 대표 에피소드를 다시 실행해 2차원·3차원 궤적 그림, 상세 신호 그림, 2차원·3차원 GIF를 만든다(\S\ref{sec:results-episodes}).

%======================================================================
\section{문제 설정}\label{sec:problem}

\subsection{기체 모델}
기체 상태는 위치 $p\in\R^3$, 관성좌표 속도 $v\in\R^3$, 자세 쿼터니언 $q\in S^3$로 이루어지고,
입력은 질량으로 나눈 총추력(추력 가속도) $a_T$와 기체 각속도 명령 $\omega$이다.
\begin{equation}
x=(p,v,q),\qquad u=(a_T,\omega),\qquad 0\le a_T\le 4g=39.24\ \mathrm{m/s^2},\qquad |\omega_i|\le 18\ \mathrm{rad/s}.
\label{eq:state}
\end{equation}
시간 간격 $\Delta t=0.02$\,s(50\,Hz)의 명시적 Euler 적분으로 상태를 갱신한다. $R(q)$는 기체 좌표에서 관성좌표로의 회전행렬, $e_3=(0,0,1)^\top$, $\otimes$는 쿼터니언 곱이다.
\begin{align}
p_{k+1}&=p_k+\Delta t\,v_k, \label{eq:dyn-p}\\
v_{k+1}&=v_k+\Delta t\,\big(a_{T,k}\,R(q_k)e_3-g\,e_3\big), \label{eq:dyn-v}\\
\tilde q_{k+1}&=q_k+\tfrac{\Delta t}{2}\,q_k\otimes(0,\omega_k),\qquad q_{k+1}=\tilde q_{k+1}/\|\tilde q_{k+1}\|. \label{eq:dyn-q}
\end{align}
이를 $x_{k+1}=f_d(x_k,u_k)$로 쓴다. 입력은 항상 상자 $\mathcal U$로 포화시킨다($\sat$).
입력의 정규화 척도는 $u_{\mathrm{scale}}=(a_{\max},\omega_{\max},\omega_{\max},\omega_{\max})=(39.24,18,18,18)$이다.
이 모델은 PX4 오프보드 제어가 받는 추력--각속도 인터페이스에 맞춘 운동 모델이며, 모터 동특성, 지연, 공력, 추정 오차는 포함하지 않는다.

\subsection{경기장과 장면 생성}\label{sec:arena}
\begin{itemize}
\item \textbf{경기장}: $x\in[0,40]$\,m, $|y|\le10$\,m, 기준 바닥 고도 $z_f=1$\,m, 상한 고도 $z_{\max}=10$\,m. 기체 반경 $r_d=0.3$\,m.
\item \textbf{장애물}: 높은 기둥 3개. 확률 0.5로 원기둥(반경 $\mathcal U[0.8,2.0]$\,m), 아니면 직사각 기둥(반폭 $\mathcal U[1.5,3.0]\times\mathcal U[1.0,2.5]$\,m, 회전 $\pm30^\circ$). 중심 $x\in[7,31]$\,m.
  장애물 간 간격 3.5\,m 이상, 출발점과 표적에서 4\,m 이상, 옆벽에서 1\,m 이상 떨어지게 하고, 장애물과 옆벽에서 1.5\,m 넘게 떨어진 0.5\,m 격자점만으로 출발점과 표적을 이을 수 있는 장면만 사용한다.
\item \textbf{출발}: $x=2$\,m, $y\in[-3,3]$\,m, $z\in[2,4]$\,m, 속력 $\mathcal U[4,6]$\,m/s(표적 방향 $\pm15^\circ$), 기울기 $10^\circ$ 이하.
\item \textbf{표적}: $x\in[32,37]$\,m, $|y|\le5$\,m, $z\in[1.6,2.5]$\,m에서 시작해 수평으로 $\mathcal U[0.2,0.8]$\,m/s의 일정 속도로 떠돌며, $x\in[32,39]$, $|y|\le9$의 영역 경계에서 반사된다.
\item \textbf{추격자}: 출발점--표적 직선의 $35$--$85\,\%$ 지점에서 옆으로 $5$--$9$\,m 떨어진 곳(높이 $2$--$4$\,m)에서 시작한다.
\end{itemize}

\paragraph{추격자 운동(환경)} 추격자 $j$의 위치 $c_j$, 속도 $\nu_j$는 장애물을 밀어내는 순수 추적으로 움직인다.
\begin{align}
w&=\frac{p-c_j}{\|p-c_j\|}+2\varpi_c\begin{bmatrix}\nabla\sdf_O(c_{j,xy})\\0\end{bmatrix},\qquad
\varpi_c=\clip\!\Big(\frac{d_{\mathrm{rep}}-\sdf_O(c_{j,xy})}{d_{\mathrm{rep}}},0,1\Big),\quad d_{\mathrm{rep}}=2\ \mathrm m,\\
\nu_{j,k+1}&=\Pi_{v_c}\!\Big(\nu_{j,k}+\Delta t\,\Pi_{a_c}\big(\tfrac{v_c w/\|w\|-\nu_{j,k}}{\Delta t}\big)\Big),\qquad c_{j,k+1}=c_{j,k}+\Delta t\,\nu_{j,k+1},
\end{align}
여기서 $\Pi_r(\cdot)$는 노름을 $r$ 이하로 줄이는 연산이며, 높이는 $[z_f+r_c,z_{\max}]$로 제한한다.
필터는 이 운동을 알지 못하고, 거리 감소에 가장 불리한 \emph{가속도 제한 없는 순수 추적}(식~\eqref{eq:pursue})으로 예측한다.

\subsection{장애물의 부호 거리}
장애물 $i$는 중심 $c_i$, 반폭 $b_i=(b_x,b_y)$, 회전 $\psi_i$, 모서리 반경 $\rho_i$로 표현한다(원기둥은 $b_i=0$, $\rho_i$=반경, 직사각 기둥은 $\rho_i=0$).
$\ell=R(\psi_i)^\top(p_{xy}-c_i)$로 두면
\begin{equation}
\sdf_i(p_{xy})=\big\|\max(|\ell|-b_i,0)\big\|_2+\min\!\big(\max(|\ell_x|-b_x,\ |\ell_y|-b_y),\,0\big)-\rho_i,\qquad
\sdf_O=\min_i\sdf_i ,
\label{eq:sdf}
\end{equation}
이며 $|\cdot|$와 $\max$는 성분별로 적용한다. 밖에서 양수, 표면에서 0, 안에서 음수이다.

\subsection{종료 조건과 성과 지표}\label{sec:outcomes}
매 스텝 다음 사건을 검사한다(같은 스텝에서 접촉이 우선).
\begin{itemize}
\item \textbf{성공}(strike): $\|p-p_t\|<0.6$\,m.
\item \textbf{접촉}(contact): 장애물 $\sdf_O(p_{xy})<r_d$, 바닥 $z<z_f$, 추격자 $\|p-c_j\|<r_c+r_d=0.7$\,m 중 하나.
\item \textbf{영역 이탈}(out): $x\notin[-1,41]$, $|y|>11$ 또는 $z>z_{\max}$. 임무 실패이지만 접촉으로 세지 않는다.
\item \textbf{시간초과}(time-up): 15\,s(750스텝) 동안 위 사건이 없음. 10\,s 결과는 같은 에피소드를 10\,s에서 절단한 것이다.
\end{itemize}

%======================================================================
\section{안전 집합}\label{sec:safe}

\subsection{기하 제약}
\begin{equation}
h_f=z-z_f,\qquad h_o=\sdf_O(p_{xy})-r_d,\qquad h_w=d_{\mathrm{wall}}(p_{xy})-r_d,\qquad h_z=z_{\max}-z-r_d,
\label{eq:geo}
\end{equation}
여기서 $d_{\mathrm{wall}}$은 네 외벽($x=0$, $x=40$, $y=\pm10$)까지의 최소 거리다.
바닥 조건은 과제의 고도 정의($z\ge z_f$)를 그대로 따르므로 $r_d$를 빼지 않는다.

\subsection{추격자 여유도 $h_c$의 유도}\label{sec:hc}
추격자 $j$에 대해 거리와 시선(추격자에서 기체로 향하는 단위벡터)을
\begin{equation}
R_j=\|p-c_j\|,\qquad \hat e_j=\frac{p-c_j}{R_j}
\end{equation}
로 두면 기체가 거리를 벌리는 속도는 $\hat e_j^\top v$이다. 추격자는 거리를 최대 $v_c$로 줄이므로(기체를 정면으로 향할 때) 최악의 접근 속도는
\begin{equation}
s_j=v_c-\hat e_j^\top v,\qquad s_j^+=\max(s_j,0)
\label{eq:closing}
\end{equation}
이다. 기체가 반응 지연 $\delta$ 뒤에 유효 이탈 가속도 $a_e$로 거리를 벌리는 속도를 높인다고 하자.
지연 동안에는 거리를 최대 $s_j^+\delta$ 잃는다. 그 뒤 접근 속도가 $s(t)=s_j^+-a_e t$로 줄어 0이 될 때까지 잃는 거리는
\begin{equation}
\int_0^{s_j^+/a_e}(s_j^+-a_e t)\,dt=\frac{(s_j^+)^2}{2a_e}
\qquad\Big(\text{등가로}\ 0=(s_j^+)^2-2a_e d_{\mathrm{brake}}\Big)
\end{equation}
이다. 충돌 반경과 두 손실을 현재 거리에서 빼면 추격자 여유도를 얻는다.
\begin{equation}
\boxed{\;h_{c,j}=R_j-(r_c+r_d)-s_j^+\,\delta-\frac{(s_j^+)^2}{2a_e}\;},\qquad \delta=0.2\ \mathrm s,\quad a_e=5\ \mathrm{m/s^2}.
\label{eq:hc}
\end{equation}
$a_e$는 기체의 최대 가속도보다 훨씬 작게 잡은 값이다. 요구 거리 $d_{\mathrm{req}}=r_c+r_d+s^+\delta+(s^+)^2/(2a_e)$의 예는 다음과 같다.

\begin{center}\small
\begin{tabular}{lccl}\toprule
상황 & $s^+$ [m/s] & $d_{\mathrm{req}}$ [m] & 구성 \\\midrule
추격자로부터 $v_c$ 이상으로 멀어짐 & 0 & 0.7 & 충돌 반경만 \\
정지(호버) & 2 & 1.5 & $0.7+0.4+0.4$ \\
추격자 쪽으로 6\,m/s 정면 접근 & 8 & 8.7 & $0.7+1.6+6.4$ \\\bottomrule
\end{tabular}
\end{center}

\noindent
같은 거리 4\,m라도 정지 중이면 $h_c=2.5$\,m이지만 6\,m/s로 정면 접근 중이면 $h_c=-4.7$\,m이다.
즉 아직 충돌하지 않은 상태도 안전 집합 밖일 수 있고, 이것이 거리만 보는 조건과의 차이다.

\paragraph{열린 공간에서의 불변성(근사)} 지연 뒤 기체가 $\hat e_j$ 방향으로 $a_e$의 가속을 유지하고 추격자가 정면으로 $v_c$로 다가오면 $\dot R_j=-s_j$, $\dot s_j=-a_e$이므로 $s_j>0$에서
\begin{equation}
\dot h_{c,j}=\dot R_j-\dot s_j\,\delta-\frac{s_j\dot s_j}{a_e}=-s_j+a_e\delta+s_j=a_e\delta\ \ge0
\end{equation}
이다. 따라서 막힌 곳이 없으면 $\{h_c\ge0\}$은 ``도피''에 대해 유지된다. 시선 회전에 의한 항(크기 $\lesssim v_c^2/(4R)$)은 $a_e$를 실제 최대보다 작게 잡아 흡수한다고 본다.
장애물과 벽이 도피 방향을 막을 수 있다는 점이 이 식만으로는 부족한 이유이고, 필터가 회복 궤적을 직접 검사하는 이유이다.

\subsection{안전 집합}
장면 정보 $\zeta$(장애물과 추격자 상태)에 대해
\begin{equation}
\Safe(\zeta)=\big\{x:\ h_f\ge0,\ h_o\ge0,\ h_w\ge0,\ h_z\ge0,\ h_{c,j}\ge0\ \ \forall j\big\}.
\label{eq:safe}
\end{equation}

\subsection{비교: 구형 근사(초기 방식)}\label{sec:ball}
초기 방식은 추격자를 현재 위치에 고정된 반경 $R_{\mathrm{ball}}=r_c+r_d+v_cT=3.7$\,m($T=1.5$\,s)의 구로 보고 $\|p-c_j\|\ge R_{\mathrm{ball}}$을 요구한다.
$T$초 안에 추격자가 닿을 수 있는 영역을 덮지만, 그 뒤 추격자의 위치와 기체의 접근 속도를 반영하지 못한다.
Phase I 학습(\S\ref{sec:phase1})은 이 구형 근사를 사용하며, 구형 근사 필터는 비교 실험(표~\ref{tab:table1_15s}, \ref{tab:table3_ablation})에 쓰인다.

%======================================================================
\section{기준 집합: 호버 LQR}\label{sec:base}

\subsection{선형화와 LQR}
호버 목표점 $a$ 주변에서 오차 상태 $e=(p-a,\ v,\ \phi)\in\R^9$를 쓴다. $\phi=2\,\mathrm{sgn}(q_w)\,q_{xyz}$는 작은 회전벡터(자세 좌표)이다.
호버 입력 $\ustar=(g,0,0,0)$ 주변의 연속시간 선형 모델은
\begin{equation}
\dot p=v,\qquad \dot v_x=g\,\phi_y,\qquad \dot v_y=-g\,\phi_x,\qquad \dot v_z=\Delta a_T,\qquad \dot\phi=\omega
\end{equation}
이고, Euler 이산화 $A_d=I+\Delta tA$, $B_d=\Delta tB$에 대해 이산 LQR을 설계한다.
\begin{equation}
e_{k+1}=A_de_k+B_d\Delta u_k,\qquad \Delta u_k=-Ke_k,\qquad K=(R_u+B_d^\top PB_d)^{-1}B_d^\top PA_d,\qquad V(e)=e^\top Pe.
\label{eq:lqr}
\end{equation}
$P$는 이산 Riccati 방정식의 해이다. 가중치는 $Q=\mathrm{diag}(2,2,8,\ 0.01,0.01,0.5,\ 8,8,2)$, $R_u=\mathrm{diag}(0.06,0.15,0.15,0.12)$이다.

\subsection{정지점의 재설정}
속도가 있는 기체는 제자리에 멈출 수 없으므로, $V$를 최소로 만드는 목표점을 해석적으로 구한다.
$r=(v,\phi)$로 두고 $P$를 위치와 나머지로 분할하면
\begin{equation}
V=(p-a)^\top P_{pp}(p-a)+2(p-a)^\top P_{pr}r+r^\top P_{rr}r .
\end{equation}
$\partial V/\partial a=-2P_{pp}(p-a)-2P_{pr}r=0$에서
\begin{equation}
a^\star(x)=p+Dr,\quad D=P_{pp}^{-1}P_{pr};\qquad
V_{\min}(r)=\min_aV=r^\top Sr,\quad S=P_{rr}-P_{rp}P_{pp}^{-1}P_{pr}
\label{eq:anchor}
\end{equation}
이다($S$는 Schur 여행렬). 이때 $e=(-Dr,r)$이므로 재설정 LQR은 $u=\ustar-\tilde Kr$, $\tilde K=K_r-K_pD$로, 어디서든 목표점 없이 정지하는 제어기가 된다.
설계값에서 $D$의 속도 성분은 약 0.385\,s로, 수평 속도 1\,m/s마다 정지점이 0.385\,m 앞에 놓인다.

\subsection{호버 타원체와 기준 집합}
$\{e:e^\top Pe\le c\}$에서 위치 성분의 방향 $n$으로의 최대 변위는 Cauchy--Schwarz 부등식으로 $\sqrt{c\,n^\top[P^{-1}]_{pp}n}$이다. 따라서
\begin{equation}
\rho_{xy}=\sqrt{c_B\,\lambda_{\max}\big([P^{-1}]_{xy}\big)}=1.22\ \mathrm m,\qquad
z_{\mathrm{ext}}=\sqrt{c_B\,[P^{-1}]_{zz}}=0.63\ \mathrm m,\qquad c_B=50 .
\label{eq:extent}
\end{equation}
필터가 쓰는 기준 집합은 정적 환경의 호버 회복 가능 상태로, 추격자 조건을 포함하지 않는다(추격자는 회피 반사와 종단 조건으로 다룬다, \S\ref{sec:family}, \ref{sec:terminal}).
\begin{equation}
\Base=\left\{x\ \middle|\
\begin{aligned}
&h_f,h_o,h_w,h_z\ge0,\qquad r^\top Sr\le c_B,\\
&\sdf_O(a^\star_{xy})\ge\rho_{xy}+r_d,\qquad d_{\mathrm{wall}}(a^\star_{xy})\ge\rho_{xy}+r_d,\\
&a^\star_z-z_{\mathrm{ext}}\ge z_f,\qquad z_{\max}-a^\star_z\ge\rho_{xy}+r_d
\end{aligned}\right\}.
\label{eq:base}
\end{equation}
즉 충분히 느리고($V_{\min}\le c_B$), LQR이 정착하는 동안 움직일 수 있는 타원체 전체가 장애물·바닥·벽에서 떨어져 있다.
천장 조건은 구현대로 $\rho_{xy}$를 쓰며, $\rho_{xy}\ge z_{\mathrm{ext}}$이므로 보수적이다.
Phase I의 기준 집합 $\BaseP$는 벽과 천장 대신 추격자 구 조건 $\|a^\star-c_j\|\ge R_{\mathrm{ball}}+\rho$($\rho=1.22$\,m, 3차원 범위)를 쓴다.

\subsection{$c_B$의 인증}
목표점을 $a_0$에 고정한 LQR $u=\ustar-K(p-a_0,r)$이 $\{V\le c_B\}$를 유지하려면 $c_B$가 다음 세 한계보다 작아야 한다.
\begin{itemize}
\item 입력 한계: 타원체 위에서 $\max|K_ie|=\sqrt{c\,K_iP^{-1}K_i^\top}\le m_i$이므로 $c_U=\min_i m_i^2/(K_iP^{-1}K_i^\top)=56.1$ ($m=(g,18,18,18)$).
\item 자세 좌표: $|\phi|<2$이면 좌표가 일대일이므로 $c_{\mathrm{chart}}=4/\lambda_{\max}([P^{-1}]_{\phi\phi})=102.1$.
\item 비선형 감소: 비선형 Euler 모델에서 $\{V=c\}$ 위 4000개 표본의 한 스텝 감소 $V^+/V<1$을 확인한다. $c=\min(c_U,c_{\mathrm{chart}})=56.1$에서 이미 모든 표본이 감소하므로(최악 방향 0.3\,m/s$^2$ 바람 포함) $c_{\mathrm{Lyap}}\ge56.1$이다. 감소하지 않으면 이분법으로 줄인다.
\end{itemize}
$c_B=50$에서 표본 최대 $V^+/V$는 0.980(바람 포함 0.986)이다. $\Base$ 안에서 수평 속력은 2.67\,m/s, 수직 속력은 2.12\,m/s, 기울기는 $52.7^\circ$ 이하이다.
표본 검사이므로 연속 상태공간 전체의 엄밀한 증명은 아니다.

\subsection{호버 latch}
계획이 $\Base$에 들어가고 회피 반사가 꺼져 있으면($\sigma=0$, \S\ref{sec:family}) 그때의 목표점 $a_0$을 고정하고 고정 목표점 LQR로 호버한다.
인증이 다루는 제어기가 바로 이 고정 목표점 LQR이기 때문이다. 매 스텝 재설정하면 목표점이 밀리미터 단위로 흔들려, 경계에서 들어온 상태가 다시 $\Base$ 밖으로 나가는 일이 생긴다.
고정된 호버는 $V(p-a_0,r)\le c_B$인 동안 종단 조건으로 인정된다.

%======================================================================
\section{Phase I: 안전 도착 회복 정책}\label{sec:phase1}

\paragraph{정책} LQR 입력에 잔차를 더한다.
\begin{equation}
\pi_\theta(x,\zeta)=\sat\big(u_{\mathrm{LQR}}(x)+\Delta u_\theta(o)\big),\qquad u_{\mathrm{LQR}}(x)=\ustar-\tilde Kr .
\label{eq:phase1-policy}
\end{equation}
관측 $o\in\R^{98}$은 $z-z_f$, $v$, $q$(8개), 가까운 순서의 장애물 슬롯 8개(최근접점 벡터 2, 중심 오프셋 2, 반폭 2, $\cos2\psi$, $\sin2\psi$, 모서리 반경, 존재 여부: 각 10개), 추격자 구 슬롯 2개(중심 오프셋 3, 반경 $R_{\mathrm{ball}}$, 존재 여부: 각 5개)로 구성된다.
슬롯 순서는 회복 전개를 시작할 때의 순서로 고정하여 전개 중 입력이 불연속으로 바뀌지 않게 한다. $\Safe$와 $\BaseP$의 소속은 $o$의 정확한 함수이다.

\paragraph{목적함수와 Bellman 관계} 에피소드는 표본 (상태, 장면)에서 시작한 회복 전개 하나이며, $\BaseP$ 도달(성공), $\Safe$ 이탈(실패), 112스텝(2.24\,s) 중 먼저 일어나는 사건에서 끝난다.
\begin{equation}
J(\theta)=\mathbb E\big[\beta^{\tau_B}\,\mathbf 1\{\tau_B<\tau_{\mathrm{out}}\}\big],\qquad \beta=0.99 .
\label{eq:sa}
\end{equation}
장면을 포함한 상태 $y=(x,\zeta)$와 결합 전이 $y^+=G(y,u)$에 대해 최적값은
\begin{equation}
W^\star(y)=1\ (y\in\BaseP\cap\Safe),\qquad W^\star(y)=0\ (y\notin\Safe),\qquad W^\star(y)=\beta\max_{u\in\mathcal U}W^\star(G(y,u))\ \ (\text{그 밖})
\end{equation}
를 만족하고, $W^\star(y)=\beta^{d(y)}$($d$는 안전하게 도착하는 최소 스텝 수)이다. TD3의 critic 목표는 이 관계를 그대로 쓴다.
\begin{equation}
y_{\mathrm{TD}}=b(o)+\beta\,c(o)\,\min_{i=1,2}Q'_i\big(o',\mu'(o')\big),\qquad b=\mathbf 1\{\BaseP\},\ c=\mathbf 1\{\Safe\setminus\BaseP\},
\end{equation}
여기서 $Q'(o',\cdot)$는 $o'\in\BaseP$이면 1, $o'\notin\Safe$이면 0으로 둔다. 엔트로피 항은 critic 목표에 넣지 않는다(도착 순서 ``빨리 $>$ 늦게 $>$ 실패''를 보존하기 위함).
actor 손실은 $-\mathbb E[\min_iQ_i(o,\Delta u_\theta(o))]+0.2\,\mathbb E\big[\tfrac14\|\Delta u/u_{\mathrm{scale}}\|^2\big]$이고, 둘째 항이 잔차를 LQR 쪽으로 붙잡는다. actor는 LQR 모방($\Delta u\approx0$)으로 초기화한다.

\paragraph{학습 분포} 일반(속력 $\le10$\,m/s, 기울기 $\le45^\circ$, 장애물 0--6개, 추격자 구 0--2개), 정지 경로 막힘(속력 $\ge3$\,m/s, LQR 정지점 부근에 장애물), 추격자 근접(속력 $\ge3$\,m/s, LQR 정지점 부근에 추격자 구), 호버 근방(느리고 정지점이 비어 있음) 영역을 가중치 $(0.4,0.2,0.2,0.2)$로 섞는다.
평가에는 저속 상태 근처에 추격자 구가 있는 저속 탈출 영역을 추가로 쓴다(가중치 0). 영역마다 5만 개 시작점을 미리 뽑고, 검증 1000개와 시험 2000개는 따로 뽑는다.

\paragraph{평가 지표와 결과} 시험 상태가 $C_N$에 속하는 비율 $\hat m$을 쓴다. $C_N$은 $N=T/\Delta t=75$스텝(1.5\,s) 안에 $\Safe$를 벗어나지 않고 $\BaseP$에 처음 도달하는 상태의 집합이다.
학습 회복 정책은 가중 평균 $70.1\pm0.3\,\%$(세 시드)로 호버 LQR의 $43.8\,\%$보다 높다(표~\ref{tab:table4_phase1}, 그림~\ref{fig:phase1}).
이 평가는 추격자를 정지한 구로 근사한 Phase I 조건이며, 최종 필터의 이동 추격자 성능과는 구분해야 한다.

%======================================================================
\section{회복 정책군}\label{sec:family}

\paragraph{기본 회복 입력} 기준 집합 안이거나 호버가 고정되어 있으면 고정 목표점 LQR, 밖이면 학습 회복 정책(또는 호버 LQR 변형에서는 재설정 LQR)을 쓴다.
\begin{equation}
\ubase^{(j)}(x)=\begin{cases}\ustar-K(p-a_0,r) & x\in\Base\ \text{또는 호버 고정},\\ \pi_\theta(x,\zeta_{\mathrm{view}}) & \text{그 밖(학습 변형)},\\ u_{\mathrm{LQR}}(x) & \text{그 밖(LQR 변형)}.\end{cases}
\end{equation}
호버가 고정되기 전에는 $a_0=a^\star(x)$이므로 첫 줄은 재설정 LQR과 같다. 학습 회복 정책의 장면 입력 $\zeta_{\mathrm{view}}$는 계획을 만든 시점의 장면(가까운 장애물 8개와 추격자 구)으로 고정한다. Phase I에서 학습한 범위(장애물 10\,m, 추격자 12\,m 이내)만 넣는다.

\paragraph{회피 반사} 최소 추격자 여유도 $h_{\min}=\min_jh_{c,j}$에 따라 회피 입력을 섞는다.
\begin{equation}
\sigma(x)=\clip\!\Big(\frac{h_{\mathrm{off}}-h_{\min}(x)}{h_{\mathrm{off}}-h_{\mathrm{on}}},0,1\Big),\qquad
\pi_b^{(j)}(x)=\sat\big((1-\sigma)\,\ubase^{(j)}(x)+\sigma\,\uesc^{(j)}(x)\big),\qquad h_{\mathrm{on}}=0.5,\ h_{\mathrm{off}}=2.0\ \mathrm m.
\label{eq:reflex}
\end{equation}
$h_{\min}\ge2$\,m이면 기본 회복 입력만, $h_{\min}=1.25$\,m이면 절반씩, $h_{\min}\le0.5$\,m이면 회피 입력만 쓴다. 반사가 작동하면($\sigma>0$) 호버 고정을 푼다.

\paragraph{회피 입력} 가장 위험한 추격자 $j^\star=\arg\min_jh_{c,j}$의 반대 방향 $\hat e_{xy}$를 변형 각 $\vartheta_j$만큼 돌리고, 정지점 $a^\star$ 주변의 장애물과 벽의 바깥 법선으로 민다.
\begin{align}
d&=\nrm\Big(R(\vartheta_j)\,\hat e_{xy}+2\varpi_o\,n_o+2\textstyle\sum_{i=1}^4\varpi_{w,i}\,n_{w,i}\Big),\qquad
\varpi=\clip\!\Big(\frac{d_{\mathrm{push}}-(\mathrm{dist}(a^\star_{xy})-\rho_{xy}-r_d)}{d_{\mathrm{push}}},0,1\Big),\\
v_{\mathrm{des}}&=\big(v_{\mathrm{esc},j}\,d,\ 1.5\,(z_{\mathrm{ref}}-z)\big),\qquad z_{\mathrm{ref}}=\clip(z,\ z_f+1,\ z_{\max}-3),
\end{align}
여기서 $n_o=\nabla\sdf_O(a^\star_{xy})/\|\nabla\sdf_O\|$, $n_{w,i}$는 벽의 안쪽 법선, $\mathrm{dist}$는 각 장애물·벽까지의 거리, $d_{\mathrm{push}}=1.5$\,m이다.
속도 루프는 $a=k_v(v_{\mathrm{des}}-v)$(수평 성분은 $\|a_{xy}\|\le g\tan45^\circ$로 제한), $a_{\mathrm{des}}=(a_{xy},a_z+g)$에서
$a_T=a_{\mathrm{des}}^\top z_b$, $\omega=k_{\mathrm{att}}R^\top(z_b\times\hat a_{\mathrm{des}})-k_\psi\psi\,e_3$로 입력을 만들고 포화시킨다($k_v,k_{\mathrm{att}},k_\psi)=(2,10,3)$, $z_b=R(q)e_3$, $\psi$는 요각).
추격자 여유가 충분할 때는 학습 회복 정책이 장애물을 처리하고, $h_c$가 줄어들수록 회복 정책은 ``$v_{\mathrm{esc}}$로 도피''로 바뀐다. 이는 열린 공간에서 $h_c$를 줄이지 않는 행동이다(\S\ref{sec:hc}).

\paragraph{일곱 변형} 회피 속력비 $v_{\mathrm{esc}}/v_c$와 회피 각 $\vartheta$의 조합으로 회복 정책을 늘린다(표~\ref{tab:variants}).
필터는 어느 하나라도 인증하는 후보를 받아들이므로, 회복 가능 집합은 각 변형 회복 가능 집합의 합집합이 된다.

\begin{table}[h]\centering\small
\caption{필터의 회복 정책 일곱 변형. 번호는 동률일 때의 우선순위이다.}\label{tab:variants}
\begin{tabular}{cccc}\toprule
번호 $j$ & 기본 회복 입력(기준 집합 밖) & $v_{\mathrm{esc}}/v_c$ & $\vartheta$ \\\midrule
0 & 학습 회복 정책 $\pi_\theta$ & 2 & $0^\circ$ \\
1 & 재설정 호버 LQR & 2 & $0^\circ$ \\
2, 3 & 학습 회복 정책 & 2 & $+60^\circ$, $-60^\circ$ \\
4 & 학습 회복 정책 & 3.5 & $0^\circ$ \\
5, 6 & 학습 회복 정책 & 3.5 & $+60^\circ$, $-60^\circ$ \\\bottomrule
\end{tabular}
\end{table}

%======================================================================
\section{이산 다중 회복 안전 필터}\label{sec:filter}



\subsection{후보 입력}
임무 정책의 명목 입력 $\unom$을 각 회복 정책의 현재 입력과 섞어 후보를 만든다.
\begin{equation}
u^{(j)}_\lambda=\sat\big(\lambda\,\unom+(1-\lambda)\,\pi_b^{(j)}(x)\big),\qquad \lambda\in\Lambda=\{1,0.75,0.5,0.25,0\},\quad j=0,\dots,6 .
\label{eq:cand}
\end{equation}
$7\times5=35$개 후보를 검사한다($\lambda=1$에서는 일곱 개가 같은 입력이다).
이는 입력 집합 전체를 최적화하는 것이 아니라 명목 입력과 각 회복 입력을 잇는 선분 위의 유한한 점을 검사하는 것이다.

\subsection{회복 궤적의 전개와 추격자 예측}
후보를 한 스텝 실행한 다음 상태에서 시작해 그 변형의 회복 정책을 $N=150$스텝(3\,s) 전개한다. 추격자는 예측 기체를 향해 $v_c$로 직진한다고 가정한다(거리 감소에 가장 불리한 운동).
\begin{align}
x_1&=f_d\big(x,u^{(j)}_\lambda\big),\qquad x_{i+1}=f_d\big(x_i,\ \pi_b^{(j)}(x_i;c_i)\big),\qquad i=1,\dots,N, \label{eq:rollout}\\
c_{\ell,i+1}&=c_{\ell,i}+\min\!\big(v_c\Delta t,\ \|p_i-c_{\ell,i}\|\big)\,\frac{p_i-c_{\ell,i}}{\|p_i-c_{\ell,i}\|},\qquad c_{\ell,0}=c_\ell(\text{현재}),\qquad c_{\ell,i,z}\ge z_f+r_c . \label{eq:pursue}
\end{align}
$\ell$은 추격자 번호이다. 회피 반사는 예측 추격자 위치 $c_i$를 보고, 학습 회복 정책의 장면 입력은 계획 시점의 것으로 고정된다. 종단 첨자는 $H=N+1$이다.

\subsection{궤적 조건(신규 계획 여유)}
새 계획은 모든 예측 노드 $i=1,\dots,H$에서 여유를 더한 안전 집합에 있어야 한다.
\begin{equation}
\min\{h_f,h_o,h_w,h_z\}(x_i)\ge m_o=0.15\ \mathrm m,\qquad \min_\ell h_{c,\ell}(x_i;c_{\ell,i})\ge m_c=0.2\ \mathrm m .
\label{eq:fresh}
\end{equation}
여유는 반사 방향의 재조준과 실제 추격자가 예측과 다르게 움직이는 것에 대비한 것이다. 계획 유지(\S\ref{sec:commit})의 재검사는 여유 없이($m_o=m_c=0$) 안전 집합만 요구한다.

\subsection{종단 집합}\label{sec:terminal}
종단 상태 $x_H$는 호버하거나 직선 제동으로 멈출 수 있어야 한다.
\begin{align}
\Term_{\mathrm{hover}}&=\big\{x_H:\ \text{호버 고정이면 } V(p_H-a_0,r_H)\le c_B,\ \text{아니면 } x_H\in\Base\big\},\\
\Term_{\mathrm{flee}}&=\big\{x_H:\ \sdf_O(p_{\mathrm{brake}})\ge r_d,\ d_{\mathrm{wall}}(p_{\mathrm{brake}})\ge r_d\big\},\qquad
p_{\mathrm{brake}}=p_{H,xy}+\frac{v_{H,xy}\|v_{H,xy}\|}{2a_e},\\
\Term&=\big(\Term_{\mathrm{hover}}\cup\Term_{\mathrm{flee}}\big)\cap\big\{\min_\ell h_{c,\ell}(x_H;c_{\ell,H})\ge m_c\big\}.
\label{eq:terminal}
\end{align}
$p_{\mathrm{brake}}$는 수평 속도 방향으로 $a_e$로 제동할 때의 정지점이다. 따라서 종단이 항상 $\Base$인 것은 아니며, 도피 중인 상태도 제동 공간이 남아 있으면 허용한다.
회복 정책 $j$의 유한시간 회복 가능 집합은
\begin{equation}
C^{(j)}=\big\{(x,\zeta):\ x_i\in\Safe^{+}\ (i=1,\dots,H),\ x_H\in\Term\big\},\qquad C=\bigcup_{j=0}^{6}C^{(j)}
\end{equation}
이며($\Safe^+$는 식~\eqref{eq:fresh}의 여유를 더한 안전 집합), 후보 실행 직후 상태가 $C^{(j)}$에 속하는지를 검사한다.
합집합은 각 집합을 포함하므로 같은 검사 조건에서 한 회복 정책만 쓰는 경우보다 선택지가 줄지 않는다. 다만 모든 회피 가능 상태를 찾는다는 뜻은 아니다.

\subsection{선택 규칙}
인증된 조합의 집합을 $\mathcal A$라 할 때 명목 입력의 비중이 가장 큰 후보를 실행한다.
\begin{equation}
(\lambda^\star,j^\star)\in\arg\max_{(\lambda,j)\in\mathcal A}\lambda,\qquad u=u^{(j^\star)}_{\lambda^\star}\quad(\text{동률이면 작은 }j).
\label{eq:select}
\end{equation}
변형마다 회복 입력이 다르므로 이는 $\|u-\unom\|$의 최소화와 일반적으로 같지 않다. ``명목 입력이 회복 가능성을 유지하면 그대로 통과시키고, 필요할 때만 회복 쪽으로 움직인다''는 최소 수정의 이산 형태이다.
그림~\ref{fig:snapshot}은 실제 에피소드에서 필터가 한 스텝에 검사한 내용을 다시 계산한 것이다. 명목 입력($\lambda=1$)은 앞의 장애물 때문에 모든 회복 정책에서 거부되고, 회복 정책 0과 $\lambda=0.75$의 조합이 인증되어 실행되었다(평가 기록과 일치).
\IfFileExists{\snapdir/snapshot_ep845_step091.png}{%
\begin{figure}[H]\centering
\includegraphics[width=\linewidth]{\snapdir/snapshot_ep845_step091.png}
\caption{필터의 한 스텝 결정(시드 0, 에피소드 845, $t=1.8$\,s). 왼쪽: 일곱 회복 정책의 3\,s 계획($\lambda=0$)과 끝나는 방식(점선 원: 호버 타원체, $\times$: 제동 정지점), 선택한 후보의 계획(굵은 파랑)과 추격자 예측(초록 점선), 추격자의 $h_c=0$ 거리(옅은 원). 오른쪽: 후보 35개의 인증 결과(초록: 인증, 빨강: 궤적 조건 위반, 주황: 종단 조건 위반)와 선택한 칸.}\label{fig:snapshot}
\end{figure}}{}

\subsection{계획 유지와 추격자 예측 재생}\label{sec:commit}
학습 회복 정책의 입력(가까운 장애물 순서, 범위 안의 장애물)은 스텝마다 바뀔 수 있어 직전에 인증한 계획이 다음 스텝에 달라질 수 있다.
그래서 필터는 기억 $m=(\zeta_{\mathrm{view}},a_0,\text{호버 고정 여부},j_c,\text{추격자 슬롯 순서},\hat c)$를 저장한다. $\hat c$는 인증한 계획이 예측한 다음 스텝의 추격자 위치이다.
$\mathcal A=\varnothing$이면 저장한 변형 $j_c$의 회복 입력을 \emph{저장한 장면 입력과 예측 추격자 위치} $\hat c$로 계산해 실행한다.
모델이 정확하면 이 입력은 인증된 회복 궤적의 다음 입력과 정확히 같다(committed continuation). 이 입력도 실제 추격자 위치에 대해 여유 없이 다시 검사하고, 실패하면 그대로 실행하되 ``인증 없음''으로 기록한다.
예측 추격자 위치는 식~\eqref{eq:pursue}로 한 스텝씩 진행한다.

\subsection{알고리즘}
\begin{algorithm}[H]
\caption{이산 다중 회복 안전 필터 $F$ (한 제어 스텝)}\label{alg:filter}
\KwIn{상태 $x$, 장면(장애물, 추격자 $c$), 명목 입력 $\unom$}
\KwMem{$m=(\zeta_{\mathrm{view}},a_0,\text{고정 여부},j_c,\text{슬롯 순서},\hat c)$}
\KwOut{실행 입력 $u$, 갱신된 기억 $m'$, 정보 $(\lambda^\star,\text{인증 여부})$}
현재 장면에서 새 장면 입력 $\zeta_{\mathrm{new}}$를 만들고, 추격자 예측 $c_1\gets$ 식~\eqref{eq:pursue}를 한 스텝 적용\;
\For{$j=0,\dots,6$}{
  $\pi_b^{(j)}(x)$ 계산(식~\eqref{eq:reflex})\;
  \For{$\lambda\in\{1,0.75,0.5,0.25,0\}$}{
    $u^{(j)}_\lambda\gets\sat(\lambda\unom+(1-\lambda)\pi_b^{(j)}(x))$, \ $x_1\gets f_d(x,u^{(j)}_\lambda)$\;
    회복 정책 $j$를 $N=150$스텝 전개(식~\eqref{eq:rollout}--\eqref{eq:pursue})\;
    모든 노드가 식~\eqref{eq:fresh}를 만족하고 $x_H\in\Term$이면 $(\lambda,j)$를 $\mathcal A$에 추가\;
  }
}
\eIf{$\mathcal A\neq\varnothing$}{
  $(\lambda^\star,j^\star)\gets\arg\max_{\mathcal A}\lambda$ (동률이면 작은 $j$), \ $u\gets u^{(j^\star)}_{\lambda^\star}$\;
  $m'\gets(\zeta_{\mathrm{new}},a_0',\text{고정 여부}',j^\star,\text{새 슬롯 순서},c_1)$\;
}{
  $u\gets\pi_b^{(j_c)}(x;\zeta_{\mathrm{view}},\hat c)$ \tcp*{직전 인증 계획의 다음 입력}
  여유 없이 실제 추격자에 대해 재검사하고, 실패하면 ``인증 없음''으로 기록\;
  $m'\gets(\zeta_{\mathrm{view}},a_0',\text{고정 여부}',j_c,\text{슬롯 순서},\Pi(\hat c))$\;
}
\KwRet{$u,\ m'$}
\end{algorithm}

\subsection{보장되는 것과 보장되지 않는 것}\label{sec:guarantee}
\begin{keybox}{명제 (가정 아래의 진술)}
가정: (A1) 기체 모델이 정확하고 외란이 없다. (A2) 각 추격자의 속력은 $v_c$ 이하이다. (A3) 3\,s 예측 구간에서 속력 $v_c$의 순수 추적이 $h_c$에 대해 최악의 경우이다(거리에 대해서는 성립, 시선 방향에 대해서는 근사).

이때 실행되는 모든 입력은 인증된 후보이거나 직전 인증 계획의 이어짐이다. 인증된 계획은 3\,s 동안 상태를 여유를 둔 $\Safe$ 안에 두고 $\Term$에서 끝난다.
추격자 예측 재생 덕분에 이어짐은 인증된 회복 궤적의 꼬리를 정확히 재현하며, 실제 추격자 위치에 대해 다시 검사된다.
\end{keybox}
남는 간극은 $\Term$이 추격자에 대해 불변 집합이 아니라는 점이다. 호버는 결국 추격자에게 따라잡히고, 도피는 열린 공간이 남아 있는 동안만 안전하다.
또한 이어짐을 재검사할 때 예측 구간이 한 스텝 늘어나므로, 그 마지막 상태가 $\Term$을 만족하지 못해 인증을 잃을 수 있다.
따라서 이 필터는 3\,s 앞을 인증하고 20\,ms마다 다시 계획하는 장치이며, 이동 추격자에 대한 무한시간 불변성의 증명은 아니다.
실제 측정에서는 인증 상실이 PS2 정책 에피소드의 1.5\,\%에서 나타났다(표~\ref{tab:table2_final}).

\subsection{계산량}
한 제어 스텝에서 후보 35개와 이어짐 7개, 즉 42개의 150스텝 전개(제어 스텝당 약 6300회의 회복 입력 계산)를 수행하며, 학습 변형은 그때마다 98차원 입력의 신경망을 평가한다.
GPU(RTX 4050 노트북)에서 1000개 에피소드를 묶어 평가하면 15\,s 에피소드 전체가 약 350\,s 걸린다(묶음 기준 약 0.46\,ms/기체-스텝).
단일 기체의 실시간 최악 실행 시간은 측정하지 않았다.

%======================================================================
\section{비교 필터: bCBF-QP}\label{sec:bcbf}

PS2-RL의 backup CBF QP를 같은 집합에 적용한 비교 구성이다. 기체와 추격자를 포함한 상태를 $y=(p,v,q,c_1,\dots,c_C)$, 결합 동역학을 $\dot y=f(y)+g(y)u$,
고정된 회복 정책(반사 포함)의 흐름을 $\varphi(\tau,y)$, 그 벡터장을 $f_\pi$라 하자.
민감도 $Q_i=\partial\varphi(\tau_i,y)/\partial y$는 Euler 전개의 정확한 Jacobian으로 전파한다.
절대 예측 시각 $t_i$를 고정하면 $\tau_i=t_i-t$가 시간에 따라 줄어들므로, 매끄러운 구간에서
\begin{equation}
\frac{d}{dt}h\big(\varphi(t_i-t,y(t))\big)=\nabla h(\varphi_i)\big[Q_i\big(f(y)+g(y)u\big)-f_\pi(\varphi_i)\big]
\end{equation}
이다. 둘째 항이 남은 예측 시간이 줄어드는 데서 나온다. 이를 $-\alpha h(\varphi_i)$ 이상으로 제한한 선형 제약을 노드 $\tau_i$(1.5\,s 구간, 3스텝 간격)와 모든 안전 함수 $h\in\{h_o,h_f,h_w,h_z,h_c\}$에 대해 둔다($\alpha=4$).
종단 노드에서는 기준 집합 함수(반사가 작동 중이면 제동점 함수)와 $h_c$에 같은 형태의 제약을 $\alpha_B=2$로 둔다.
\begin{equation}
\min_{u,\ s\ge0}\ \tfrac12\sum_kw_k(u_k-u_{\mathrm{nom},k})^2+\tfrac12w_ss^2\quad\text{s.t. 완화된 CBF 제약, }u\in\mathcal U,\qquad w=(1,0.1,0.1,1),\ w_s=10^6 .
\end{equation}
roll/pitch 각속도의 가중치가 작아 추력보다 기울기로 경로를 바꾸는 해를 선호한다.
회복 입력 $u=\pi_b(x)$가 제약을 $10^{-2}$보다 크게 위반하면(정규화된 행 기준) QP를 건너뛰고, QP 출력은 다시 이산 필터(최종 회복 정책군)를 통과시킨다.
따라서 이 구성의 접촉 억제는 QP 단독의 결과가 아니며, 선형화·이산화·slack 때문에 표준 연속시간 CBF 불변성 정리를 그대로 적용할 수 없다.

%======================================================================
\section{임무 정책}\label{sec:striker}

\paragraph{기준 접근 제어} 표적 위치 $p_t$와 속도 $v_t$에 대해
\begin{equation}
v_{\mathrm{des}}=v_{\mathrm{cr}}\frac{p_t-p}{\|p_t-p\|}+v_t,\qquad v_{\mathrm{cr}}=6\ \mathrm{m/s}
\end{equation}
를 회피 반사와 같은 속도·자세 루프($k_v=2.5$, 기울기 $50^\circ$, $k_{\mathrm{att}}=10$, $k_\psi=3$)로 추종해 $\upur$를 만든다. 장애물과 추격자를 고려하지 않는다.

\paragraph{잔차 정책}
\begin{equation}
\unom=\sat\big(\upur(x)+\Delta u_\xi(o)\big),\qquad |\Delta a_T|\le0.1\,a_{\max}=3.92\ \mathrm{m/s^2},\quad |\Delta\omega_i|\le0.25\,\omega_{\max}=4.5\ \mathrm{rad/s}.
\label{eq:residual}
\end{equation}
관측은 135차원이다: 회복 정책 형식의 장면 관측 98개, 표적 오프셋/10, 표적 속도, 추격자 속도, 경기장 내 위치(14개),
유도 특징 19개(시선 방향 3, $d/10$, $\log(d+0.1)$, 상대 속도/5 3, 접근 속도/5, 시선각 변화율/2 3, 요 좌표계의 상대 위치/10 3과 상대 속도/5 3, 표적 높이/2),
그리고 보정 대상인 기준 입력 $(\upur-\ustar)/u_{\mathrm{scale}}$(4개).

\paragraph{보상} 표적 거리 $d_k=\|p_k-p_{t,k}\|$에 대해
\begin{equation}
r_k=w_p(d_{k-1}-d_k)-w_t-\frac{w_a}{4}\Big\|\frac{\Delta u_k}{u_{\mathrm{scale}}}\Big\|^2+100\cdot\mathbf 1_{\text{성공}}-150\cdot\mathbf 1_{\text{접촉 또는 이탈}},\qquad w_p=2,\ w_t=w_a=0.05 .
\label{eq:reward}
\end{equation}
거리 항은 $w_p(d_0-d_{\mathrm{end}})$로 합쳐지고, 시간 비용은 빠른 성공을 장려한다. 시간초과는 critic 목표에서 종료로 보지 않는다(부트스트랩 유지).

\paragraph{학습} SAC, $\gamma=0.995$, 보상 척도 0.1, 병렬 환경 64개, 배치 256, 재생 버퍼 $10^6$, 학습률 $3\times10^{-4}$, $\tau=0.005$, 온도 초기값 0.1(하한 $10^{-3}$, 목표 엔트로피 $-4$), 은닉층 256$\times$2(ELU).
처음 20만 스텝은 actor를 고정하고 critic만 학습한다(critic warm-up). actor 손실에 고정점 벌점 $0.2\,\mathbb E\big[\tfrac14\|\Delta u/b\|^2\big]$($b$는 잔차 한계)을 더한다.
1000만 스텝 학습하며, 25만 스텝마다 128개 에피소드로 평가해 반환값이 가장 큰 체크포인트를 쓴다. 처음부터 학습한 정책은 실패했고(성공 0, 이후 0.41--0.49), 기준 접근 제어의 잔차로 두어야 안정적으로 학습되었다.

%======================================================================
\section{PS2: 필터를 포함한 임무 학습}\label{sec:ps2}

필터를 환경의 일부로 두고 임무 정책을 다시 학습한다.
\begin{align}
u_{\mathrm{nom},k}&=\sat\big(\upur(x_k)+\Delta u_\xi(o_k)\big),\qquad (u_k,m_{k+1})=F\big(x_k,\zeta_k,m_k,u_{\mathrm{nom},k}\big),\qquad x_{k+1}=f_d(x_k,u_k), \label{eq:ps2}\\
r^{\mathrm{PS2}}_k&=r_k-w_f\|\eta_k\|^2,\qquad \eta_k=\frac{u_k-u_{\mathrm{nom},k}}{u_{\mathrm{scale}}},\qquad w_f=0.5 . \label{eq:ps2-reward}
\end{align}
정책은 자신이 제안한 입력이 그대로 실행된다고 가정하지 않고, 필터 개입 이후의 상태와 보상으로 학습한다. 필터가 자주 막는 행동보다 실제로 진전이 생기고 수정 비용이 작은 행동이 유리해진다.
필터 기억 $m$까지 포함한 전이 $(x,\zeta,m)\mapsto(x^+,\zeta^+,m^+)$는 정책에서 보면 환경의 일부이므로, 이산 필터를 미분하지 않고도 off-policy SAC로 학습할 수 있다.
다만 $m$은 관측되지 않으므로 관측이 완전한 상태라고 주장하지 않는다. 필터 출력을 복사하는 모방학습이 아니라, 수정된 전이에서 같은 강화학습 목표를 다시 최적화하는 것이다.

\paragraph{설정} 같은 시드의 잔차 정책 최고 체크포인트에서 시작하여, 처음 10만 스텝은 actor를 고정하고 critic을 재적응시키며, 온도는 시작 값으로 고정한다(학습률 0).
배치할 필터(최종 회복 정책군, 추격자 예측 재생)와 같은 필터로 300만 스텝 학습하고(학습 = 배치), 25만 스텝마다 64개 에피소드로 평가해 반환값이 가장 큰 체크포인트를 쓴다.
PS2 정책은 필터와 함께 쓰도록 학습한 정책이며, 필터 없이 실행하면 성공률이 잔차 정책보다 낮다(62.1\,\% 대 74.0\,\%). 필터가 장애물에서 비켜 주는 것에 의존하기 때문이다.

\paragraph{학습 시간} RTX 4050 노트북 GPU에서 시드당 Phase I 약 10분(1000만 스텝), 잔차 정책 28--34분(1000만 스텝), PS2 72--76분(300만 스텝)이 걸렸다.

%======================================================================
\section{평가 방법}\label{sec:protocol}

\paragraph{절차} 시드 $s\in\{0,1,2\}$마다 전체 파이프라인(Phase I $\to$ 잔차 정책 $\to$ PS2)을 그 시드로 학습한다.
모든 구성을 같은 1000개 평가 장면(장면 시드 9001, 학습 장면 풀과 별도)으로 평가하며, 장면마다 출발 교란과 표적·추격자 초기 조건까지 같다.
에피소드는 15\,s이고, 10\,s 결과는 같은 에피소드를 10\,s에서 절단한 것이다(10\,s까지 끝나지 않은 에피소드는 시간초과로 센다).
각 시드는 그 시드의 임무 정책을 학습한 추격자 설정(원래 추격자)으로 평가한다.

\paragraph{지표} 성공률, 접촉률(장애물·바닥·추격자), 시간초과율, 영역 이탈률, 성공 에피소드의 평균 도달 시간,
평균 명목 유지율 $\bar\lambda$(에피소드 스텝 평균 $\lambda^\star$, 새 인증이 없는 스텝은 0),
인증 상실률(한 스텝 이상 유효한 인증 계획이 없었던 에피소드의 비율). 인증 상실은 접촉과 다른 지표이다.

\paragraph{통계} 시드별 비율의 평균 $\pm$ 표준편차를 보고하고, 시드를 합친 횟수 $k/n$에 대해 양측 95\,\% Clopper--Pearson 구간을 쓴다.
\begin{equation}
p_L=\mathrm{B}^{-1}\!\big(\tfrac{\alpha}{2};\,k,\,n-k+1\big),\qquad p_U=\mathrm{B}^{-1}\!\big(1-\tfrac{\alpha}{2};\,k+1,\,n-k\big),\qquad \alpha=0.05,
\end{equation}
여기서 $\mathrm B^{-1}$는 베타 분포의 분위수이다. $k=0$이면 $p_U=1-(\alpha/2)^{1/n}$이고, 단측 95\,\% 상한은 $1-0.05^{1/n}\approx3/n$이다.
$n=1000$에서 0.30\,\%(단측)와 0.37\,\%(양측), $n=3000$에서 0.10\,\%와 0.12\,\%이다.
모든 구성이 같은 장면을 쓰므로 두 구성의 비교는 \emph{쌍 비교}이다. 구성 A만 성공한 장면 수 $b$와 B만 성공한 장면 수 $c$로 정확 McNemar 검정을 한다.
\begin{equation}
p_{\mathrm{McNemar}}=\min\Big\{1,\ 2\sum_{i=0}^{\min(b,c)}\binom{b+c}{i}2^{-(b+c)}\Big\}.
\end{equation}
구성끼리 비교하는 표(표~\ref{tab:table2_final}, \ref{tab:table3_ablation})는 모든 행이 평가된 공통 시드만 쓴다.

%======================================================================
\section{결과}\label{sec:results}

\begin{warnbox}{수치의 기준}
아래 표와 그림은 업로드된 평가 결과(시드 0 전체, 시드 1 일부, 시드 2 미평가)에서 만든 것이다. 각 행의 ``시드 $\times$ 장면'' 열이 사용한 에피소드 수를 나타낸다.
세 시드 평가를 마친 뒤 \file{STEPS=tables bash docker\_batch/run\_strike\_paper.sh}로 다시 만들고 이 문서를 다시 컴파일하면 표와 그림이 갱신된다.
\end{warnbox}

\subsection{주요 비교}
\tabin{table1_10s}
\tabin{table1_15s}

\begin{figure}[H]\centering
\figin{\figpath{fig_outcomes.pdf}}
\caption{구성별 결과 비율(왼쪽 10\,s, 오른쪽 15\,s 제한). 막대 안의 수는 성공률(\%)이다.}\label{fig:outcomes}
\end{figure}

\paragraph{필터의 효과} 같은 임무 정책에서 최종 필터는 접촉을 거의 없앤다. 잔차 정책은 26.1\,\%(521/2000)에서 0/2000으로, 기준 접근 제어는 57.5\,\%에서 1/2000으로 줄었다.
대신 우회와 정지 때문에 10\,s 시간초과가 늘어난다(잔차 정책 0\,\%$\to$15.5\,\%). 15\,s 제한에서는 시간초과 대부분이 늦은 성공으로 바뀐다.

\paragraph{PS2의 효과} 같은 최종 필터를 쓸 때 PS2 정책은 잔차 정책보다 성공률이 높다(10\,s 84.5$\to$90.4\,\%, 15\,s 94.4$\to$97.6\,\%).
같은 장면 쌍 비교에서 PS2만 성공한 장면이 172개, 잔차 정책만 성공한 장면이 54개로(10\,s) 차이는 유의하다($p<10^{-6}$, 표~\ref{tab:table5_paired}).
평균 도달 시간도 6.53\,s에서 6.11\,s로 줄었다. 필터를 고정하고 임무 정책만 바꾼 비교이므로 필터 포함 재학습의 효과를 보여 주는 핵심 근거이다.

\begin{figure}[H]\centering
\begin{subfigure}{0.49\linewidth}\centering\figin{\figpath{fig_strike_time.pdf}}\caption{누적 성공률}\label{fig:strike-time}\end{subfigure}\hfill
\begin{subfigure}{0.49\linewidth}\centering\figin{\figpath{fig_safety.pdf}}\caption{추격자 최소 여유 거리의 누적 분포}\label{fig:safety}\end{subfigure}
\caption{(a) 시간에 따른 누적 성공률(음영은 시드를 합친 95\,\% 구간). 점선은 10\,s. (b) 에피소드별 추격자 최소 여유 거리 $\min_t(\|p-c\|-r_c-r_d)$의 누적 분포. 필터를 쓴 구성에서 0.2\,m의 계단은 추격자가 바짝 따라오는 동안 필터가 $h_c$를 신규 계획 여유 $m_c=0.2$\,m에 유지한 에피소드들이다(멀어지는 중에는 $h_c$가 거리 여유와 같다).}
\end{figure}

\subsection{최종 필터 선택과 구성 요소}
\tabin{table2_final}
\tabin{table3_ablation}
\tabin{table5_paired}

\begin{figure}[H]\centering
\figin{\figpath{fig_ablation.pdf}}
\caption{구성 요소 제거 실험(PS2 정책, 15\,s). 오차 막대는 95\,\% Clopper--Pearson 구간이다.}\label{fig:ablation}
\end{figure}

\begin{itemize}
\item \textbf{회피 변형}: 단일 회피 필터는 접촉 1.7\,\%(17/1000)로, 최종 필터(1/1000)와의 쌍 비교에서 차이가 유의하다($p=3.1\times10^{-5}$). 성공률도 낮다(15\,s 95.7\,\% 대 98.4\,\%).
\item \textbf{bCBF-QP}: 성공률은 비슷하지만 최종 필터가 10\,s와 15\,s 모두 높고($p=0.0015$, $0.043$), 평균 도달 시간이 짧다(5.93\,s 대 6.33\,s, 공통 시드 기준). bCBF는 명목 유지율이 더 높지만(0.97) 이후 이산 검사에서 인증을 더 자주 잃는다(2.5\,\% 대 1.5\,\%).
\item \textbf{학습 회복 정책}: 회복 정책을 호버 LQR로 바꾸면 15\,s 성공률이 82--84\,\%로 크게 떨어지고 시간초과가 15--17\,\%로 늘어난다(시드 2개와 공통 시드 기준). Phase I 회복 정책이 필터의 허용 범위를 넓히는 것이 성능의 주요 원인이다.
\item \textbf{구형 근사}: 추격자를 정지한 구로 보면 접촉이 26.3\,\%로, 상대 접근 속도를 반영한 $h_c$가 필수적이다.
\item \textbf{계획 유지의 추격자 예측 재생과 호버 latch}: 1000개 장면에서 성공·접촉의 차이는 관측되지 않았다(98.6\,\%, 98.3\,\% 대 98.4\,\%). 두 요소는 성능보다 이어짐이 인증된 계획을 정확히 재현하고 인증 대상 제어기와 일치하게 하는 일관성을 위한 것이다.
\end{itemize}

\begin{figure}[H]\centering
\figin{\figpath{fig_filter.pdf}}
\caption{필터 개입. 왼쪽: 에피소드 평균 명목 유지율 $\bar\lambda$의 누적 분포(1에서의 계단은 필터가 한 번도 수정하지 않은 에피소드). 오른쪽: 인증 상실 에피소드 비율.}\label{fig:filter}
\end{figure}

\subsection{학습 과정과 Phase I}
\begin{figure}[H]\centering
\figin{\figpath{fig_training.pdf}}
\caption{학습 중 평가 성공률(시드별). 원은 평가에 사용한 체크포인트. 필터 없이 학습한 잔차 정책은 시드 1, 2에서 500--650만 스텝 이후 붕괴했으며, 평가에는 최고 반환값 체크포인트(시드 0: 850만, 시드 1: 125만, 시드 2: 175만 스텝)를 썼다. PS2 학습은 필터 안에서 안정적이다.}\label{fig:training}
\end{figure}
\tabin{table4_phase1}
\begin{figure}[H]\centering
\figin[0.55\linewidth]{\figpath{fig_phase1.pdf}}
\caption{Phase I 회복 성공률(1.5\,s 안에 안전하게 기준 집합 도달). 학습 회복 정책은 세 시드 평균과 표준편차.}\label{fig:phase1}
\end{figure}

\subsection{대표 에피소드}\label{sec:results-episodes}
\file{run\_strike\_paper.sh}는 시드 0의 평가에서 필터가 장애물 충돌을 막은 장면, 추격자 접촉을 막은 장면, 추격자 근접 회피, 유일한 접촉 장면,
기준 제어가 실패한 장면 등을 자동으로 고르고(\file{episode\_picks.json}), 각 장면을 여러 구성으로 다시 실행해 다음을 만든다.
\begin{itemize}
\item 2차원 궤적 비교(구성별 패널), 3차원 궤적(속력 색), 3차원 비교 그림;
\item 상세 신호 그림: 표적 거리, 장애물·벽·바닥·천장 여유, $h_c$와 반사 가중치 $\sigma$, 속력과 고도, $\lambda^\star$와 인증 상태(신규/이어짐/인증 없음), 사용한 회복 변형, 입력과 명목 입력;
\item 2차원 GIF와 회전하는 3차원 GIF(추격자 $h_c$ 요구 거리 구 표시);
\item 필터 스냅샷(\file{scripts/strike\_filter\_snapshot.py}): 처음 명목 입력을 줄인 스텝, 이전 계획을 이어 간 스텝, 인증을 잃은 스텝에서 필터의 검사를 다시 계산한 그림(그림~\ref{fig:snapshot}).
\end{itemize}
\IfFileExists{\episodedir/overview_main.png}{%
\begin{figure}[H]\centering
\includegraphics[width=\linewidth]{\episodedir/overview_main.png}
\caption{대표 에피소드(시드 0)의 수평 궤적. 실선은 최종 필터를 쓴 PS2 정책과 잔차 정책, 점선은 필터 없는 PS2 정책이다. 별은 성공, X는 접촉, 회색 원은 시간초과.}\label{fig:episodes-overview}
\end{figure}}{\IfFileExists{\episodedir/overview.png}{%
\begin{figure}[H]\centering
\includegraphics[width=\linewidth]{\episodedir/overview.png}
\caption{대표 에피소드(시드 0)의 수평 궤적.}\label{fig:episodes-overview}
\end{figure}}{}}

\paragraph{유일한 접촉 사례(시드 0, 에피소드 845)} 기체는 큰 원기둥과 위쪽 벽 사이의 좁은 공간($y\approx9$\,m)으로 들어간 뒤 약 6\,s부터 새 인증을 찾지 못해
이전 계획(호버 LQR + 회피 변형)을 이어 갔고, 약 7.5\,s부터는 그 계획도 재검사를 통과하지 못했다(인증 없음).
추격자가 다가오면서 회피 반사가 작동했지만 벽과 장애물이 도피 방향을 막아 9.2\,s에 접촉했다. 같은 장면에서 bCBF 구성도 같은 방식으로 실패했다.
이 사례는 \S\ref{sec:guarantee}의 남는 간극, 즉 호버 종단이 다가오는 추격자에 대해 불변 집합이 아니고 막힌 공간에서는 도피도 보장되지 않는다는 점이 실제로 나타난 경우이다.
인증을 잃은 7.5\,s의 검사를 다시 계산하면(그림~\ref{fig:snapshot-lost}) 35개 후보 모두 궤적 조건을 위반했다. 위쪽 벽($y=10$\,m) 쪽으로 움직이던 기체에서 일곱 회복 정책의 3\,s 계획이 모두 벽을 넘었기 때문이다.
Phase I 회복 정책은 벽 없이 학습되었고, 회피 반사는 추격자 여유가 $h_{\mathrm{off}}=2$\,m 이상이면 섞이지 않는다. 벽을 포함한 Phase I 학습이 남은 개선 방향이다.
\IfFileExists{\snapdir/snapshot_ep845_step374.png}{%
\begin{figure}[H]\centering
\includegraphics[width=\linewidth]{\snapdir/snapshot_ep845_step374.png}
\caption{인증을 잃은 시점의 검사(에피소드 845, $t=7.5$\,s). 일곱 회복 정책의 계획이 모두 위쪽 벽을 넘어 새 인증이 없고, 필터는 이전 계획을 이어 실행했다.}\label{fig:snapshot-lost}
\end{figure}}{}
\IfFileExists{\episodedir/ep845/compare2d.png}{%
\begin{figure}[H]\centering
\includegraphics[width=\linewidth]{\episodedir/ep845/compare2d.png}
\caption{에피소드 845를 여러 구성으로 다시 실행한 수평 궤적(CPU 재실행).}\label{fig:ep845-2d}
\end{figure}}{}
\IfFileExists{\episodedir/ep845/detail_ps2_ours.png}{%
\begin{figure}[H]\centering
\includegraphics[width=0.86\linewidth]{\episodedir/ep845/detail_ps2_ours.png}
\caption{에피소드 845, PS2 정책 + 최종 필터의 상세 신호. 위에서부터 표적 거리, 기하 여유, $h_c$와 반사 가중치 $\sigma$, 속력과 고도, $\lambda^\star$와 인증 상태(노랑: 이전 계획 유지, 빨강: 인증 없음), 사용한 회복 변형, 입력과 명목 입력.}\label{fig:ep845-detail}
\end{figure}}{}

%======================================================================
\section{한계와 적용 범위}\label{sec:limits}
\begin{enumerate}
\item \textbf{종단 집합의 불변성}: 호버 기준 집합은 정적 장애물에 대한 회복 수단이다. 추격자가 계속 다가오면 호버는 영구히 안전하지 않고, 도피 종단도 유한시간 예측의 조건이다. 정적 기준 집합에 기대는 원래 이론의 보장을 그대로 옮길 수 없다.
\item \textbf{추격자 모델}: 예측은 현재 시선 방향의 순수 추적이며, 장애물 우회와 시선각 변화를 포함한 모든 추격 전략의 최악값이라는 증명은 없다. 반응 지연과 이탈 가속도에 기반한 $h_c$도 근사이다.
\item \textbf{모델 오차}: 입력 지연, 공력, 센서 오차, 바람이 있으면 3\,s 검사와 실제 경로가 달라질 수 있다. 모든 오차를 포함하는 강건 경계는 제시하지 않았다.
\item \textbf{시간 이산 검사}: 검사는 20\,ms 노드에서 한다. 노드 사이 연속시간 충돌까지 배제하려면 이동량이나 보간 오차를 이용한 추가 조건이 필요하다.
\item \textbf{인증 상실}: 계획 유지 중에도 인증을 잃을 수 있다(PS2 에피소드의 1.5\,\%). 이 시점은 명제의 전제를 벗어난다.
\item \textbf{통계}: 1000개 장면과 2--3개 시드의 결과이다. 접촉 1/2000은 접촉확률이 0이라는 증명이 아니며 95\,\% 상한은 0.28\,\%이다.
\item \textbf{실기체}: 결과는 시뮬레이션의 성공과 접촉이다. 실시간 최악 실행 시간, PX4 추력·각속도 추종 오차, 실기체 운용 성능은 검증하지 않았다.
\end{enumerate}

%======================================================================
\appendix
\section{하이퍼파라미터}\label{app:hparams}
\begin{center}\small
\begin{tabularx}{\linewidth}{>{\raggedright\arraybackslash}p{3.0cm}>{\raggedright\arraybackslash}X}\toprule
구분 & 값 \\\midrule
모델 & $\Delta t=0.02$\,s, $a_T\in[0,4g]$, $|\omega_i|\le18$\,rad/s, $r_d=0.3$\,m, $z_f=1$\,m, $z_{\max}=10$\,m \\
경기장·장면 & 40$\times$20\,m, 장애물 3개, 추격자 1개($r_c=0.4$\,m, $v_c=2$\,m/s, $a_c=2$\,m/s$^2$, 반발 거리 2\,m), 표적 0.2--0.8\,m/s, 성공 반경 0.6\,m, 15\,s \\
$h_c$ & $\delta=0.2$\,s, $a_e=5$\,m/s$^2$ \\
호버 LQR & $Q=\mathrm{diag}(2,2,8,0.01,0.01,0.5,8,8,2)$, $R_u=\mathrm{diag}(0.06,0.15,0.15,0.12)$, $c_B=50$, $\rho_{xy}=1.22$\,m, $z_{\mathrm{ext}}=0.63$\,m \\
Phase I & TD3, 1000만 스텝, 병렬 64, 배치 256, 버퍼 $10^6$, $\beta=0.99$, 학습률 actor $10^{-4}$/critic $3\times10^{-4}$, $\tau=0.0025$, 정책 지연 2, 탐색 잡음 0.1, 잔차 벌점 0.2, 에피소드 112스텝, 구형 근사 $T=1.5$\,s, 영역 가중치 $(0.4,0.2,0.2,0.2)$ \\
회피 반사 & $h_{\mathrm{on}}=0.5$\,m, $h_{\mathrm{off}}=2.0$\,m, $(k_v,k_{\mathrm{att}},k_\psi)=(2,10,3)$, 기울기 $45^\circ$, $d_{\mathrm{push}}=1.5$\,m, $z_{\mathrm{ref}}\in[z_f+1,z_{\max}-3]$ \\
회복 정책군 & $(v_{\mathrm{esc}}/v_c,\vartheta)\in\{(2,0^\circ)\text{ 학습},(2,0^\circ)\text{ LQR},(2,\pm60^\circ),(3.5,0^\circ),(3.5,\pm60^\circ)\}$ \\
이산 필터 & $\Lambda=\{1,0.75,0.5,0.25,0\}$, 예측 3\,s(150스텝), 신규 계획 여유 $m_o=0.15$\,m·$m_c=0.2$\,m, 종단 $\Term_{\mathrm{hover}}\cup\Term_{\mathrm{flee}}$, 계획 유지(추격자 예측 재생), 호버 latch \\
bCBF-QP & $\alpha=4$, $\alpha_B=2$, 1.5\,s, 3스텝 간격, $w=(1,0.1,0.1,1)$, $w_s=10^6$, 건너뛰기 허용치 $10^{-2}$, 보조 필터 = 최종 이산 필터 \\
기준 접근 제어 & $v_{\mathrm{cr}}=6$\,m/s, $k_v=2.5$, 기울기 $50^\circ$ \\
잔차 정책 & 잔차 한계 추력 10\,\%, 각속도 25\,\%; 관측 135; SAC $\gamma=0.995$, 보상 척도 0.1, 병렬 64, 배치 256, 학습률 $3\times10^{-4}$, $\tau=0.005$, 온도 0.1(하한 $10^{-3}$), critic warm-up 20만, 고정점 벌점 0.2, 1000만 스텝 \\
보상 & $w_p=2$, $w_t=0.05$, $w_a=0.05$, 성공 $+100$, 접촉·이탈 $-150$ \\
PS2 & $w_f=0.5$, 잔차 정책에서 시작, critic 재적응 10만, 온도 고정, 300만 스텝 \\
평가 & 시드 3개, 장면 1000개(시드 9001), 15\,s와 10\,s 절단, Clopper--Pearson 95\,\%, 정확 McNemar \\\bottomrule
\end{tabularx}
\end{center}

\section{기호표}\label{app:symbols}
\begin{center}\small
\begin{tabularx}{\linewidth}{>{\raggedright\arraybackslash}p{3.4cm}X}\toprule
기호 & 의미 \\\midrule
$x=(p,v,q)$, $u=(a_T,\omega)$ & 기체 상태, 입력(추력 가속도, 기체 각속도) \\
$f_d$, $\Delta t$ & 이산 동역학, 시간 간격 \\
$\zeta$, $c_\ell$, $\nu_\ell$ & 장면 정보, 추격자 위치, 추격자 속도 \\
$r_d$, $r_c$, $v_c$, $a_c$ & 기체 반경, 추격자 반경, 추격자 속력·가속도 상한 \\
$h_f,h_o,h_w,h_z$ & 바닥, 장애물, 벽, 천장 여유 \\
$h_{c,j}$, $s_j^+$, $\delta$, $a_e$ & 추격자 여유도, 양의 접근 속도, 반응 지연, 유효 이탈 가속도 \\
$\Safe$, $\Base$, $\BaseP$ & 안전 집합, 필터의 기준 집합, Phase I 기준 집합 \\
$e$, $P$, $K$, $V$ & 호버 오차, Riccati 해, LQR 이득, $V=e^\top Pe$ \\
$a^\star$, $D$, $S$ & 최적 정지점, 정지점 행렬, Schur 여행렬 ($V_{\min}=r^\top Sr$) \\
$c_B$, $\rho_{xy}$, $z_{\mathrm{ext}}$ & 기준 집합 수준, 타원체의 수평·수직 범위 \\
$\pi_\theta$, $\Delta u_\theta$ & Phase I 학습 회복 정책과 그 잔차 \\
$\pi_b^{(j)}$, $\uesc^{(j)}$, $\sigma$ & 회복 정책 변형 $j$, 회피 입력, 반사 가중치 \\
$\Lambda$, $\lambda^\star$, $\bar\lambda$ & 혼합률 집합, 선택한 혼합률, 에피소드 평균 \\
$N$, $H$, $m_o$, $m_c$ & 예측 스텝(150), 종단 첨자($N+1$), 신규 계획 여유 \\
$\Term$, $p_{\mathrm{brake}}$ & 종단 집합, 제동 정지점 \\
$F$, $m$ & 안전 필터, 필터 기억(계획 유지) \\
$\upur$, $\Delta u_\xi$, $\unom$ & 기준 접근 입력, 임무 정책 잔차, 명목 입력 \\
$r_k$, $r^{\mathrm{PS2}}_k$, $\eta_k$, $w_f$ & 보상, PS2 보상, 정규화한 필터 수정량, 수정 비용 가중치 \\\bottomrule
\end{tabularx}
\end{center}

\section{원고 초안과 다른 점}\label{app:draft}
초안을 이 버전의 결과로 바꿀 때 확인할 항목이다.
\begin{enumerate}
\item \textbf{평가 규모}: 초안의 ``단일 시드, 128개 장면'' 수치(44.5, 82.8, 93.8, 100\,\% 등)는 이전 실험이다. 현재 표는 1000개 장면, 시드 2--3개이며 15\,s 에피소드와 10\,s 절단을 함께 보고한다.
\item \textbf{접촉 0건 표현}: PS2 정책 + 최종 필터는 접촉 1/2000(시드 0의 한 장면)이다. ``접촉이 관측되지 않았다'' 대신 ``2000회 중 1회(95\,\% 상한 0.28\,\%)''로 쓴다. 잔차 정책 + 최종 필터는 0/2000이다.
\item \textbf{15\,s 성공률}: 초안의 ``128개 장면 모두 성공(100\,\%)''은 현재 97.6\,\%이다(시간초과 2.4\,\%).
\item \textbf{PS2 정책의 학습 필터}: 초안 표 1은 단일 회피 필터로 학습한 PS2 정책을 썼다. 최종 파이프라인의 PS2 정책은 배치하는 최종 필터(회피 변형, 추격자 예측 재생)로 학습했다.
\item \textbf{bCBF slack 가중치}: 초안 식 (34) 뒤의 $w_s=10$은 $w_s=10^6$이다.
\item \textbf{신규 계획 여유}: 0.15\,m는 장애물과 벽뿐 아니라 바닥과 천장에도 적용된다.
\item \textbf{기준 집합의 추격자 조건}: 필터의 기준 집합 $\Base$에는 추격자 조건이 없고, 추격자는 종단 조건 $h_c\ge0.2$\,m와 회피 반사로 다룬다. Phase I의 기준 집합만 추격자 구 조건을 가진다.
\item \textbf{천장 조건}: 식 (20)의 $z_{\max}-a^\star_z\ge\rho_{\mathrm{top}}+r_d$에서 구현은 $\rho_{\mathrm{top}}=\rho_{xy}$이다(초안과 같음, 보수적).
\item \textbf{제동 정지점}: 식 (30)은 구현에서 수평 성분만 쓴다: $p_{\mathrm{brake}}=p_{H,xy}+v_{H,xy}\|v_{H,xy}\|/(2a_e)$. 이 점이 장애물과 벽에서 $r_d$ 이상 떨어져야 한다.
\item \textbf{추격자 운동}: 환경의 추격자는 속력 2\,m/s와 함께 가속도 2\,m/s$^2$ 상한을 가지며, 장애물 2\,m 이내에서 밀려난다(\S\ref{sec:arena}). 필터의 예측(가속도 제한 없는 순수 추적)과 다르다는 점을 밝혀 둔다.
\item \textbf{Phase I 수치}: 가중 성공률은 세 시드 평균 70.1$\pm$0.3\,\%(초안 69.9\,\%), LQR 43.8\,\%이다.
\item \textbf{이동 추격자 회복 학습 변형}: 초안 표 2의 해당 행은 최종 파이프라인 평가에 포함되지 않은 128개 장면 실험이다. 쓰려면 이전 실험임을 밝히거나 표에서 뺀다.
\item \textbf{수식 표기}: 초안 식 (2)·(6)의 $|\omega_i)$, $|\ell_i)$는 절댓값 $|\omega_i|$, $|\ell_i|$이고, 식 (4)의 괄호는 $v_{k+1}=v_k+\Delta t\,(a_{T,k}R(q_k)e_3-ge_3)$이다. 회피 각 $\vartheta$와 임무 정책 매개변수(초안 $\Delta u_\vartheta$)의 기호가 겹치므로 이 문서는 임무 정책 잔차를 $\Delta u_\xi$로 썼다.
\end{enumerate}

\section{한글(HWP) 수식 편집기 입력}\label{app:hwp}
본문의 주요 수식을 한글 수식 편집기 문법으로 옮긴 것이다(같은 내용: \file{results/strike\_paper/hwp\_equations.txt}). 붙여넣은 뒤 모양을 확인한다.
\IfFileExists{hwp_equations.txt}{\VerbatimInput[fontsize=\scriptsize,breaklines,breakanywhere,firstline=4]{hwp_equations.txt}}{\fbox{hwp\_equations.txt 없음}}

\end{document}
"
\providecommand{\paperdir}{../strike_reduced/final/paper}   % the report's results in this branch
\providecommand{\episodedir}{\paperdir/episodes}
\providecommand{\snapdir}{\paperdir/traj4/snapshots}   % filter snapshots
\documentclass[10pt,a4paper]{article}
\usepackage[margin=2.1cm]{geometry}
\usepackage{fontspec}
\setmainfont{Noto Serif CJK KR}
\setsansfont{Noto Sans CJK KR}
\setmonofont{Noto Sans Mono CJK KR}
\usepackage{amsmath,amssymb,bm,mathtools}
\usepackage{graphicx,float}
\usepackage{booktabs,array,tabularx,multirow}
\usepackage[dvipsnames]{xcolor}
\usepackage{enumitem}
\usepackage{caption,subcaption}
\usepackage{titlesec}
\usepackage{fancyhdr}
\usepackage[most]{tcolorbox}
\usepackage[ruled,vlined,linesnumbered]{algorithm2e}
\usepackage{fvextra}
\usepackage{tikz}
\usetikzlibrary{arrows.meta,positioning,fit,calc}
\usepackage[colorlinks=true,linkcolor=NavyBlue,urlcolor=NavyBlue,citecolor=NavyBlue]{hyperref}

\linespread{1.18}
\setcounter{tocdepth}{1}
\floatplacement{table}{H}
\setlength{\parindent}{1em}
\setlength{\parskip}{2pt}
\emergencystretch=2em
\setlist{nosep,leftmargin=1.5em}
\captionsetup{font=small,labelfont=bf,labelsep=period}
\renewcommand{\figurename}{그림}
\renewcommand{\tablename}{표}
\renewcommand{\contentsname}{목차}
\renewcommand{\refname}{참고 문헌}
\renewcommand{\abstractname}{요약}
\renewcommand{\appendixname}{부록}
\SetAlgorithmName{알고리즘}{알고리즘}{알고리즘 목록}
\SetKwInput{KwIn}{입력}
\SetKwInput{KwOut}{출력}
\SetKwInput{KwMem}{기억}
\SetKw{KwRet}{반환}
\SetKwIF{If}{ElseIf}{Else}{만약}{이면}{그렇지 않고 만약}{그렇지 않으면}{끝}
\SetKwFor{For}{각}{에 대해}{끝}
\titleformat{\section}{\large\bfseries\sffamily\color{NavyBlue}}{\thesection}{0.6em}{}
\titleformat{\subsection}{\normalsize\bfseries\sffamily}{\thesubsection}{0.6em}{}
\titleformat{\paragraph}[runin]{\normalsize\bfseries}{}{}{}[.]
\pagestyle{fancy}\fancyhf{}
\lhead{\small STRIKE: 방법과 결과 상세 정리}
\rhead{\small 2026.\ 10.\ 23}
\cfoot{\small\thepage}
\renewcommand{\headrulewidth}{0.3pt}

\newtcolorbox{keybox}[1]{colback=NavyBlue!4,colframe=NavyBlue!55,fonttitle=\bfseries\sffamily\small,title={#1},
  left=4pt,right=4pt,top=2pt,bottom=2pt,fontupper=\small}
\newtcolorbox{warnbox}[1]{colback=Orange!6,colframe=Orange!70,fonttitle=\bfseries\sffamily\small,title={#1},
  left=4pt,right=4pt,top=2pt,bottom=2pt,fontupper=\small}

\newcommand{\R}{\mathbb{R}}
\newcommand{\sat}{\operatorname{sat}}
\newcommand{\clip}{\operatorname{clip}}
\newcommand{\sdf}{\operatorname{sdf}}
\newcommand{\nrm}{\operatorname{normalize}}
\newcommand{\Safe}{\mathcal{S}}
\newcommand{\Base}{\bar{\mathcal{B}}}
\newcommand{\BaseP}{\mathcal{B}_{\mathrm{I}}}
\newcommand{\Term}{\mathcal{T}}
\newcommand{\unom}{u_{\mathrm{nom}}}
\newcommand{\upur}{u_{\mathrm{pur}}}
\newcommand{\uesc}{u_{\mathrm{esc}}}
\newcommand{\ubase}{u_{\mathrm{base}}}
\newcommand{\ustar}{u^{\star}}
\newcommand{\file}[1]{\texttt{\small #1}}
\newcommand{\figpath}[1]{\paperdir/figures/#1}
\newcommand{\tabin}[1]{\IfFileExists{\paperdir/tables/#1.tex}{\input{\paperdir/tables/#1.tex}}{\begin{center}\fbox{표 파일 없음: #1.tex}\end{center}}}
\newcommand{\figin}[2][\linewidth]{\IfFileExists{#2}{\includegraphics[width=#1]{#2}}{\fbox{그림 파일 없음}}}

\title{\vspace{-1.0cm}\textbf{\sffamily STRIKE: 다중 회복 정책 안전 필터와\\ 필터 포함 임무 학습(PS2)을 이용한 쿼드로터 저고도 표적 접근}\\[4pt]
\large 방법, 수식, 알고리즘, 평가 방법과 결과의 상세 정리 (논문 작성용)}
\author{김중원 \quad (정리 보조: Claude)}
\date{\small 2026년 10월 23일. 대상: 희소 경기장, 원래 추격자(반경 0.4 m, 2 m/s, 2 m/s$^2$)로 학습한 최종 구성}

\begin{document}
\maketitle
\thispagestyle{fancy}

\begin{abstract}
\noindent
장애물과 이동 추격자가 함께 있는 저고도 경기장에서 쿼드로터가 이동 표적의 0.6\,m 이내로 접근하는 문제(STRIKE)를 다룬다.
방법은 두 층으로 이루어진다.
\emph{안전 필터}는 임무 정책의 명목 입력을 일곱 개 회복 정책과 다섯 단계로 혼합한 35개 후보 각각을 한 스텝 실행한 뒤,
그 다음 상태에서 회복 정책을 3\,s 동안 가상으로 전개하여 바닥·장애물·벽·천장 제약과 상대 접근 속도를 반영한 추격자 여유도 $h_c$를 확인하고,
종단 상태가 호버하거나 제동할 수 있는 상태인지 검사한다. 통과한 후보 가운데 명목 입력을 가장 많이 유지하는 입력을 실행하고,
통과 후보가 없으면 직전에 인증한 계획을 그 추격자 예측 그대로 이어서 실행한다.
\emph{임무 정책}은 6\,m/s 기준 접근 제어에 잔차를 더하는 SAC 정책이며, 실제 필터를 포함한 환경에서 필터의 수정량을 벌점으로 받으며 재학습한다(PS2).
1000개 평가 장면과 15\,s 제한에서 PS2 정책과 최종 필터의 조합은 성공률 97.6\,\%, 접촉 2000회 중 1회(0.05\,\%, 95\,\% 상한 0.28\,\%)를 보였고,
같은 필터를 쓰는 잔차 정책(94.4\,\%, 접촉 0/2000)보다 성공률이 높았다(같은 장면 쌍 비교, 정확 McNemar 검정 $p<10^{-6}$).
수치는 현재 업로드된 평가(시드 0 전체, 시드 1 일부)에 기반하며, 세 시드 평가가 끝나면 같은 스크립트로 표와 그림이 갱신된다.
\end{abstract}

\tableofcontents
\vspace{4pt}

%======================================================================
\section{이 문서의 범위와 읽는 법}\label{sec:scope}

이 문서는 학술대회 원고(\emph{STRIKE 다중 회복 정책과 안전 필터 학습을 이용한 쿼드로터의 저고도 표적 접근}) 작성을 위해,
최종 구성의 수식과 알고리즘을 구현(\file{ps2rl/})과 대조하여 정리한 것이다.
모든 수식과 상수는 저장소의 코드와 학습에 사용된 설정 파일(\file{configs.json})에서 확인했다.
원고 초안과 다른 부분은 부록~\ref{app:draft}에, 주요 수식의 한글(HWP) 수식 편집기 입력은 부록~\ref{app:hwp}에 정리했다.

\begin{keybox}{이 버전의 구성 (원고의 ``최종 구성'')}
\begin{itemize}
\item \textbf{경기장}: 40\,m$\times$20\,m, 높은 장애물 3개, 이동 추격자 1개(반경 $r_c=0.4$\,m, 속력 $\le v_c=2$\,m/s, 가속도 $\le a_c=2$\,m/s$^2$), 이동 표적, 15\,s 제한.
\item \textbf{안전 필터}: 이산 다중 회복 필터. 회복 정책 7개(학습 회복 정책 또는 호버 LQR과 회피 반사의 결합), 혼합률 5단계, 3\,s 예측, 신규 계획 여유, 확장 종단 집합, 추격자 예측을 재생하는 계획 유지.
\item \textbf{임무 정책}: 6\,m/s 기준 접근 제어 + SAC 잔차(관측 135차원). PS2: 배치하는 필터를 그대로 포함한 환경에서 300만 스텝 재학습.
\item \textbf{평가}: 세 시드(전체 파이프라인을 시드별로 학습), 1000개 평가 장면(시드 9001), 15\,s 에피소드와 10\,s 절단 결과, Clopper--Pearson 구간과 같은 장면 쌍 비교.
\end{itemize}
\end{keybox}

\begin{figure}[H]\centering
\begin{tikzpicture}[node distance=5mm and 7mm, every node/.style={font=\small},
  box/.style={draw, rounded corners, align=center, minimum height=10mm, inner sep=4pt, fill=#1!8}, box/.default=gray,
  >={Stealth[length=2mm]}]
\node[box=blue] (p1) {Phase I\\ 학습 회복 정책 $\pi_\theta$\\ (안전 도착, TD3)};
\node[box=orange, right=of p1] (fam) {회복 정책군 $\{\pi_b^{(j)}\}_{j=0}^{6}$\\ 기본 회복 + 회피 반사\\ (속력비, 회피 각)};
\node[box=red, right=of fam] (filt) {이산 안전 필터 $F$\\ 35개 후보, 3\,s 전개,\\ 종단 검사, 계획 유지};
\node[box=gray, right=of filt] (plant) {쿼드로터,\\ 장애물, 추격자,\\ 표적};
\node[box=green, below=8mm of filt] (str) {임무 정책\\ $\unom=\sat(\upur+\Delta u_\xi)$};
\draw[->] (p1) -- (fam);
\draw[->] (fam) -- (filt);
\draw[->] (str) -- node[right]{$\unom$} (filt);
\draw[->] (filt) -- node[above]{$u$} (plant);
\draw[->] (plant) |- node[pos=0.75,below]{$o$ (상태, 장면)} (str);
\node[draw, dashed, rounded corners, fit=(filt)(str)(plant), inner sep=5pt,
  label={[font=\scriptsize]below:PS2: 필터를 포함한 환경에서 임무 정책을 학습($r^{\mathrm{PS2}}=r-w_f\|\eta\|^2$)}] {};
\end{tikzpicture}
\caption{전체 구성. 안전 필터는 회복 가능성을, 임무 정책은 표적 접근 성능을 담당한다.}\label{fig:overview}
\end{figure}

\paragraph{재현 경로} 표·그림·본문 수치는 \file{docker\_batch/run\_strike\_paper.sh}가 만든다.
\file{STEPS=tables}는 몇 초 안에 표(\file{.md/.csv/.tsv/.tex}), 그림 7종, 본문 인용 수치(\file{numbers.md})를 만든다.
나머지 단계는 대표 에피소드를 다시 실행해 2차원·3차원 궤적 그림, 상세 신호 그림, 2차원·3차원 GIF를 만든다(\S\ref{sec:results-episodes}).

%======================================================================
\section{문제 설정}\label{sec:problem}

\subsection{기체 모델}
기체 상태는 위치 $p\in\R^3$, 관성좌표 속도 $v\in\R^3$, 자세 쿼터니언 $q\in S^3$로 이루어지고,
입력은 질량으로 나눈 총추력(추력 가속도) $a_T$와 기체 각속도 명령 $\omega$이다.
\begin{equation}
x=(p,v,q),\qquad u=(a_T,\omega),\qquad 0\le a_T\le 4g=39.24\ \mathrm{m/s^2},\qquad |\omega_i|\le 18\ \mathrm{rad/s}.
\label{eq:state}
\end{equation}
시간 간격 $\Delta t=0.02$\,s(50\,Hz)의 명시적 Euler 적분으로 상태를 갱신한다. $R(q)$는 기체 좌표에서 관성좌표로의 회전행렬, $e_3=(0,0,1)^\top$, $\otimes$는 쿼터니언 곱이다.
\begin{align}
p_{k+1}&=p_k+\Delta t\,v_k, \label{eq:dyn-p}\\
v_{k+1}&=v_k+\Delta t\,\big(a_{T,k}\,R(q_k)e_3-g\,e_3\big), \label{eq:dyn-v}\\
\tilde q_{k+1}&=q_k+\tfrac{\Delta t}{2}\,q_k\otimes(0,\omega_k),\qquad q_{k+1}=\tilde q_{k+1}/\|\tilde q_{k+1}\|. \label{eq:dyn-q}
\end{align}
이를 $x_{k+1}=f_d(x_k,u_k)$로 쓴다. 입력은 항상 상자 $\mathcal U$로 포화시킨다($\sat$).
입력의 정규화 척도는 $u_{\mathrm{scale}}=(a_{\max},\omega_{\max},\omega_{\max},\omega_{\max})=(39.24,18,18,18)$이다.
이 모델은 PX4 오프보드 제어가 받는 추력--각속도 인터페이스에 맞춘 운동 모델이며, 모터 동특성, 지연, 공력, 추정 오차는 포함하지 않는다.

\subsection{경기장과 장면 생성}\label{sec:arena}
\begin{itemize}
\item \textbf{경기장}: $x\in[0,40]$\,m, $|y|\le10$\,m, 기준 바닥 고도 $z_f=1$\,m, 상한 고도 $z_{\max}=10$\,m. 기체 반경 $r_d=0.3$\,m.
\item \textbf{장애물}: 높은 기둥 3개. 확률 0.5로 원기둥(반경 $\mathcal U[0.8,2.0]$\,m), 아니면 직사각 기둥(반폭 $\mathcal U[1.5,3.0]\times\mathcal U[1.0,2.5]$\,m, 회전 $\pm30^\circ$). 중심 $x\in[7,31]$\,m.
  장애물 간 간격 3.5\,m 이상, 출발점과 표적에서 4\,m 이상, 옆벽에서 1\,m 이상 떨어지게 하고, 장애물과 옆벽에서 1.5\,m 넘게 떨어진 0.5\,m 격자점만으로 출발점과 표적을 이을 수 있는 장면만 사용한다.
\item \textbf{출발}: $x=2$\,m, $y\in[-3,3]$\,m, $z\in[2,4]$\,m, 속력 $\mathcal U[4,6]$\,m/s(표적 방향 $\pm15^\circ$), 기울기 $10^\circ$ 이하.
\item \textbf{표적}: $x\in[32,37]$\,m, $|y|\le5$\,m, $z\in[1.6,2.5]$\,m에서 시작해 수평으로 $\mathcal U[0.2,0.8]$\,m/s의 일정 속도로 떠돌며, $x\in[32,39]$, $|y|\le9$의 영역 경계에서 반사된다.
\item \textbf{추격자}: 출발점--표적 직선의 $35$--$85\,\%$ 지점에서 옆으로 $5$--$9$\,m 떨어진 곳(높이 $2$--$4$\,m)에서 시작한다.
\end{itemize}

\paragraph{추격자 운동(환경)} 추격자 $j$의 위치 $c_j$, 속도 $\nu_j$는 장애물을 밀어내는 순수 추적으로 움직인다.
\begin{align}
w&=\frac{p-c_j}{\|p-c_j\|}+2\varpi_c\begin{bmatrix}\nabla\sdf_O(c_{j,xy})\\0\end{bmatrix},\qquad
\varpi_c=\clip\!\Big(\frac{d_{\mathrm{rep}}-\sdf_O(c_{j,xy})}{d_{\mathrm{rep}}},0,1\Big),\quad d_{\mathrm{rep}}=2\ \mathrm m,\\
\nu_{j,k+1}&=\Pi_{v_c}\!\Big(\nu_{j,k}+\Delta t\,\Pi_{a_c}\big(\tfrac{v_c w/\|w\|-\nu_{j,k}}{\Delta t}\big)\Big),\qquad c_{j,k+1}=c_{j,k}+\Delta t\,\nu_{j,k+1},
\end{align}
여기서 $\Pi_r(\cdot)$는 노름을 $r$ 이하로 줄이는 연산이며, 높이는 $[z_f+r_c,z_{\max}]$로 제한한다.
필터는 이 운동을 알지 못하고, 거리 감소에 가장 불리한 \emph{가속도 제한 없는 순수 추적}(식~\eqref{eq:pursue})으로 예측한다.

\subsection{장애물의 부호 거리}
장애물 $i$는 중심 $c_i$, 반폭 $b_i=(b_x,b_y)$, 회전 $\psi_i$, 모서리 반경 $\rho_i$로 표현한다(원기둥은 $b_i=0$, $\rho_i$=반경, 직사각 기둥은 $\rho_i=0$).
$\ell=R(\psi_i)^\top(p_{xy}-c_i)$로 두면
\begin{equation}
\sdf_i(p_{xy})=\big\|\max(|\ell|-b_i,0)\big\|_2+\min\!\big(\max(|\ell_x|-b_x,\ |\ell_y|-b_y),\,0\big)-\rho_i,\qquad
\sdf_O=\min_i\sdf_i ,
\label{eq:sdf}
\end{equation}
이며 $|\cdot|$와 $\max$는 성분별로 적용한다. 밖에서 양수, 표면에서 0, 안에서 음수이다.

\subsection{종료 조건과 성과 지표}\label{sec:outcomes}
매 스텝 다음 사건을 검사한다(같은 스텝에서 접촉이 우선).
\begin{itemize}
\item \textbf{성공}(strike): $\|p-p_t\|<0.6$\,m.
\item \textbf{접촉}(contact): 장애물 $\sdf_O(p_{xy})<r_d$, 바닥 $z<z_f$, 추격자 $\|p-c_j\|<r_c+r_d=0.7$\,m 중 하나.
\item \textbf{영역 이탈}(out): $x\notin[-1,41]$, $|y|>11$ 또는 $z>z_{\max}$. 임무 실패이지만 접촉으로 세지 않는다.
\item \textbf{시간초과}(time-up): 15\,s(750스텝) 동안 위 사건이 없음. 10\,s 결과는 같은 에피소드를 10\,s에서 절단한 것이다.
\end{itemize}

%======================================================================
\section{안전 집합}\label{sec:safe}

\subsection{기하 제약}
\begin{equation}
h_f=z-z_f,\qquad h_o=\sdf_O(p_{xy})-r_d,\qquad h_w=d_{\mathrm{wall}}(p_{xy})-r_d,\qquad h_z=z_{\max}-z-r_d,
\label{eq:geo}
\end{equation}
여기서 $d_{\mathrm{wall}}$은 네 외벽($x=0$, $x=40$, $y=\pm10$)까지의 최소 거리다.
바닥 조건은 과제의 고도 정의($z\ge z_f$)를 그대로 따르므로 $r_d$를 빼지 않는다.

\subsection{추격자 여유도 $h_c$의 유도}\label{sec:hc}
추격자 $j$에 대해 거리와 시선(추격자에서 기체로 향하는 단위벡터)을
\begin{equation}
R_j=\|p-c_j\|,\qquad \hat e_j=\frac{p-c_j}{R_j}
\end{equation}
로 두면 기체가 거리를 벌리는 속도는 $\hat e_j^\top v$이다. 추격자는 거리를 최대 $v_c$로 줄이므로(기체를 정면으로 향할 때) 최악의 접근 속도는
\begin{equation}
s_j=v_c-\hat e_j^\top v,\qquad s_j^+=\max(s_j,0)
\label{eq:closing}
\end{equation}
이다. 기체가 반응 지연 $\delta$ 뒤에 유효 이탈 가속도 $a_e$로 거리를 벌리는 속도를 높인다고 하자.
지연 동안에는 거리를 최대 $s_j^+\delta$ 잃는다. 그 뒤 접근 속도가 $s(t)=s_j^+-a_e t$로 줄어 0이 될 때까지 잃는 거리는
\begin{equation}
\int_0^{s_j^+/a_e}(s_j^+-a_e t)\,dt=\frac{(s_j^+)^2}{2a_e}
\qquad\Big(\text{등가로}\ 0=(s_j^+)^2-2a_e d_{\mathrm{brake}}\Big)
\end{equation}
이다. 충돌 반경과 두 손실을 현재 거리에서 빼면 추격자 여유도를 얻는다.
\begin{equation}
\boxed{\;h_{c,j}=R_j-(r_c+r_d)-s_j^+\,\delta-\frac{(s_j^+)^2}{2a_e}\;},\qquad \delta=0.2\ \mathrm s,\quad a_e=5\ \mathrm{m/s^2}.
\label{eq:hc}
\end{equation}
$a_e$는 기체의 최대 가속도보다 훨씬 작게 잡은 값이다. 요구 거리 $d_{\mathrm{req}}=r_c+r_d+s^+\delta+(s^+)^2/(2a_e)$의 예는 다음과 같다.

\begin{center}\small
\begin{tabular}{lccl}\toprule
상황 & $s^+$ [m/s] & $d_{\mathrm{req}}$ [m] & 구성 \\\midrule
추격자로부터 $v_c$ 이상으로 멀어짐 & 0 & 0.7 & 충돌 반경만 \\
정지(호버) & 2 & 1.5 & $0.7+0.4+0.4$ \\
추격자 쪽으로 6\,m/s 정면 접근 & 8 & 8.7 & $0.7+1.6+6.4$ \\\bottomrule
\end{tabular}
\end{center}

\noindent
같은 거리 4\,m라도 정지 중이면 $h_c=2.5$\,m이지만 6\,m/s로 정면 접근 중이면 $h_c=-4.7$\,m이다.
즉 아직 충돌하지 않은 상태도 안전 집합 밖일 수 있고, 이것이 거리만 보는 조건과의 차이다.

\paragraph{열린 공간에서의 불변성(근사)} 지연 뒤 기체가 $\hat e_j$ 방향으로 $a_e$의 가속을 유지하고 추격자가 정면으로 $v_c$로 다가오면 $\dot R_j=-s_j$, $\dot s_j=-a_e$이므로 $s_j>0$에서
\begin{equation}
\dot h_{c,j}=\dot R_j-\dot s_j\,\delta-\frac{s_j\dot s_j}{a_e}=-s_j+a_e\delta+s_j=a_e\delta\ \ge0
\end{equation}
이다. 따라서 막힌 곳이 없으면 $\{h_c\ge0\}$은 ``도피''에 대해 유지된다. 시선 회전에 의한 항(크기 $\lesssim v_c^2/(4R)$)은 $a_e$를 실제 최대보다 작게 잡아 흡수한다고 본다.
장애물과 벽이 도피 방향을 막을 수 있다는 점이 이 식만으로는 부족한 이유이고, 필터가 회복 궤적을 직접 검사하는 이유이다.

\subsection{안전 집합}
장면 정보 $\zeta$(장애물과 추격자 상태)에 대해
\begin{equation}
\Safe(\zeta)=\big\{x:\ h_f\ge0,\ h_o\ge0,\ h_w\ge0,\ h_z\ge0,\ h_{c,j}\ge0\ \ \forall j\big\}.
\label{eq:safe}
\end{equation}

\subsection{비교: 구형 근사(초기 방식)}\label{sec:ball}
초기 방식은 추격자를 현재 위치에 고정된 반경 $R_{\mathrm{ball}}=r_c+r_d+v_cT=3.7$\,m($T=1.5$\,s)의 구로 보고 $\|p-c_j\|\ge R_{\mathrm{ball}}$을 요구한다.
$T$초 안에 추격자가 닿을 수 있는 영역을 덮지만, 그 뒤 추격자의 위치와 기체의 접근 속도를 반영하지 못한다.
Phase I 학습(\S\ref{sec:phase1})은 이 구형 근사를 사용하며, 구형 근사 필터는 비교 실험(표~\ref{tab:table1_15s}, \ref{tab:table3_ablation})에 쓰인다.

%======================================================================
\section{기준 집합: 호버 LQR}\label{sec:base}

\subsection{선형화와 LQR}
호버 목표점 $a$ 주변에서 오차 상태 $e=(p-a,\ v,\ \phi)\in\R^9$를 쓴다. $\phi=2\,\mathrm{sgn}(q_w)\,q_{xyz}$는 작은 회전벡터(자세 좌표)이다.
호버 입력 $\ustar=(g,0,0,0)$ 주변의 연속시간 선형 모델은
\begin{equation}
\dot p=v,\qquad \dot v_x=g\,\phi_y,\qquad \dot v_y=-g\,\phi_x,\qquad \dot v_z=\Delta a_T,\qquad \dot\phi=\omega
\end{equation}
이고, Euler 이산화 $A_d=I+\Delta tA$, $B_d=\Delta tB$에 대해 이산 LQR을 설계한다.
\begin{equation}
e_{k+1}=A_de_k+B_d\Delta u_k,\qquad \Delta u_k=-Ke_k,\qquad K=(R_u+B_d^\top PB_d)^{-1}B_d^\top PA_d,\qquad V(e)=e^\top Pe.
\label{eq:lqr}
\end{equation}
$P$는 이산 Riccati 방정식의 해이다. 가중치는 $Q=\mathrm{diag}(2,2,8,\ 0.01,0.01,0.5,\ 8,8,2)$, $R_u=\mathrm{diag}(0.06,0.15,0.15,0.12)$이다.

\subsection{정지점의 재설정}
속도가 있는 기체는 제자리에 멈출 수 없으므로, $V$를 최소로 만드는 목표점을 해석적으로 구한다.
$r=(v,\phi)$로 두고 $P$를 위치와 나머지로 분할하면
\begin{equation}
V=(p-a)^\top P_{pp}(p-a)+2(p-a)^\top P_{pr}r+r^\top P_{rr}r .
\end{equation}
$\partial V/\partial a=-2P_{pp}(p-a)-2P_{pr}r=0$에서
\begin{equation}
a^\star(x)=p+Dr,\quad D=P_{pp}^{-1}P_{pr};\qquad
V_{\min}(r)=\min_aV=r^\top Sr,\quad S=P_{rr}-P_{rp}P_{pp}^{-1}P_{pr}
\label{eq:anchor}
\end{equation}
이다($S$는 Schur 여행렬). 이때 $e=(-Dr,r)$이므로 재설정 LQR은 $u=\ustar-\tilde Kr$, $\tilde K=K_r-K_pD$로, 어디서든 목표점 없이 정지하는 제어기가 된다.
설계값에서 $D$의 속도 성분은 약 0.385\,s로, 수평 속도 1\,m/s마다 정지점이 0.385\,m 앞에 놓인다.

\subsection{호버 타원체와 기준 집합}
$\{e:e^\top Pe\le c\}$에서 위치 성분의 방향 $n$으로의 최대 변위는 Cauchy--Schwarz 부등식으로 $\sqrt{c\,n^\top[P^{-1}]_{pp}n}$이다. 따라서
\begin{equation}
\rho_{xy}=\sqrt{c_B\,\lambda_{\max}\big([P^{-1}]_{xy}\big)}=1.22\ \mathrm m,\qquad
z_{\mathrm{ext}}=\sqrt{c_B\,[P^{-1}]_{zz}}=0.63\ \mathrm m,\qquad c_B=50 .
\label{eq:extent}
\end{equation}
필터가 쓰는 기준 집합은 정적 환경의 호버 회복 가능 상태로, 추격자 조건을 포함하지 않는다(추격자는 회피 반사와 종단 조건으로 다룬다, \S\ref{sec:family}, \ref{sec:terminal}).
\begin{equation}
\Base=\left\{x\ \middle|\
\begin{aligned}
&h_f,h_o,h_w,h_z\ge0,\qquad r^\top Sr\le c_B,\\
&\sdf_O(a^\star_{xy})\ge\rho_{xy}+r_d,\qquad d_{\mathrm{wall}}(a^\star_{xy})\ge\rho_{xy}+r_d,\\
&a^\star_z-z_{\mathrm{ext}}\ge z_f,\qquad z_{\max}-a^\star_z\ge\rho_{xy}+r_d
\end{aligned}\right\}.
\label{eq:base}
\end{equation}
즉 충분히 느리고($V_{\min}\le c_B$), LQR이 정착하는 동안 움직일 수 있는 타원체 전체가 장애물·바닥·벽에서 떨어져 있다.
천장 조건은 구현대로 $\rho_{xy}$를 쓰며, $\rho_{xy}\ge z_{\mathrm{ext}}$이므로 보수적이다.
Phase I의 기준 집합 $\BaseP$는 벽과 천장 대신 추격자 구 조건 $\|a^\star-c_j\|\ge R_{\mathrm{ball}}+\rho$($\rho=1.22$\,m, 3차원 범위)를 쓴다.

\subsection{$c_B$의 인증}
목표점을 $a_0$에 고정한 LQR $u=\ustar-K(p-a_0,r)$이 $\{V\le c_B\}$를 유지하려면 $c_B$가 다음 세 한계보다 작아야 한다.
\begin{itemize}
\item 입력 한계: 타원체 위에서 $\max|K_ie|=\sqrt{c\,K_iP^{-1}K_i^\top}\le m_i$이므로 $c_U=\min_i m_i^2/(K_iP^{-1}K_i^\top)=56.1$ ($m=(g,18,18,18)$).
\item 자세 좌표: $|\phi|<2$이면 좌표가 일대일이므로 $c_{\mathrm{chart}}=4/\lambda_{\max}([P^{-1}]_{\phi\phi})=102.1$.
\item 비선형 감소: 비선형 Euler 모델에서 $\{V=c\}$ 위 4000개 표본의 한 스텝 감소 $V^+/V<1$을 확인한다. $c=\min(c_U,c_{\mathrm{chart}})=56.1$에서 이미 모든 표본이 감소하므로(최악 방향 0.3\,m/s$^2$ 바람 포함) $c_{\mathrm{Lyap}}\ge56.1$이다. 감소하지 않으면 이분법으로 줄인다.
\end{itemize}
$c_B=50$에서 표본 최대 $V^+/V$는 0.980(바람 포함 0.986)이다. $\Base$ 안에서 수평 속력은 2.67\,m/s, 수직 속력은 2.12\,m/s, 기울기는 $52.7^\circ$ 이하이다.
표본 검사이므로 연속 상태공간 전체의 엄밀한 증명은 아니다.

\subsection{호버 latch}
계획이 $\Base$에 들어가고 회피 반사가 꺼져 있으면($\sigma=0$, \S\ref{sec:family}) 그때의 목표점 $a_0$을 고정하고 고정 목표점 LQR로 호버한다.
인증이 다루는 제어기가 바로 이 고정 목표점 LQR이기 때문이다. 매 스텝 재설정하면 목표점이 밀리미터 단위로 흔들려, 경계에서 들어온 상태가 다시 $\Base$ 밖으로 나가는 일이 생긴다.
고정된 호버는 $V(p-a_0,r)\le c_B$인 동안 종단 조건으로 인정된다.

%======================================================================
\section{Phase I: 안전 도착 회복 정책}\label{sec:phase1}

\paragraph{정책} LQR 입력에 잔차를 더한다.
\begin{equation}
\pi_\theta(x,\zeta)=\sat\big(u_{\mathrm{LQR}}(x)+\Delta u_\theta(o)\big),\qquad u_{\mathrm{LQR}}(x)=\ustar-\tilde Kr .
\label{eq:phase1-policy}
\end{equation}
관측 $o\in\R^{98}$은 $z-z_f$, $v$, $q$(8개), 가까운 순서의 장애물 슬롯 8개(최근접점 벡터 2, 중심 오프셋 2, 반폭 2, $\cos2\psi$, $\sin2\psi$, 모서리 반경, 존재 여부: 각 10개), 추격자 구 슬롯 2개(중심 오프셋 3, 반경 $R_{\mathrm{ball}}$, 존재 여부: 각 5개)로 구성된다.
슬롯 순서는 회복 전개를 시작할 때의 순서로 고정하여 전개 중 입력이 불연속으로 바뀌지 않게 한다. $\Safe$와 $\BaseP$의 소속은 $o$의 정확한 함수이다.

\paragraph{목적함수와 Bellman 관계} 에피소드는 표본 (상태, 장면)에서 시작한 회복 전개 하나이며, $\BaseP$ 도달(성공), $\Safe$ 이탈(실패), 112스텝(2.24\,s) 중 먼저 일어나는 사건에서 끝난다.
\begin{equation}
J(\theta)=\mathbb E\big[\beta^{\tau_B}\,\mathbf 1\{\tau_B<\tau_{\mathrm{out}}\}\big],\qquad \beta=0.99 .
\label{eq:sa}
\end{equation}
장면을 포함한 상태 $y=(x,\zeta)$와 결합 전이 $y^+=G(y,u)$에 대해 최적값은
\begin{equation}
W^\star(y)=1\ (y\in\BaseP\cap\Safe),\qquad W^\star(y)=0\ (y\notin\Safe),\qquad W^\star(y)=\beta\max_{u\in\mathcal U}W^\star(G(y,u))\ \ (\text{그 밖})
\end{equation}
를 만족하고, $W^\star(y)=\beta^{d(y)}$($d$는 안전하게 도착하는 최소 스텝 수)이다. TD3의 critic 목표는 이 관계를 그대로 쓴다.
\begin{equation}
y_{\mathrm{TD}}=b(o)+\beta\,c(o)\,\min_{i=1,2}Q'_i\big(o',\mu'(o')\big),\qquad b=\mathbf 1\{\BaseP\},\ c=\mathbf 1\{\Safe\setminus\BaseP\},
\end{equation}
여기서 $Q'(o',\cdot)$는 $o'\in\BaseP$이면 1, $o'\notin\Safe$이면 0으로 둔다. 엔트로피 항은 critic 목표에 넣지 않는다(도착 순서 ``빨리 $>$ 늦게 $>$ 실패''를 보존하기 위함).
actor 손실은 $-\mathbb E[\min_iQ_i(o,\Delta u_\theta(o))]+0.2\,\mathbb E\big[\tfrac14\|\Delta u/u_{\mathrm{scale}}\|^2\big]$이고, 둘째 항이 잔차를 LQR 쪽으로 붙잡는다. actor는 LQR 모방($\Delta u\approx0$)으로 초기화한다.

\paragraph{학습 분포} 일반(속력 $\le10$\,m/s, 기울기 $\le45^\circ$, 장애물 0--6개, 추격자 구 0--2개), 정지 경로 막힘(속력 $\ge3$\,m/s, LQR 정지점 부근에 장애물), 추격자 근접(속력 $\ge3$\,m/s, LQR 정지점 부근에 추격자 구), 호버 근방(느리고 정지점이 비어 있음) 영역을 가중치 $(0.4,0.2,0.2,0.2)$로 섞는다.
평가에는 저속 상태 근처에 추격자 구가 있는 저속 탈출 영역을 추가로 쓴다(가중치 0). 영역마다 5만 개 시작점을 미리 뽑고, 검증 1000개와 시험 2000개는 따로 뽑는다.

\paragraph{평가 지표와 결과} 시험 상태가 $C_N$에 속하는 비율 $\hat m$을 쓴다. $C_N$은 $N=T/\Delta t=75$스텝(1.5\,s) 안에 $\Safe$를 벗어나지 않고 $\BaseP$에 처음 도달하는 상태의 집합이다.
학습 회복 정책은 가중 평균 $70.1\pm0.3\,\%$(세 시드)로 호버 LQR의 $43.8\,\%$보다 높다(표~\ref{tab:table4_phase1}, 그림~\ref{fig:phase1}).
이 평가는 추격자를 정지한 구로 근사한 Phase I 조건이며, 최종 필터의 이동 추격자 성능과는 구분해야 한다.

%======================================================================
\section{회복 정책군}\label{sec:family}

\paragraph{기본 회복 입력} 기준 집합 안이거나 호버가 고정되어 있으면 고정 목표점 LQR, 밖이면 학습 회복 정책(또는 호버 LQR 변형에서는 재설정 LQR)을 쓴다.
\begin{equation}
\ubase^{(j)}(x)=\begin{cases}\ustar-K(p-a_0,r) & x\in\Base\ \text{또는 호버 고정},\\ \pi_\theta(x,\zeta_{\mathrm{view}}) & \text{그 밖(학습 변형)},\\ u_{\mathrm{LQR}}(x) & \text{그 밖(LQR 변형)}.\end{cases}
\end{equation}
호버가 고정되기 전에는 $a_0=a^\star(x)$이므로 첫 줄은 재설정 LQR과 같다. 학습 회복 정책의 장면 입력 $\zeta_{\mathrm{view}}$는 계획을 만든 시점의 장면(가까운 장애물 8개와 추격자 구)으로 고정한다. Phase I에서 학습한 범위(장애물 10\,m, 추격자 12\,m 이내)만 넣는다.

\paragraph{회피 반사} 최소 추격자 여유도 $h_{\min}=\min_jh_{c,j}$에 따라 회피 입력을 섞는다.
\begin{equation}
\sigma(x)=\clip\!\Big(\frac{h_{\mathrm{off}}-h_{\min}(x)}{h_{\mathrm{off}}-h_{\mathrm{on}}},0,1\Big),\qquad
\pi_b^{(j)}(x)=\sat\big((1-\sigma)\,\ubase^{(j)}(x)+\sigma\,\uesc^{(j)}(x)\big),\qquad h_{\mathrm{on}}=0.5,\ h_{\mathrm{off}}=2.0\ \mathrm m.
\label{eq:reflex}
\end{equation}
$h_{\min}\ge2$\,m이면 기본 회복 입력만, $h_{\min}=1.25$\,m이면 절반씩, $h_{\min}\le0.5$\,m이면 회피 입력만 쓴다. 반사가 작동하면($\sigma>0$) 호버 고정을 푼다.

\paragraph{회피 입력} 가장 위험한 추격자 $j^\star=\arg\min_jh_{c,j}$의 반대 방향 $\hat e_{xy}$를 변형 각 $\vartheta_j$만큼 돌리고, 정지점 $a^\star$ 주변의 장애물과 벽의 바깥 법선으로 민다.
\begin{align}
d&=\nrm\Big(R(\vartheta_j)\,\hat e_{xy}+2\varpi_o\,n_o+2\textstyle\sum_{i=1}^4\varpi_{w,i}\,n_{w,i}\Big),\qquad
\varpi=\clip\!\Big(\frac{d_{\mathrm{push}}-(\mathrm{dist}(a^\star_{xy})-\rho_{xy}-r_d)}{d_{\mathrm{push}}},0,1\Big),\\
v_{\mathrm{des}}&=\big(v_{\mathrm{esc},j}\,d,\ 1.5\,(z_{\mathrm{ref}}-z)\big),\qquad z_{\mathrm{ref}}=\clip(z,\ z_f+1,\ z_{\max}-3),
\end{align}
여기서 $n_o=\nabla\sdf_O(a^\star_{xy})/\|\nabla\sdf_O\|$, $n_{w,i}$는 벽의 안쪽 법선, $\mathrm{dist}$는 각 장애물·벽까지의 거리, $d_{\mathrm{push}}=1.5$\,m이다.
속도 루프는 $a=k_v(v_{\mathrm{des}}-v)$(수평 성분은 $\|a_{xy}\|\le g\tan45^\circ$로 제한), $a_{\mathrm{des}}=(a_{xy},a_z+g)$에서
$a_T=a_{\mathrm{des}}^\top z_b$, $\omega=k_{\mathrm{att}}R^\top(z_b\times\hat a_{\mathrm{des}})-k_\psi\psi\,e_3$로 입력을 만들고 포화시킨다($k_v,k_{\mathrm{att}},k_\psi)=(2,10,3)$, $z_b=R(q)e_3$, $\psi$는 요각).
추격자 여유가 충분할 때는 학습 회복 정책이 장애물을 처리하고, $h_c$가 줄어들수록 회복 정책은 ``$v_{\mathrm{esc}}$로 도피''로 바뀐다. 이는 열린 공간에서 $h_c$를 줄이지 않는 행동이다(\S\ref{sec:hc}).

\paragraph{일곱 변형} 회피 속력비 $v_{\mathrm{esc}}/v_c$와 회피 각 $\vartheta$의 조합으로 회복 정책을 늘린다(표~\ref{tab:variants}).
필터는 어느 하나라도 인증하는 후보를 받아들이므로, 회복 가능 집합은 각 변형 회복 가능 집합의 합집합이 된다.

\begin{table}[h]\centering\small
\caption{필터의 회복 정책 일곱 변형. 번호는 동률일 때의 우선순위이다.}\label{tab:variants}
\begin{tabular}{cccc}\toprule
번호 $j$ & 기본 회복 입력(기준 집합 밖) & $v_{\mathrm{esc}}/v_c$ & $\vartheta$ \\\midrule
0 & 학습 회복 정책 $\pi_\theta$ & 2 & $0^\circ$ \\
1 & 재설정 호버 LQR & 2 & $0^\circ$ \\
2, 3 & 학습 회복 정책 & 2 & $+60^\circ$, $-60^\circ$ \\
4 & 학습 회복 정책 & 3.5 & $0^\circ$ \\
5, 6 & 학습 회복 정책 & 3.5 & $+60^\circ$, $-60^\circ$ \\\bottomrule
\end{tabular}
\end{table}

%======================================================================
\section{이산 다중 회복 안전 필터}\label{sec:filter}



\subsection{후보 입력}
임무 정책의 명목 입력 $\unom$을 각 회복 정책의 현재 입력과 섞어 후보를 만든다.
\begin{equation}
u^{(j)}_\lambda=\sat\big(\lambda\,\unom+(1-\lambda)\,\pi_b^{(j)}(x)\big),\qquad \lambda\in\Lambda=\{1,0.75,0.5,0.25,0\},\quad j=0,\dots,6 .
\label{eq:cand}
\end{equation}
$7\times5=35$개 후보를 검사한다($\lambda=1$에서는 일곱 개가 같은 입력이다).
이는 입력 집합 전체를 최적화하는 것이 아니라 명목 입력과 각 회복 입력을 잇는 선분 위의 유한한 점을 검사하는 것이다.

\subsection{회복 궤적의 전개와 추격자 예측}
후보를 한 스텝 실행한 다음 상태에서 시작해 그 변형의 회복 정책을 $N=150$스텝(3\,s) 전개한다. 추격자는 예측 기체를 향해 $v_c$로 직진한다고 가정한다(거리 감소에 가장 불리한 운동).
\begin{align}
x_1&=f_d\big(x,u^{(j)}_\lambda\big),\qquad x_{i+1}=f_d\big(x_i,\ \pi_b^{(j)}(x_i;c_i)\big),\qquad i=1,\dots,N, \label{eq:rollout}\\
c_{\ell,i+1}&=c_{\ell,i}+\min\!\big(v_c\Delta t,\ \|p_i-c_{\ell,i}\|\big)\,\frac{p_i-c_{\ell,i}}{\|p_i-c_{\ell,i}\|},\qquad c_{\ell,0}=c_\ell(\text{현재}),\qquad c_{\ell,i,z}\ge z_f+r_c . \label{eq:pursue}
\end{align}
$\ell$은 추격자 번호이다. 회피 반사는 예측 추격자 위치 $c_i$를 보고, 학습 회복 정책의 장면 입력은 계획 시점의 것으로 고정된다. 종단 첨자는 $H=N+1$이다.

\subsection{궤적 조건(신규 계획 여유)}
새 계획은 모든 예측 노드 $i=1,\dots,H$에서 여유를 더한 안전 집합에 있어야 한다.
\begin{equation}
\min\{h_f,h_o,h_w,h_z\}(x_i)\ge m_o=0.15\ \mathrm m,\qquad \min_\ell h_{c,\ell}(x_i;c_{\ell,i})\ge m_c=0.2\ \mathrm m .
\label{eq:fresh}
\end{equation}
여유는 반사 방향의 재조준과 실제 추격자가 예측과 다르게 움직이는 것에 대비한 것이다. 계획 유지(\S\ref{sec:commit})의 재검사는 여유 없이($m_o=m_c=0$) 안전 집합만 요구한다.

\subsection{종단 집합}\label{sec:terminal}
종단 상태 $x_H$는 호버하거나 직선 제동으로 멈출 수 있어야 한다.
\begin{align}
\Term_{\mathrm{hover}}&=\big\{x_H:\ \text{호버 고정이면 } V(p_H-a_0,r_H)\le c_B,\ \text{아니면 } x_H\in\Base\big\},\\
\Term_{\mathrm{flee}}&=\big\{x_H:\ \sdf_O(p_{\mathrm{brake}})\ge r_d,\ d_{\mathrm{wall}}(p_{\mathrm{brake}})\ge r_d\big\},\qquad
p_{\mathrm{brake}}=p_{H,xy}+\frac{v_{H,xy}\|v_{H,xy}\|}{2a_e},\\
\Term&=\big(\Term_{\mathrm{hover}}\cup\Term_{\mathrm{flee}}\big)\cap\big\{\min_\ell h_{c,\ell}(x_H;c_{\ell,H})\ge m_c\big\}.
\label{eq:terminal}
\end{align}
$p_{\mathrm{brake}}$는 수평 속도 방향으로 $a_e$로 제동할 때의 정지점이다. 따라서 종단이 항상 $\Base$인 것은 아니며, 도피 중인 상태도 제동 공간이 남아 있으면 허용한다.
회복 정책 $j$의 유한시간 회복 가능 집합은
\begin{equation}
C^{(j)}=\big\{(x,\zeta):\ x_i\in\Safe^{+}\ (i=1,\dots,H),\ x_H\in\Term\big\},\qquad C=\bigcup_{j=0}^{6}C^{(j)}
\end{equation}
이며($\Safe^+$는 식~\eqref{eq:fresh}의 여유를 더한 안전 집합), 후보 실행 직후 상태가 $C^{(j)}$에 속하는지를 검사한다.
합집합은 각 집합을 포함하므로 같은 검사 조건에서 한 회복 정책만 쓰는 경우보다 선택지가 줄지 않는다. 다만 모든 회피 가능 상태를 찾는다는 뜻은 아니다.

\subsection{선택 규칙}
인증된 조합의 집합을 $\mathcal A$라 할 때 명목 입력의 비중이 가장 큰 후보를 실행한다.
\begin{equation}
(\lambda^\star,j^\star)\in\arg\max_{(\lambda,j)\in\mathcal A}\lambda,\qquad u=u^{(j^\star)}_{\lambda^\star}\quad(\text{동률이면 작은 }j).
\label{eq:select}
\end{equation}
변형마다 회복 입력이 다르므로 이는 $\|u-\unom\|$의 최소화와 일반적으로 같지 않다. ``명목 입력이 회복 가능성을 유지하면 그대로 통과시키고, 필요할 때만 회복 쪽으로 움직인다''는 최소 수정의 이산 형태이다.
그림~\ref{fig:snapshot}은 실제 에피소드에서 필터가 한 스텝에 검사한 내용을 다시 계산한 것이다. 명목 입력($\lambda=1$)은 앞의 장애물 때문에 모든 회복 정책에서 거부되고, 회복 정책 0과 $\lambda=0.75$의 조합이 인증되어 실행되었다(평가 기록과 일치).
\IfFileExists{\snapdir/snapshot_ep845_step091.png}{%
\begin{figure}[H]\centering
\includegraphics[width=\linewidth]{\snapdir/snapshot_ep845_step091.png}
\caption{필터의 한 스텝 결정(시드 0, 에피소드 845, $t=1.8$\,s). 왼쪽: 일곱 회복 정책의 3\,s 계획($\lambda=0$)과 끝나는 방식(점선 원: 호버 타원체, $\times$: 제동 정지점), 선택한 후보의 계획(굵은 파랑)과 추격자 예측(초록 점선), 추격자의 $h_c=0$ 거리(옅은 원). 오른쪽: 후보 35개의 인증 결과(초록: 인증, 빨강: 궤적 조건 위반, 주황: 종단 조건 위반)와 선택한 칸.}\label{fig:snapshot}
\end{figure}}{}

\subsection{계획 유지와 추격자 예측 재생}\label{sec:commit}
학습 회복 정책의 입력(가까운 장애물 순서, 범위 안의 장애물)은 스텝마다 바뀔 수 있어 직전에 인증한 계획이 다음 스텝에 달라질 수 있다.
그래서 필터는 기억 $m=(\zeta_{\mathrm{view}},a_0,\text{호버 고정 여부},j_c,\text{추격자 슬롯 순서},\hat c)$를 저장한다. $\hat c$는 인증한 계획이 예측한 다음 스텝의 추격자 위치이다.
$\mathcal A=\varnothing$이면 저장한 변형 $j_c$의 회복 입력을 \emph{저장한 장면 입력과 예측 추격자 위치} $\hat c$로 계산해 실행한다.
모델이 정확하면 이 입력은 인증된 회복 궤적의 다음 입력과 정확히 같다(committed continuation). 이 입력도 실제 추격자 위치에 대해 여유 없이 다시 검사하고, 실패하면 그대로 실행하되 ``인증 없음''으로 기록한다.
예측 추격자 위치는 식~\eqref{eq:pursue}로 한 스텝씩 진행한다.

\subsection{알고리즘}
\begin{algorithm}[H]
\caption{이산 다중 회복 안전 필터 $F$ (한 제어 스텝)}\label{alg:filter}
\KwIn{상태 $x$, 장면(장애물, 추격자 $c$), 명목 입력 $\unom$}
\KwMem{$m=(\zeta_{\mathrm{view}},a_0,\text{고정 여부},j_c,\text{슬롯 순서},\hat c)$}
\KwOut{실행 입력 $u$, 갱신된 기억 $m'$, 정보 $(\lambda^\star,\text{인증 여부})$}
현재 장면에서 새 장면 입력 $\zeta_{\mathrm{new}}$를 만들고, 추격자 예측 $c_1\gets$ 식~\eqref{eq:pursue}를 한 스텝 적용\;
\For{$j=0,\dots,6$}{
  $\pi_b^{(j)}(x)$ 계산(식~\eqref{eq:reflex})\;
  \For{$\lambda\in\{1,0.75,0.5,0.25,0\}$}{
    $u^{(j)}_\lambda\gets\sat(\lambda\unom+(1-\lambda)\pi_b^{(j)}(x))$, \ $x_1\gets f_d(x,u^{(j)}_\lambda)$\;
    회복 정책 $j$를 $N=150$스텝 전개(식~\eqref{eq:rollout}--\eqref{eq:pursue})\;
    모든 노드가 식~\eqref{eq:fresh}를 만족하고 $x_H\in\Term$이면 $(\lambda,j)$를 $\mathcal A$에 추가\;
  }
}
\eIf{$\mathcal A\neq\varnothing$}{
  $(\lambda^\star,j^\star)\gets\arg\max_{\mathcal A}\lambda$ (동률이면 작은 $j$), \ $u\gets u^{(j^\star)}_{\lambda^\star}$\;
  $m'\gets(\zeta_{\mathrm{new}},a_0',\text{고정 여부}',j^\star,\text{새 슬롯 순서},c_1)$\;
}{
  $u\gets\pi_b^{(j_c)}(x;\zeta_{\mathrm{view}},\hat c)$ \tcp*{직전 인증 계획의 다음 입력}
  여유 없이 실제 추격자에 대해 재검사하고, 실패하면 ``인증 없음''으로 기록\;
  $m'\gets(\zeta_{\mathrm{view}},a_0',\text{고정 여부}',j_c,\text{슬롯 순서},\Pi(\hat c))$\;
}
\KwRet{$u,\ m'$}
\end{algorithm}

\subsection{보장되는 것과 보장되지 않는 것}\label{sec:guarantee}
\begin{keybox}{명제 (가정 아래의 진술)}
가정: (A1) 기체 모델이 정확하고 외란이 없다. (A2) 각 추격자의 속력은 $v_c$ 이하이다. (A3) 3\,s 예측 구간에서 속력 $v_c$의 순수 추적이 $h_c$에 대해 최악의 경우이다(거리에 대해서는 성립, 시선 방향에 대해서는 근사).

이때 실행되는 모든 입력은 인증된 후보이거나 직전 인증 계획의 이어짐이다. 인증된 계획은 3\,s 동안 상태를 여유를 둔 $\Safe$ 안에 두고 $\Term$에서 끝난다.
추격자 예측 재생 덕분에 이어짐은 인증된 회복 궤적의 꼬리를 정확히 재현하며, 실제 추격자 위치에 대해 다시 검사된다.
\end{keybox}
남는 간극은 $\Term$이 추격자에 대해 불변 집합이 아니라는 점이다. 호버는 결국 추격자에게 따라잡히고, 도피는 열린 공간이 남아 있는 동안만 안전하다.
또한 이어짐을 재검사할 때 예측 구간이 한 스텝 늘어나므로, 그 마지막 상태가 $\Term$을 만족하지 못해 인증을 잃을 수 있다.
따라서 이 필터는 3\,s 앞을 인증하고 20\,ms마다 다시 계획하는 장치이며, 이동 추격자에 대한 무한시간 불변성의 증명은 아니다.
실제 측정에서는 인증 상실이 PS2 정책 에피소드의 1.5\,\%에서 나타났다(표~\ref{tab:table2_final}).

\subsection{계산량}
한 제어 스텝에서 후보 35개와 이어짐 7개, 즉 42개의 150스텝 전개(제어 스텝당 약 6300회의 회복 입력 계산)를 수행하며, 학습 변형은 그때마다 98차원 입력의 신경망을 평가한다.
GPU(RTX 4050 노트북)에서 1000개 에피소드를 묶어 평가하면 15\,s 에피소드 전체가 약 350\,s 걸린다(묶음 기준 약 0.46\,ms/기체-스텝).
단일 기체의 실시간 최악 실행 시간은 측정하지 않았다.

%======================================================================
\section{비교 필터: bCBF-QP}\label{sec:bcbf}

PS2-RL의 backup CBF QP를 같은 집합에 적용한 비교 구성이다. 기체와 추격자를 포함한 상태를 $y=(p,v,q,c_1,\dots,c_C)$, 결합 동역학을 $\dot y=f(y)+g(y)u$,
고정된 회복 정책(반사 포함)의 흐름을 $\varphi(\tau,y)$, 그 벡터장을 $f_\pi$라 하자.
민감도 $Q_i=\partial\varphi(\tau_i,y)/\partial y$는 Euler 전개의 정확한 Jacobian으로 전파한다.
절대 예측 시각 $t_i$를 고정하면 $\tau_i=t_i-t$가 시간에 따라 줄어들므로, 매끄러운 구간에서
\begin{equation}
\frac{d}{dt}h\big(\varphi(t_i-t,y(t))\big)=\nabla h(\varphi_i)\big[Q_i\big(f(y)+g(y)u\big)-f_\pi(\varphi_i)\big]
\end{equation}
이다. 둘째 항이 남은 예측 시간이 줄어드는 데서 나온다. 이를 $-\alpha h(\varphi_i)$ 이상으로 제한한 선형 제약을 노드 $\tau_i$(1.5\,s 구간, 3스텝 간격)와 모든 안전 함수 $h\in\{h_o,h_f,h_w,h_z,h_c\}$에 대해 둔다($\alpha=4$).
종단 노드에서는 기준 집합 함수(반사가 작동 중이면 제동점 함수)와 $h_c$에 같은 형태의 제약을 $\alpha_B=2$로 둔다.
\begin{equation}
\min_{u,\ s\ge0}\ \tfrac12\sum_kw_k(u_k-u_{\mathrm{nom},k})^2+\tfrac12w_ss^2\quad\text{s.t. 완화된 CBF 제약, }u\in\mathcal U,\qquad w=(1,0.1,0.1,1),\ w_s=10^6 .
\end{equation}
roll/pitch 각속도의 가중치가 작아 추력보다 기울기로 경로를 바꾸는 해를 선호한다.
회복 입력 $u=\pi_b(x)$가 제약을 $10^{-2}$보다 크게 위반하면(정규화된 행 기준) QP를 건너뛰고, QP 출력은 다시 이산 필터(최종 회복 정책군)를 통과시킨다.
따라서 이 구성의 접촉 억제는 QP 단독의 결과가 아니며, 선형화·이산화·slack 때문에 표준 연속시간 CBF 불변성 정리를 그대로 적용할 수 없다.

%======================================================================
\section{임무 정책}\label{sec:striker}

\paragraph{기준 접근 제어} 표적 위치 $p_t$와 속도 $v_t$에 대해
\begin{equation}
v_{\mathrm{des}}=v_{\mathrm{cr}}\frac{p_t-p}{\|p_t-p\|}+v_t,\qquad v_{\mathrm{cr}}=6\ \mathrm{m/s}
\end{equation}
를 회피 반사와 같은 속도·자세 루프($k_v=2.5$, 기울기 $50^\circ$, $k_{\mathrm{att}}=10$, $k_\psi=3$)로 추종해 $\upur$를 만든다. 장애물과 추격자를 고려하지 않는다.

\paragraph{잔차 정책}
\begin{equation}
\unom=\sat\big(\upur(x)+\Delta u_\xi(o)\big),\qquad |\Delta a_T|\le0.1\,a_{\max}=3.92\ \mathrm{m/s^2},\quad |\Delta\omega_i|\le0.25\,\omega_{\max}=4.5\ \mathrm{rad/s}.
\label{eq:residual}
\end{equation}
관측은 135차원이다: 회복 정책 형식의 장면 관측 98개, 표적 오프셋/10, 표적 속도, 추격자 속도, 경기장 내 위치(14개),
유도 특징 19개(시선 방향 3, $d/10$, $\log(d+0.1)$, 상대 속도/5 3, 접근 속도/5, 시선각 변화율/2 3, 요 좌표계의 상대 위치/10 3과 상대 속도/5 3, 표적 높이/2),
그리고 보정 대상인 기준 입력 $(\upur-\ustar)/u_{\mathrm{scale}}$(4개).

\paragraph{보상} 표적 거리 $d_k=\|p_k-p_{t,k}\|$에 대해
\begin{equation}
r_k=w_p(d_{k-1}-d_k)-w_t-\frac{w_a}{4}\Big\|\frac{\Delta u_k}{u_{\mathrm{scale}}}\Big\|^2+100\cdot\mathbf 1_{\text{성공}}-150\cdot\mathbf 1_{\text{접촉 또는 이탈}},\qquad w_p=2,\ w_t=w_a=0.05 .
\label{eq:reward}
\end{equation}
거리 항은 $w_p(d_0-d_{\mathrm{end}})$로 합쳐지고, 시간 비용은 빠른 성공을 장려한다. 시간초과는 critic 목표에서 종료로 보지 않는다(부트스트랩 유지).

\paragraph{학습} SAC, $\gamma=0.995$, 보상 척도 0.1, 병렬 환경 64개, 배치 256, 재생 버퍼 $10^6$, 학습률 $3\times10^{-4}$, $\tau=0.005$, 온도 초기값 0.1(하한 $10^{-3}$, 목표 엔트로피 $-4$), 은닉층 256$\times$2(ELU).
처음 20만 스텝은 actor를 고정하고 critic만 학습한다(critic warm-up). actor 손실에 고정점 벌점 $0.2\,\mathbb E\big[\tfrac14\|\Delta u/b\|^2\big]$($b$는 잔차 한계)을 더한다.
1000만 스텝 학습하며, 25만 스텝마다 128개 에피소드로 평가해 반환값이 가장 큰 체크포인트를 쓴다. 처음부터 학습한 정책은 실패했고(성공 0, 이후 0.41--0.49), 기준 접근 제어의 잔차로 두어야 안정적으로 학습되었다.

%======================================================================
\section{PS2: 필터를 포함한 임무 학습}\label{sec:ps2}

필터를 환경의 일부로 두고 임무 정책을 다시 학습한다.
\begin{align}
u_{\mathrm{nom},k}&=\sat\big(\upur(x_k)+\Delta u_\xi(o_k)\big),\qquad (u_k,m_{k+1})=F\big(x_k,\zeta_k,m_k,u_{\mathrm{nom},k}\big),\qquad x_{k+1}=f_d(x_k,u_k), \label{eq:ps2}\\
r^{\mathrm{PS2}}_k&=r_k-w_f\|\eta_k\|^2,\qquad \eta_k=\frac{u_k-u_{\mathrm{nom},k}}{u_{\mathrm{scale}}},\qquad w_f=0.5 . \label{eq:ps2-reward}
\end{align}
정책은 자신이 제안한 입력이 그대로 실행된다고 가정하지 않고, 필터 개입 이후의 상태와 보상으로 학습한다. 필터가 자주 막는 행동보다 실제로 진전이 생기고 수정 비용이 작은 행동이 유리해진다.
필터 기억 $m$까지 포함한 전이 $(x,\zeta,m)\mapsto(x^+,\zeta^+,m^+)$는 정책에서 보면 환경의 일부이므로, 이산 필터를 미분하지 않고도 off-policy SAC로 학습할 수 있다.
다만 $m$은 관측되지 않으므로 관측이 완전한 상태라고 주장하지 않는다. 필터 출력을 복사하는 모방학습이 아니라, 수정된 전이에서 같은 강화학습 목표를 다시 최적화하는 것이다.

\paragraph{설정} 같은 시드의 잔차 정책 최고 체크포인트에서 시작하여, 처음 10만 스텝은 actor를 고정하고 critic을 재적응시키며, 온도는 시작 값으로 고정한다(학습률 0).
배치할 필터(최종 회복 정책군, 추격자 예측 재생)와 같은 필터로 300만 스텝 학습하고(학습 = 배치), 25만 스텝마다 64개 에피소드로 평가해 반환값이 가장 큰 체크포인트를 쓴다.
PS2 정책은 필터와 함께 쓰도록 학습한 정책이며, 필터 없이 실행하면 성공률이 잔차 정책보다 낮다(62.1\,\% 대 74.0\,\%). 필터가 장애물에서 비켜 주는 것에 의존하기 때문이다.

\paragraph{학습 시간} RTX 4050 노트북 GPU에서 시드당 Phase I 약 10분(1000만 스텝), 잔차 정책 28--34분(1000만 스텝), PS2 72--76분(300만 스텝)이 걸렸다.

%======================================================================
\section{평가 방법}\label{sec:protocol}

\paragraph{절차} 시드 $s\in\{0,1,2\}$마다 전체 파이프라인(Phase I $\to$ 잔차 정책 $\to$ PS2)을 그 시드로 학습한다.
모든 구성을 같은 1000개 평가 장면(장면 시드 9001, 학습 장면 풀과 별도)으로 평가하며, 장면마다 출발 교란과 표적·추격자 초기 조건까지 같다.
에피소드는 15\,s이고, 10\,s 결과는 같은 에피소드를 10\,s에서 절단한 것이다(10\,s까지 끝나지 않은 에피소드는 시간초과로 센다).
각 시드는 그 시드의 임무 정책을 학습한 추격자 설정(원래 추격자)으로 평가한다.

\paragraph{지표} 성공률, 접촉률(장애물·바닥·추격자), 시간초과율, 영역 이탈률, 성공 에피소드의 평균 도달 시간,
평균 명목 유지율 $\bar\lambda$(에피소드 스텝 평균 $\lambda^\star$, 새 인증이 없는 스텝은 0),
인증 상실률(한 스텝 이상 유효한 인증 계획이 없었던 에피소드의 비율). 인증 상실은 접촉과 다른 지표이다.

\paragraph{통계} 시드별 비율의 평균 $\pm$ 표준편차를 보고하고, 시드를 합친 횟수 $k/n$에 대해 양측 95\,\% Clopper--Pearson 구간을 쓴다.
\begin{equation}
p_L=\mathrm{B}^{-1}\!\big(\tfrac{\alpha}{2};\,k,\,n-k+1\big),\qquad p_U=\mathrm{B}^{-1}\!\big(1-\tfrac{\alpha}{2};\,k+1,\,n-k\big),\qquad \alpha=0.05,
\end{equation}
여기서 $\mathrm B^{-1}$는 베타 분포의 분위수이다. $k=0$이면 $p_U=1-(\alpha/2)^{1/n}$이고, 단측 95\,\% 상한은 $1-0.05^{1/n}\approx3/n$이다.
$n=1000$에서 0.30\,\%(단측)와 0.37\,\%(양측), $n=3000$에서 0.10\,\%와 0.12\,\%이다.
모든 구성이 같은 장면을 쓰므로 두 구성의 비교는 \emph{쌍 비교}이다. 구성 A만 성공한 장면 수 $b$와 B만 성공한 장면 수 $c$로 정확 McNemar 검정을 한다.
\begin{equation}
p_{\mathrm{McNemar}}=\min\Big\{1,\ 2\sum_{i=0}^{\min(b,c)}\binom{b+c}{i}2^{-(b+c)}\Big\}.
\end{equation}
구성끼리 비교하는 표(표~\ref{tab:table2_final}, \ref{tab:table3_ablation})는 모든 행이 평가된 공통 시드만 쓴다.

%======================================================================
\section{결과}\label{sec:results}

\begin{warnbox}{수치의 기준}
아래 표와 그림은 업로드된 평가 결과(시드 0 전체, 시드 1 일부, 시드 2 미평가)에서 만든 것이다. 각 행의 ``시드 $\times$ 장면'' 열이 사용한 에피소드 수를 나타낸다.
세 시드 평가를 마친 뒤 \file{STEPS=tables bash docker\_batch/run\_strike\_paper.sh}로 다시 만들고 이 문서를 다시 컴파일하면 표와 그림이 갱신된다.
\end{warnbox}

\subsection{주요 비교}
\tabin{table1_10s}
\tabin{table1_15s}

\begin{figure}[H]\centering
\figin{\figpath{fig_outcomes.pdf}}
\caption{구성별 결과 비율(왼쪽 10\,s, 오른쪽 15\,s 제한). 막대 안의 수는 성공률(\%)이다.}\label{fig:outcomes}
\end{figure}

\paragraph{필터의 효과} 같은 임무 정책에서 최종 필터는 접촉을 거의 없앤다. 잔차 정책은 26.1\,\%(521/2000)에서 0/2000으로, 기준 접근 제어는 57.5\,\%에서 1/2000으로 줄었다.
대신 우회와 정지 때문에 10\,s 시간초과가 늘어난다(잔차 정책 0\,\%$\to$15.5\,\%). 15\,s 제한에서는 시간초과 대부분이 늦은 성공으로 바뀐다.

\paragraph{PS2의 효과} 같은 최종 필터를 쓸 때 PS2 정책은 잔차 정책보다 성공률이 높다(10\,s 84.5$\to$90.4\,\%, 15\,s 94.4$\to$97.6\,\%).
같은 장면 쌍 비교에서 PS2만 성공한 장면이 172개, 잔차 정책만 성공한 장면이 54개로(10\,s) 차이는 유의하다($p<10^{-6}$, 표~\ref{tab:table5_paired}).
평균 도달 시간도 6.53\,s에서 6.11\,s로 줄었다. 필터를 고정하고 임무 정책만 바꾼 비교이므로 필터 포함 재학습의 효과를 보여 주는 핵심 근거이다.

\begin{figure}[H]\centering
\begin{subfigure}{0.49\linewidth}\centering\figin{\figpath{fig_strike_time.pdf}}\caption{누적 성공률}\label{fig:strike-time}\end{subfigure}\hfill
\begin{subfigure}{0.49\linewidth}\centering\figin{\figpath{fig_safety.pdf}}\caption{추격자 최소 여유 거리의 누적 분포}\label{fig:safety}\end{subfigure}
\caption{(a) 시간에 따른 누적 성공률(음영은 시드를 합친 95\,\% 구간). 점선은 10\,s. (b) 에피소드별 추격자 최소 여유 거리 $\min_t(\|p-c\|-r_c-r_d)$의 누적 분포. 필터를 쓴 구성에서 0.2\,m의 계단은 추격자가 바짝 따라오는 동안 필터가 $h_c$를 신규 계획 여유 $m_c=0.2$\,m에 유지한 에피소드들이다(멀어지는 중에는 $h_c$가 거리 여유와 같다).}
\end{figure}

\subsection{최종 필터 선택과 구성 요소}
\tabin{table2_final}
\tabin{table3_ablation}
\tabin{table5_paired}

\begin{figure}[H]\centering
\figin{\figpath{fig_ablation.pdf}}
\caption{구성 요소 제거 실험(PS2 정책, 15\,s). 오차 막대는 95\,\% Clopper--Pearson 구간이다.}\label{fig:ablation}
\end{figure}

\begin{itemize}
\item \textbf{회피 변형}: 단일 회피 필터는 접촉 1.7\,\%(17/1000)로, 최종 필터(1/1000)와의 쌍 비교에서 차이가 유의하다($p=3.1\times10^{-5}$). 성공률도 낮다(15\,s 95.7\,\% 대 98.4\,\%).
\item \textbf{bCBF-QP}: 성공률은 비슷하지만 최종 필터가 10\,s와 15\,s 모두 높고($p=0.0015$, $0.043$), 평균 도달 시간이 짧다(5.93\,s 대 6.33\,s, 공통 시드 기준). bCBF는 명목 유지율이 더 높지만(0.97) 이후 이산 검사에서 인증을 더 자주 잃는다(2.5\,\% 대 1.5\,\%).
\item \textbf{학습 회복 정책}: 회복 정책을 호버 LQR로 바꾸면 15\,s 성공률이 82--84\,\%로 크게 떨어지고 시간초과가 15--17\,\%로 늘어난다(시드 2개와 공통 시드 기준). Phase I 회복 정책이 필터의 허용 범위를 넓히는 것이 성능의 주요 원인이다.
\item \textbf{구형 근사}: 추격자를 정지한 구로 보면 접촉이 26.3\,\%로, 상대 접근 속도를 반영한 $h_c$가 필수적이다.
\item \textbf{계획 유지의 추격자 예측 재생과 호버 latch}: 1000개 장면에서 성공·접촉의 차이는 관측되지 않았다(98.6\,\%, 98.3\,\% 대 98.4\,\%). 두 요소는 성능보다 이어짐이 인증된 계획을 정확히 재현하고 인증 대상 제어기와 일치하게 하는 일관성을 위한 것이다.
\end{itemize}

\begin{figure}[H]\centering
\figin{\figpath{fig_filter.pdf}}
\caption{필터 개입. 왼쪽: 에피소드 평균 명목 유지율 $\bar\lambda$의 누적 분포(1에서의 계단은 필터가 한 번도 수정하지 않은 에피소드). 오른쪽: 인증 상실 에피소드 비율.}\label{fig:filter}
\end{figure}

\subsection{학습 과정과 Phase I}
\begin{figure}[H]\centering
\figin{\figpath{fig_training.pdf}}
\caption{학습 중 평가 성공률(시드별). 원은 평가에 사용한 체크포인트. 필터 없이 학습한 잔차 정책은 시드 1, 2에서 500--650만 스텝 이후 붕괴했으며, 평가에는 최고 반환값 체크포인트(시드 0: 850만, 시드 1: 125만, 시드 2: 175만 스텝)를 썼다. PS2 학습은 필터 안에서 안정적이다.}\label{fig:training}
\end{figure}
\tabin{table4_phase1}
\begin{figure}[H]\centering
\figin[0.55\linewidth]{\figpath{fig_phase1.pdf}}
\caption{Phase I 회복 성공률(1.5\,s 안에 안전하게 기준 집합 도달). 학습 회복 정책은 세 시드 평균과 표준편차.}\label{fig:phase1}
\end{figure}

\subsection{대표 에피소드}\label{sec:results-episodes}
\file{run\_strike\_paper.sh}는 시드 0의 평가에서 필터가 장애물 충돌을 막은 장면, 추격자 접촉을 막은 장면, 추격자 근접 회피, 유일한 접촉 장면,
기준 제어가 실패한 장면 등을 자동으로 고르고(\file{episode\_picks.json}), 각 장면을 여러 구성으로 다시 실행해 다음을 만든다.
\begin{itemize}
\item 2차원 궤적 비교(구성별 패널), 3차원 궤적(속력 색), 3차원 비교 그림;
\item 상세 신호 그림: 표적 거리, 장애물·벽·바닥·천장 여유, $h_c$와 반사 가중치 $\sigma$, 속력과 고도, $\lambda^\star$와 인증 상태(신규/이어짐/인증 없음), 사용한 회복 변형, 입력과 명목 입력;
\item 2차원 GIF와 회전하는 3차원 GIF(추격자 $h_c$ 요구 거리 구 표시);
\item 필터 스냅샷(\file{scripts/strike\_filter\_snapshot.py}): 처음 명목 입력을 줄인 스텝, 이전 계획을 이어 간 스텝, 인증을 잃은 스텝에서 필터의 검사를 다시 계산한 그림(그림~\ref{fig:snapshot}).
\end{itemize}
\IfFileExists{\episodedir/overview_main.png}{%
\begin{figure}[H]\centering
\includegraphics[width=\linewidth]{\episodedir/overview_main.png}
\caption{대표 에피소드(시드 0)의 수평 궤적. 실선은 최종 필터를 쓴 PS2 정책과 잔차 정책, 점선은 필터 없는 PS2 정책이다. 별은 성공, X는 접촉, 회색 원은 시간초과.}\label{fig:episodes-overview}
\end{figure}}{\IfFileExists{\episodedir/overview.png}{%
\begin{figure}[H]\centering
\includegraphics[width=\linewidth]{\episodedir/overview.png}
\caption{대표 에피소드(시드 0)의 수평 궤적.}\label{fig:episodes-overview}
\end{figure}}{}}

\paragraph{유일한 접촉 사례(시드 0, 에피소드 845)} 기체는 큰 원기둥과 위쪽 벽 사이의 좁은 공간($y\approx9$\,m)으로 들어간 뒤 약 6\,s부터 새 인증을 찾지 못해
이전 계획(호버 LQR + 회피 변형)을 이어 갔고, 약 7.5\,s부터는 그 계획도 재검사를 통과하지 못했다(인증 없음).
추격자가 다가오면서 회피 반사가 작동했지만 벽과 장애물이 도피 방향을 막아 9.2\,s에 접촉했다. 같은 장면에서 bCBF 구성도 같은 방식으로 실패했다.
이 사례는 \S\ref{sec:guarantee}의 남는 간극, 즉 호버 종단이 다가오는 추격자에 대해 불변 집합이 아니고 막힌 공간에서는 도피도 보장되지 않는다는 점이 실제로 나타난 경우이다.
인증을 잃은 7.5\,s의 검사를 다시 계산하면(그림~\ref{fig:snapshot-lost}) 35개 후보 모두 궤적 조건을 위반했다. 위쪽 벽($y=10$\,m) 쪽으로 움직이던 기체에서 일곱 회복 정책의 3\,s 계획이 모두 벽을 넘었기 때문이다.
Phase I 회복 정책은 벽 없이 학습되었고, 회피 반사는 추격자 여유가 $h_{\mathrm{off}}=2$\,m 이상이면 섞이지 않는다. 벽을 포함한 Phase I 학습이 남은 개선 방향이다.
\IfFileExists{\snapdir/snapshot_ep845_step374.png}{%
\begin{figure}[H]\centering
\includegraphics[width=\linewidth]{\snapdir/snapshot_ep845_step374.png}
\caption{인증을 잃은 시점의 검사(에피소드 845, $t=7.5$\,s). 일곱 회복 정책의 계획이 모두 위쪽 벽을 넘어 새 인증이 없고, 필터는 이전 계획을 이어 실행했다.}\label{fig:snapshot-lost}
\end{figure}}{}
\IfFileExists{\episodedir/ep845/compare2d.png}{%
\begin{figure}[H]\centering
\includegraphics[width=\linewidth]{\episodedir/ep845/compare2d.png}
\caption{에피소드 845를 여러 구성으로 다시 실행한 수평 궤적(CPU 재실행).}\label{fig:ep845-2d}
\end{figure}}{}
\IfFileExists{\episodedir/ep845/detail_ps2_ours.png}{%
\begin{figure}[H]\centering
\includegraphics[width=0.86\linewidth]{\episodedir/ep845/detail_ps2_ours.png}
\caption{에피소드 845, PS2 정책 + 최종 필터의 상세 신호. 위에서부터 표적 거리, 기하 여유, $h_c$와 반사 가중치 $\sigma$, 속력과 고도, $\lambda^\star$와 인증 상태(노랑: 이전 계획 유지, 빨강: 인증 없음), 사용한 회복 변형, 입력과 명목 입력.}\label{fig:ep845-detail}
\end{figure}}{}

%======================================================================
\section{한계와 적용 범위}\label{sec:limits}
\begin{enumerate}
\item \textbf{종단 집합의 불변성}: 호버 기준 집합은 정적 장애물에 대한 회복 수단이다. 추격자가 계속 다가오면 호버는 영구히 안전하지 않고, 도피 종단도 유한시간 예측의 조건이다. 정적 기준 집합에 기대는 원래 이론의 보장을 그대로 옮길 수 없다.
\item \textbf{추격자 모델}: 예측은 현재 시선 방향의 순수 추적이며, 장애물 우회와 시선각 변화를 포함한 모든 추격 전략의 최악값이라는 증명은 없다. 반응 지연과 이탈 가속도에 기반한 $h_c$도 근사이다.
\item \textbf{모델 오차}: 입력 지연, 공력, 센서 오차, 바람이 있으면 3\,s 검사와 실제 경로가 달라질 수 있다. 모든 오차를 포함하는 강건 경계는 제시하지 않았다.
\item \textbf{시간 이산 검사}: 검사는 20\,ms 노드에서 한다. 노드 사이 연속시간 충돌까지 배제하려면 이동량이나 보간 오차를 이용한 추가 조건이 필요하다.
\item \textbf{인증 상실}: 계획 유지 중에도 인증을 잃을 수 있다(PS2 에피소드의 1.5\,\%). 이 시점은 명제의 전제를 벗어난다.
\item \textbf{통계}: 1000개 장면과 2--3개 시드의 결과이다. 접촉 1/2000은 접촉확률이 0이라는 증명이 아니며 95\,\% 상한은 0.28\,\%이다.
\item \textbf{실기체}: 결과는 시뮬레이션의 성공과 접촉이다. 실시간 최악 실행 시간, PX4 추력·각속도 추종 오차, 실기체 운용 성능은 검증하지 않았다.
\end{enumerate}

%======================================================================
\appendix
\section{하이퍼파라미터}\label{app:hparams}
\begin{center}\small
\begin{tabularx}{\linewidth}{>{\raggedright\arraybackslash}p{3.0cm}>{\raggedright\arraybackslash}X}\toprule
구분 & 값 \\\midrule
모델 & $\Delta t=0.02$\,s, $a_T\in[0,4g]$, $|\omega_i|\le18$\,rad/s, $r_d=0.3$\,m, $z_f=1$\,m, $z_{\max}=10$\,m \\
경기장·장면 & 40$\times$20\,m, 장애물 3개, 추격자 1개($r_c=0.4$\,m, $v_c=2$\,m/s, $a_c=2$\,m/s$^2$, 반발 거리 2\,m), 표적 0.2--0.8\,m/s, 성공 반경 0.6\,m, 15\,s \\
$h_c$ & $\delta=0.2$\,s, $a_e=5$\,m/s$^2$ \\
호버 LQR & $Q=\mathrm{diag}(2,2,8,0.01,0.01,0.5,8,8,2)$, $R_u=\mathrm{diag}(0.06,0.15,0.15,0.12)$, $c_B=50$, $\rho_{xy}=1.22$\,m, $z_{\mathrm{ext}}=0.63$\,m \\
Phase I & TD3, 1000만 스텝, 병렬 64, 배치 256, 버퍼 $10^6$, $\beta=0.99$, 학습률 actor $10^{-4}$/critic $3\times10^{-4}$, $\tau=0.0025$, 정책 지연 2, 탐색 잡음 0.1, 잔차 벌점 0.2, 에피소드 112스텝, 구형 근사 $T=1.5$\,s, 영역 가중치 $(0.4,0.2,0.2,0.2)$ \\
회피 반사 & $h_{\mathrm{on}}=0.5$\,m, $h_{\mathrm{off}}=2.0$\,m, $(k_v,k_{\mathrm{att}},k_\psi)=(2,10,3)$, 기울기 $45^\circ$, $d_{\mathrm{push}}=1.5$\,m, $z_{\mathrm{ref}}\in[z_f+1,z_{\max}-3]$ \\
회복 정책군 & $(v_{\mathrm{esc}}/v_c,\vartheta)\in\{(2,0^\circ)\text{ 학습},(2,0^\circ)\text{ LQR},(2,\pm60^\circ),(3.5,0^\circ),(3.5,\pm60^\circ)\}$ \\
이산 필터 & $\Lambda=\{1,0.75,0.5,0.25,0\}$, 예측 3\,s(150스텝), 신규 계획 여유 $m_o=0.15$\,m·$m_c=0.2$\,m, 종단 $\Term_{\mathrm{hover}}\cup\Term_{\mathrm{flee}}$, 계획 유지(추격자 예측 재생), 호버 latch \\
bCBF-QP & $\alpha=4$, $\alpha_B=2$, 1.5\,s, 3스텝 간격, $w=(1,0.1,0.1,1)$, $w_s=10^6$, 건너뛰기 허용치 $10^{-2}$, 보조 필터 = 최종 이산 필터 \\
기준 접근 제어 & $v_{\mathrm{cr}}=6$\,m/s, $k_v=2.5$, 기울기 $50^\circ$ \\
잔차 정책 & 잔차 한계 추력 10\,\%, 각속도 25\,\%; 관측 135; SAC $\gamma=0.995$, 보상 척도 0.1, 병렬 64, 배치 256, 학습률 $3\times10^{-4}$, $\tau=0.005$, 온도 0.1(하한 $10^{-3}$), critic warm-up 20만, 고정점 벌점 0.2, 1000만 스텝 \\
보상 & $w_p=2$, $w_t=0.05$, $w_a=0.05$, 성공 $+100$, 접촉·이탈 $-150$ \\
PS2 & $w_f=0.5$, 잔차 정책에서 시작, critic 재적응 10만, 온도 고정, 300만 스텝 \\
평가 & 시드 3개, 장면 1000개(시드 9001), 15\,s와 10\,s 절단, Clopper--Pearson 95\,\%, 정확 McNemar \\\bottomrule
\end{tabularx}
\end{center}

\section{기호표}\label{app:symbols}
\begin{center}\small
\begin{tabularx}{\linewidth}{>{\raggedright\arraybackslash}p{3.4cm}X}\toprule
기호 & 의미 \\\midrule
$x=(p,v,q)$, $u=(a_T,\omega)$ & 기체 상태, 입력(추력 가속도, 기체 각속도) \\
$f_d$, $\Delta t$ & 이산 동역학, 시간 간격 \\
$\zeta$, $c_\ell$, $\nu_\ell$ & 장면 정보, 추격자 위치, 추격자 속도 \\
$r_d$, $r_c$, $v_c$, $a_c$ & 기체 반경, 추격자 반경, 추격자 속력·가속도 상한 \\
$h_f,h_o,h_w,h_z$ & 바닥, 장애물, 벽, 천장 여유 \\
$h_{c,j}$, $s_j^+$, $\delta$, $a_e$ & 추격자 여유도, 양의 접근 속도, 반응 지연, 유효 이탈 가속도 \\
$\Safe$, $\Base$, $\BaseP$ & 안전 집합, 필터의 기준 집합, Phase I 기준 집합 \\
$e$, $P$, $K$, $V$ & 호버 오차, Riccati 해, LQR 이득, $V=e^\top Pe$ \\
$a^\star$, $D$, $S$ & 최적 정지점, 정지점 행렬, Schur 여행렬 ($V_{\min}=r^\top Sr$) \\
$c_B$, $\rho_{xy}$, $z_{\mathrm{ext}}$ & 기준 집합 수준, 타원체의 수평·수직 범위 \\
$\pi_\theta$, $\Delta u_\theta$ & Phase I 학습 회복 정책과 그 잔차 \\
$\pi_b^{(j)}$, $\uesc^{(j)}$, $\sigma$ & 회복 정책 변형 $j$, 회피 입력, 반사 가중치 \\
$\Lambda$, $\lambda^\star$, $\bar\lambda$ & 혼합률 집합, 선택한 혼합률, 에피소드 평균 \\
$N$, $H$, $m_o$, $m_c$ & 예측 스텝(150), 종단 첨자($N+1$), 신규 계획 여유 \\
$\Term$, $p_{\mathrm{brake}}$ & 종단 집합, 제동 정지점 \\
$F$, $m$ & 안전 필터, 필터 기억(계획 유지) \\
$\upur$, $\Delta u_\xi$, $\unom$ & 기준 접근 입력, 임무 정책 잔차, 명목 입력 \\
$r_k$, $r^{\mathrm{PS2}}_k$, $\eta_k$, $w_f$ & 보상, PS2 보상, 정규화한 필터 수정량, 수정 비용 가중치 \\\bottomrule
\end{tabularx}
\end{center}

\section{원고 초안과 다른 점}\label{app:draft}
초안을 이 버전의 결과로 바꿀 때 확인할 항목이다.
\begin{enumerate}
\item \textbf{평가 규모}: 초안의 ``단일 시드, 128개 장면'' 수치(44.5, 82.8, 93.8, 100\,\% 등)는 이전 실험이다. 현재 표는 1000개 장면, 시드 2--3개이며 15\,s 에피소드와 10\,s 절단을 함께 보고한다.
\item \textbf{접촉 0건 표현}: PS2 정책 + 최종 필터는 접촉 1/2000(시드 0의 한 장면)이다. ``접촉이 관측되지 않았다'' 대신 ``2000회 중 1회(95\,\% 상한 0.28\,\%)''로 쓴다. 잔차 정책 + 최종 필터는 0/2000이다.
\item \textbf{15\,s 성공률}: 초안의 ``128개 장면 모두 성공(100\,\%)''은 현재 97.6\,\%이다(시간초과 2.4\,\%).
\item \textbf{PS2 정책의 학습 필터}: 초안 표 1은 단일 회피 필터로 학습한 PS2 정책을 썼다. 최종 파이프라인의 PS2 정책은 배치하는 최종 필터(회피 변형, 추격자 예측 재생)로 학습했다.
\item \textbf{bCBF slack 가중치}: 초안 식 (34) 뒤의 $w_s=10$은 $w_s=10^6$이다.
\item \textbf{신규 계획 여유}: 0.15\,m는 장애물과 벽뿐 아니라 바닥과 천장에도 적용된다.
\item \textbf{기준 집합의 추격자 조건}: 필터의 기준 집합 $\Base$에는 추격자 조건이 없고, 추격자는 종단 조건 $h_c\ge0.2$\,m와 회피 반사로 다룬다. Phase I의 기준 집합만 추격자 구 조건을 가진다.
\item \textbf{천장 조건}: 식 (20)의 $z_{\max}-a^\star_z\ge\rho_{\mathrm{top}}+r_d$에서 구현은 $\rho_{\mathrm{top}}=\rho_{xy}$이다(초안과 같음, 보수적).
\item \textbf{제동 정지점}: 식 (30)은 구현에서 수평 성분만 쓴다: $p_{\mathrm{brake}}=p_{H,xy}+v_{H,xy}\|v_{H,xy}\|/(2a_e)$. 이 점이 장애물과 벽에서 $r_d$ 이상 떨어져야 한다.
\item \textbf{추격자 운동}: 환경의 추격자는 속력 2\,m/s와 함께 가속도 2\,m/s$^2$ 상한을 가지며, 장애물 2\,m 이내에서 밀려난다(\S\ref{sec:arena}). 필터의 예측(가속도 제한 없는 순수 추적)과 다르다는 점을 밝혀 둔다.
\item \textbf{Phase I 수치}: 가중 성공률은 세 시드 평균 70.1$\pm$0.3\,\%(초안 69.9\,\%), LQR 43.8\,\%이다.
\item \textbf{이동 추격자 회복 학습 변형}: 초안 표 2의 해당 행은 최종 파이프라인 평가에 포함되지 않은 128개 장면 실험이다. 쓰려면 이전 실험임을 밝히거나 표에서 뺀다.
\item \textbf{수식 표기}: 초안 식 (2)·(6)의 $|\omega_i)$, $|\ell_i)$는 절댓값 $|\omega_i|$, $|\ell_i|$이고, 식 (4)의 괄호는 $v_{k+1}=v_k+\Delta t\,(a_{T,k}R(q_k)e_3-ge_3)$이다. 회피 각 $\vartheta$와 임무 정책 매개변수(초안 $\Delta u_\vartheta$)의 기호가 겹치므로 이 문서는 임무 정책 잔차를 $\Delta u_\xi$로 썼다.
\end{enumerate}

\section{한글(HWP) 수식 편집기 입력}\label{app:hwp}
본문의 주요 수식을 한글 수식 편집기 문법으로 옮긴 것이다(같은 내용: \file{results/strike\_paper/hwp\_equations.txt}). 붙여넣은 뒤 모양을 확인한다.
\IfFileExists{hwp_equations.txt}{\VerbatimInput[fontsize=\scriptsize,breaklines,breakanywhere,firstline=4]{hwp_equations.txt}}{\fbox{hwp\_equations.txt 없음}}

\end{document}
"
\providecommand{\paperdir}{../strike_reduced/final/paper}   % the report's results in this branch
\providecommand{\episodedir}{\paperdir/episodes}
\providecommand{\snapdir}{\paperdir/traj4/snapshots}   % filter snapshots
\documentclass[10pt,a4paper]{article}
\usepackage[margin=2.1cm]{geometry}
\usepackage{fontspec}
\setmainfont{Noto Serif CJK KR}
\setsansfont{Noto Sans CJK KR}
\setmonofont{Noto Sans Mono CJK KR}
\usepackage{amsmath,amssymb,bm,mathtools}
\usepackage{graphicx,float}
\usepackage{booktabs,array,tabularx,multirow}
\usepackage[dvipsnames]{xcolor}
\usepackage{enumitem}
\usepackage{caption,subcaption}
\usepackage{titlesec}
\usepackage{fancyhdr}
\usepackage[most]{tcolorbox}
\usepackage[ruled,vlined,linesnumbered]{algorithm2e}
\usepackage{fvextra}
\usepackage{tikz}
\usetikzlibrary{arrows.meta,positioning,fit,calc}
\usepackage[colorlinks=true,linkcolor=NavyBlue,urlcolor=NavyBlue,citecolor=NavyBlue]{hyperref}

\linespread{1.18}
\setcounter{tocdepth}{1}
\floatplacement{table}{H}
\setlength{\parindent}{1em}
\setlength{\parskip}{2pt}
\emergencystretch=2em
\setlist{nosep,leftmargin=1.5em}
\captionsetup{font=small,labelfont=bf,labelsep=period}
\renewcommand{\figurename}{그림}
\renewcommand{\tablename}{표}
\renewcommand{\contentsname}{목차}
\renewcommand{\refname}{참고 문헌}
\renewcommand{\abstractname}{요약}
\renewcommand{\appendixname}{부록}
\SetAlgorithmName{알고리즘}{알고리즘}{알고리즘 목록}
\SetKwInput{KwIn}{입력}
\SetKwInput{KwOut}{출력}
\SetKwInput{KwMem}{기억}
\SetKw{KwRet}{반환}
\SetKwIF{If}{ElseIf}{Else}{만약}{이면}{그렇지 않고 만약}{그렇지 않으면}{끝}
\SetKwFor{For}{각}{에 대해}{끝}
\titleformat{\section}{\large\bfseries\sffamily\color{NavyBlue}}{\thesection}{0.6em}{}
\titleformat{\subsection}{\normalsize\bfseries\sffamily}{\thesubsection}{0.6em}{}
\titleformat{\paragraph}[runin]{\normalsize\bfseries}{}{}{}[.]
\pagestyle{fancy}\fancyhf{}
\lhead{\small STRIKE: 방법과 결과 상세 정리}
\rhead{\small 2026.\ 10.\ 23}
\cfoot{\small\thepage}
\renewcommand{\headrulewidth}{0.3pt}

\newtcolorbox{keybox}[1]{colback=NavyBlue!4,colframe=NavyBlue!55,fonttitle=\bfseries\sffamily\small,title={#1},
  left=4pt,right=4pt,top=2pt,bottom=2pt,fontupper=\small}
\newtcolorbox{warnbox}[1]{colback=Orange!6,colframe=Orange!70,fonttitle=\bfseries\sffamily\small,title={#1},
  left=4pt,right=4pt,top=2pt,bottom=2pt,fontupper=\small}

\newcommand{\R}{\mathbb{R}}
\newcommand{\sat}{\operatorname{sat}}
\newcommand{\clip}{\operatorname{clip}}
\newcommand{\sdf}{\operatorname{sdf}}
\newcommand{\nrm}{\operatorname{normalize}}
\newcommand{\Safe}{\mathcal{S}}
\newcommand{\Base}{\bar{\mathcal{B}}}
\newcommand{\BaseP}{\mathcal{B}_{\mathrm{I}}}
\newcommand{\Term}{\mathcal{T}}
\newcommand{\unom}{u_{\mathrm{nom}}}
\newcommand{\upur}{u_{\mathrm{pur}}}
\newcommand{\uesc}{u_{\mathrm{esc}}}
\newcommand{\ubase}{u_{\mathrm{base}}}
\newcommand{\ustar}{u^{\star}}
\newcommand{\file}[1]{\texttt{\small #1}}
\newcommand{\figpath}[1]{\paperdir/figures/#1}
\newcommand{\tabin}[1]{\IfFileExists{\paperdir/tables/#1.tex}{\input{\paperdir/tables/#1.tex}}{\begin{center}\fbox{표 파일 없음: #1.tex}\end{center}}}
\newcommand{\figin}[2][\linewidth]{\IfFileExists{#2}{\includegraphics[width=#1]{#2}}{\fbox{그림 파일 없음}}}

\title{\vspace{-1.0cm}\textbf{\sffamily STRIKE: 다중 회복 정책 안전 필터와\\ 필터 포함 임무 학습(PS2)을 이용한 쿼드로터 저고도 표적 접근}\\[4pt]
\large 방법, 수식, 알고리즘, 평가 방법과 결과의 상세 정리 (논문 작성용)}
\author{김중원 \quad (정리 보조: Claude)}
\date{\small 2026년 10월 23일. 대상: 희소 경기장, 원래 추격자(반경 0.4 m, 2 m/s, 2 m/s$^2$)로 학습한 최종 구성}

\begin{document}
\maketitle
\thispagestyle{fancy}

\begin{abstract}
\noindent
장애물과 이동 추격자가 함께 있는 저고도 경기장에서 쿼드로터가 이동 표적의 0.6\,m 이내로 접근하는 문제(STRIKE)를 다룬다.
방법은 두 층으로 이루어진다.
\emph{안전 필터}는 임무 정책의 명목 입력을 일곱 개 회복 정책과 다섯 단계로 혼합한 35개 후보 각각을 한 스텝 실행한 뒤,
그 다음 상태에서 회복 정책을 3\,s 동안 가상으로 전개하여 바닥·장애물·벽·천장 제약과 상대 접근 속도를 반영한 추격자 여유도 $h_c$를 확인하고,
종단 상태가 호버하거나 제동할 수 있는 상태인지 검사한다. 통과한 후보 가운데 명목 입력을 가장 많이 유지하는 입력을 실행하고,
통과 후보가 없으면 직전에 인증한 계획을 그 추격자 예측 그대로 이어서 실행한다.
\emph{임무 정책}은 6\,m/s 기준 접근 제어에 잔차를 더하는 SAC 정책이며, 실제 필터를 포함한 환경에서 필터의 수정량을 벌점으로 받으며 재학습한다(PS2).
1000개 평가 장면과 15\,s 제한에서 PS2 정책과 최종 필터의 조합은 성공률 97.6\,\%, 접촉 2000회 중 1회(0.05\,\%, 95\,\% 상한 0.28\,\%)를 보였고,
같은 필터를 쓰는 잔차 정책(94.4\,\%, 접촉 0/2000)보다 성공률이 높았다(같은 장면 쌍 비교, 정확 McNemar 검정 $p<10^{-6}$).
수치는 현재 업로드된 평가(시드 0 전체, 시드 1 일부)에 기반하며, 세 시드 평가가 끝나면 같은 스크립트로 표와 그림이 갱신된다.
\end{abstract}

\tableofcontents
\vspace{4pt}

%======================================================================
\section{이 문서의 범위와 읽는 법}\label{sec:scope}

이 문서는 학술대회 원고(\emph{STRIKE 다중 회복 정책과 안전 필터 학습을 이용한 쿼드로터의 저고도 표적 접근}) 작성을 위해,
최종 구성의 수식과 알고리즘을 구현(\file{ps2rl/})과 대조하여 정리한 것이다.
모든 수식과 상수는 저장소의 코드와 학습에 사용된 설정 파일(\file{configs.json})에서 확인했다.
원고 초안과 다른 부분은 부록~\ref{app:draft}에, 주요 수식의 한글(HWP) 수식 편집기 입력은 부록~\ref{app:hwp}에 정리했다.

\begin{keybox}{이 버전의 구성 (원고의 ``최종 구성'')}
\begin{itemize}
\item \textbf{경기장}: 40\,m$\times$20\,m, 높은 장애물 3개, 이동 추격자 1개(반경 $r_c=0.4$\,m, 속력 $\le v_c=2$\,m/s, 가속도 $\le a_c=2$\,m/s$^2$), 이동 표적, 15\,s 제한.
\item \textbf{안전 필터}: 이산 다중 회복 필터. 회복 정책 7개(학습 회복 정책 또는 호버 LQR과 회피 반사의 결합), 혼합률 5단계, 3\,s 예측, 신규 계획 여유, 확장 종단 집합, 추격자 예측을 재생하는 계획 유지.
\item \textbf{임무 정책}: 6\,m/s 기준 접근 제어 + SAC 잔차(관측 135차원). PS2: 배치하는 필터를 그대로 포함한 환경에서 300만 스텝 재학습.
\item \textbf{평가}: 세 시드(전체 파이프라인을 시드별로 학습), 1000개 평가 장면(시드 9001), 15\,s 에피소드와 10\,s 절단 결과, Clopper--Pearson 구간과 같은 장면 쌍 비교.
\end{itemize}
\end{keybox}

\begin{figure}[H]\centering
\begin{tikzpicture}[node distance=5mm and 7mm, every node/.style={font=\small},
  box/.style={draw, rounded corners, align=center, minimum height=10mm, inner sep=4pt, fill=#1!8}, box/.default=gray,
  >={Stealth[length=2mm]}]
\node[box=blue] (p1) {Phase I\\ 학습 회복 정책 $\pi_\theta$\\ (안전 도착, TD3)};
\node[box=orange, right=of p1] (fam) {회복 정책군 $\{\pi_b^{(j)}\}_{j=0}^{6}$\\ 기본 회복 + 회피 반사\\ (속력비, 회피 각)};
\node[box=red, right=of fam] (filt) {이산 안전 필터 $F$\\ 35개 후보, 3\,s 전개,\\ 종단 검사, 계획 유지};
\node[box=gray, right=of filt] (plant) {쿼드로터,\\ 장애물, 추격자,\\ 표적};
\node[box=green, below=8mm of filt] (str) {임무 정책\\ $\unom=\sat(\upur+\Delta u_\xi)$};
\draw[->] (p1) -- (fam);
\draw[->] (fam) -- (filt);
\draw[->] (str) -- node[right]{$\unom$} (filt);
\draw[->] (filt) -- node[above]{$u$} (plant);
\draw[->] (plant) |- node[pos=0.75,below]{$o$ (상태, 장면)} (str);
\node[draw, dashed, rounded corners, fit=(filt)(str)(plant), inner sep=5pt,
  label={[font=\scriptsize]below:PS2: 필터를 포함한 환경에서 임무 정책을 학습($r^{\mathrm{PS2}}=r-w_f\|\eta\|^2$)}] {};
\end{tikzpicture}
\caption{전체 구성. 안전 필터는 회복 가능성을, 임무 정책은 표적 접근 성능을 담당한다.}\label{fig:overview}
\end{figure}

\paragraph{재현 경로} 표·그림·본문 수치는 \file{docker\_batch/run\_strike\_paper.sh}가 만든다.
\file{STEPS=tables}는 몇 초 안에 표(\file{.md/.csv/.tsv/.tex}), 그림 7종, 본문 인용 수치(\file{numbers.md})를 만든다.
나머지 단계는 대표 에피소드를 다시 실행해 2차원·3차원 궤적 그림, 상세 신호 그림, 2차원·3차원 GIF를 만든다(\S\ref{sec:results-episodes}).

%======================================================================
\section{문제 설정}\label{sec:problem}

\subsection{기체 모델}
기체 상태는 위치 $p\in\R^3$, 관성좌표 속도 $v\in\R^3$, 자세 쿼터니언 $q\in S^3$로 이루어지고,
입력은 질량으로 나눈 총추력(추력 가속도) $a_T$와 기체 각속도 명령 $\omega$이다.
\begin{equation}
x=(p,v,q),\qquad u=(a_T,\omega),\qquad 0\le a_T\le 4g=39.24\ \mathrm{m/s^2},\qquad |\omega_i|\le 18\ \mathrm{rad/s}.
\label{eq:state}
\end{equation}
시간 간격 $\Delta t=0.02$\,s(50\,Hz)의 명시적 Euler 적분으로 상태를 갱신한다. $R(q)$는 기체 좌표에서 관성좌표로의 회전행렬, $e_3=(0,0,1)^\top$, $\otimes$는 쿼터니언 곱이다.
\begin{align}
p_{k+1}&=p_k+\Delta t\,v_k, \label{eq:dyn-p}\\
v_{k+1}&=v_k+\Delta t\,\big(a_{T,k}\,R(q_k)e_3-g\,e_3\big), \label{eq:dyn-v}\\
\tilde q_{k+1}&=q_k+\tfrac{\Delta t}{2}\,q_k\otimes(0,\omega_k),\qquad q_{k+1}=\tilde q_{k+1}/\|\tilde q_{k+1}\|. \label{eq:dyn-q}
\end{align}
이를 $x_{k+1}=f_d(x_k,u_k)$로 쓴다. 입력은 항상 상자 $\mathcal U$로 포화시킨다($\sat$).
입력의 정규화 척도는 $u_{\mathrm{scale}}=(a_{\max},\omega_{\max},\omega_{\max},\omega_{\max})=(39.24,18,18,18)$이다.
이 모델은 PX4 오프보드 제어가 받는 추력--각속도 인터페이스에 맞춘 운동 모델이며, 모터 동특성, 지연, 공력, 추정 오차는 포함하지 않는다.

\subsection{경기장과 장면 생성}\label{sec:arena}
\begin{itemize}
\item \textbf{경기장}: $x\in[0,40]$\,m, $|y|\le10$\,m, 기준 바닥 고도 $z_f=1$\,m, 상한 고도 $z_{\max}=10$\,m. 기체 반경 $r_d=0.3$\,m.
\item \textbf{장애물}: 높은 기둥 3개. 확률 0.5로 원기둥(반경 $\mathcal U[0.8,2.0]$\,m), 아니면 직사각 기둥(반폭 $\mathcal U[1.5,3.0]\times\mathcal U[1.0,2.5]$\,m, 회전 $\pm30^\circ$). 중심 $x\in[7,31]$\,m.
  장애물 간 간격 3.5\,m 이상, 출발점과 표적에서 4\,m 이상, 옆벽에서 1\,m 이상 떨어지게 하고, 장애물과 옆벽에서 1.5\,m 넘게 떨어진 0.5\,m 격자점만으로 출발점과 표적을 이을 수 있는 장면만 사용한다.
\item \textbf{출발}: $x=2$\,m, $y\in[-3,3]$\,m, $z\in[2,4]$\,m, 속력 $\mathcal U[4,6]$\,m/s(표적 방향 $\pm15^\circ$), 기울기 $10^\circ$ 이하.
\item \textbf{표적}: $x\in[32,37]$\,m, $|y|\le5$\,m, $z\in[1.6,2.5]$\,m에서 시작해 수평으로 $\mathcal U[0.2,0.8]$\,m/s의 일정 속도로 떠돌며, $x\in[32,39]$, $|y|\le9$의 영역 경계에서 반사된다.
\item \textbf{추격자}: 출발점--표적 직선의 $35$--$85\,\%$ 지점에서 옆으로 $5$--$9$\,m 떨어진 곳(높이 $2$--$4$\,m)에서 시작한다.
\end{itemize}

\paragraph{추격자 운동(환경)} 추격자 $j$의 위치 $c_j$, 속도 $\nu_j$는 장애물을 밀어내는 순수 추적으로 움직인다.
\begin{align}
w&=\frac{p-c_j}{\|p-c_j\|}+2\varpi_c\begin{bmatrix}\nabla\sdf_O(c_{j,xy})\\0\end{bmatrix},\qquad
\varpi_c=\clip\!\Big(\frac{d_{\mathrm{rep}}-\sdf_O(c_{j,xy})}{d_{\mathrm{rep}}},0,1\Big),\quad d_{\mathrm{rep}}=2\ \mathrm m,\\
\nu_{j,k+1}&=\Pi_{v_c}\!\Big(\nu_{j,k}+\Delta t\,\Pi_{a_c}\big(\tfrac{v_c w/\|w\|-\nu_{j,k}}{\Delta t}\big)\Big),\qquad c_{j,k+1}=c_{j,k}+\Delta t\,\nu_{j,k+1},
\end{align}
여기서 $\Pi_r(\cdot)$는 노름을 $r$ 이하로 줄이는 연산이며, 높이는 $[z_f+r_c,z_{\max}]$로 제한한다.
필터는 이 운동을 알지 못하고, 거리 감소에 가장 불리한 \emph{가속도 제한 없는 순수 추적}(식~\eqref{eq:pursue})으로 예측한다.

\subsection{장애물의 부호 거리}
장애물 $i$는 중심 $c_i$, 반폭 $b_i=(b_x,b_y)$, 회전 $\psi_i$, 모서리 반경 $\rho_i$로 표현한다(원기둥은 $b_i=0$, $\rho_i$=반경, 직사각 기둥은 $\rho_i=0$).
$\ell=R(\psi_i)^\top(p_{xy}-c_i)$로 두면
\begin{equation}
\sdf_i(p_{xy})=\big\|\max(|\ell|-b_i,0)\big\|_2+\min\!\big(\max(|\ell_x|-b_x,\ |\ell_y|-b_y),\,0\big)-\rho_i,\qquad
\sdf_O=\min_i\sdf_i ,
\label{eq:sdf}
\end{equation}
이며 $|\cdot|$와 $\max$는 성분별로 적용한다. 밖에서 양수, 표면에서 0, 안에서 음수이다.

\subsection{종료 조건과 성과 지표}\label{sec:outcomes}
매 스텝 다음 사건을 검사한다(같은 스텝에서 접촉이 우선).
\begin{itemize}
\item \textbf{성공}(strike): $\|p-p_t\|<0.6$\,m.
\item \textbf{접촉}(contact): 장애물 $\sdf_O(p_{xy})<r_d$, 바닥 $z<z_f$, 추격자 $\|p-c_j\|<r_c+r_d=0.7$\,m 중 하나.
\item \textbf{영역 이탈}(out): $x\notin[-1,41]$, $|y|>11$ 또는 $z>z_{\max}$. 임무 실패이지만 접촉으로 세지 않는다.
\item \textbf{시간초과}(time-up): 15\,s(750스텝) 동안 위 사건이 없음. 10\,s 결과는 같은 에피소드를 10\,s에서 절단한 것이다.
\end{itemize}

%======================================================================
\section{안전 집합}\label{sec:safe}

\subsection{기하 제약}
\begin{equation}
h_f=z-z_f,\qquad h_o=\sdf_O(p_{xy})-r_d,\qquad h_w=d_{\mathrm{wall}}(p_{xy})-r_d,\qquad h_z=z_{\max}-z-r_d,
\label{eq:geo}
\end{equation}
여기서 $d_{\mathrm{wall}}$은 네 외벽($x=0$, $x=40$, $y=\pm10$)까지의 최소 거리다.
바닥 조건은 과제의 고도 정의($z\ge z_f$)를 그대로 따르므로 $r_d$를 빼지 않는다.

\subsection{추격자 여유도 $h_c$의 유도}\label{sec:hc}
추격자 $j$에 대해 거리와 시선(추격자에서 기체로 향하는 단위벡터)을
\begin{equation}
R_j=\|p-c_j\|,\qquad \hat e_j=\frac{p-c_j}{R_j}
\end{equation}
로 두면 기체가 거리를 벌리는 속도는 $\hat e_j^\top v$이다. 추격자는 거리를 최대 $v_c$로 줄이므로(기체를 정면으로 향할 때) 최악의 접근 속도는
\begin{equation}
s_j=v_c-\hat e_j^\top v,\qquad s_j^+=\max(s_j,0)
\label{eq:closing}
\end{equation}
이다. 기체가 반응 지연 $\delta$ 뒤에 유효 이탈 가속도 $a_e$로 거리를 벌리는 속도를 높인다고 하자.
지연 동안에는 거리를 최대 $s_j^+\delta$ 잃는다. 그 뒤 접근 속도가 $s(t)=s_j^+-a_e t$로 줄어 0이 될 때까지 잃는 거리는
\begin{equation}
\int_0^{s_j^+/a_e}(s_j^+-a_e t)\,dt=\frac{(s_j^+)^2}{2a_e}
\qquad\Big(\text{등가로}\ 0=(s_j^+)^2-2a_e d_{\mathrm{brake}}\Big)
\end{equation}
이다. 충돌 반경과 두 손실을 현재 거리에서 빼면 추격자 여유도를 얻는다.
\begin{equation}
\boxed{\;h_{c,j}=R_j-(r_c+r_d)-s_j^+\,\delta-\frac{(s_j^+)^2}{2a_e}\;},\qquad \delta=0.2\ \mathrm s,\quad a_e=5\ \mathrm{m/s^2}.
\label{eq:hc}
\end{equation}
$a_e$는 기체의 최대 가속도보다 훨씬 작게 잡은 값이다. 요구 거리 $d_{\mathrm{req}}=r_c+r_d+s^+\delta+(s^+)^2/(2a_e)$의 예는 다음과 같다.

\begin{center}\small
\begin{tabular}{lccl}\toprule
상황 & $s^+$ [m/s] & $d_{\mathrm{req}}$ [m] & 구성 \\\midrule
추격자로부터 $v_c$ 이상으로 멀어짐 & 0 & 0.7 & 충돌 반경만 \\
정지(호버) & 2 & 1.5 & $0.7+0.4+0.4$ \\
추격자 쪽으로 6\,m/s 정면 접근 & 8 & 8.7 & $0.7+1.6+6.4$ \\\bottomrule
\end{tabular}
\end{center}

\noindent
같은 거리 4\,m라도 정지 중이면 $h_c=2.5$\,m이지만 6\,m/s로 정면 접근 중이면 $h_c=-4.7$\,m이다.
즉 아직 충돌하지 않은 상태도 안전 집합 밖일 수 있고, 이것이 거리만 보는 조건과의 차이다.

\paragraph{열린 공간에서의 불변성(근사)} 지연 뒤 기체가 $\hat e_j$ 방향으로 $a_e$의 가속을 유지하고 추격자가 정면으로 $v_c$로 다가오면 $\dot R_j=-s_j$, $\dot s_j=-a_e$이므로 $s_j>0$에서
\begin{equation}
\dot h_{c,j}=\dot R_j-\dot s_j\,\delta-\frac{s_j\dot s_j}{a_e}=-s_j+a_e\delta+s_j=a_e\delta\ \ge0
\end{equation}
이다. 따라서 막힌 곳이 없으면 $\{h_c\ge0\}$은 ``도피''에 대해 유지된다. 시선 회전에 의한 항(크기 $\lesssim v_c^2/(4R)$)은 $a_e$를 실제 최대보다 작게 잡아 흡수한다고 본다.
장애물과 벽이 도피 방향을 막을 수 있다는 점이 이 식만으로는 부족한 이유이고, 필터가 회복 궤적을 직접 검사하는 이유이다.

\subsection{안전 집합}
장면 정보 $\zeta$(장애물과 추격자 상태)에 대해
\begin{equation}
\Safe(\zeta)=\big\{x:\ h_f\ge0,\ h_o\ge0,\ h_w\ge0,\ h_z\ge0,\ h_{c,j}\ge0\ \ \forall j\big\}.
\label{eq:safe}
\end{equation}

\subsection{비교: 구형 근사(초기 방식)}\label{sec:ball}
초기 방식은 추격자를 현재 위치에 고정된 반경 $R_{\mathrm{ball}}=r_c+r_d+v_cT=3.7$\,m($T=1.5$\,s)의 구로 보고 $\|p-c_j\|\ge R_{\mathrm{ball}}$을 요구한다.
$T$초 안에 추격자가 닿을 수 있는 영역을 덮지만, 그 뒤 추격자의 위치와 기체의 접근 속도를 반영하지 못한다.
Phase I 학습(\S\ref{sec:phase1})은 이 구형 근사를 사용하며, 구형 근사 필터는 비교 실험(표~\ref{tab:table1_15s}, \ref{tab:table3_ablation})에 쓰인다.

%======================================================================
\section{기준 집합: 호버 LQR}\label{sec:base}

\subsection{선형화와 LQR}
호버 목표점 $a$ 주변에서 오차 상태 $e=(p-a,\ v,\ \phi)\in\R^9$를 쓴다. $\phi=2\,\mathrm{sgn}(q_w)\,q_{xyz}$는 작은 회전벡터(자세 좌표)이다.
호버 입력 $\ustar=(g,0,0,0)$ 주변의 연속시간 선형 모델은
\begin{equation}
\dot p=v,\qquad \dot v_x=g\,\phi_y,\qquad \dot v_y=-g\,\phi_x,\qquad \dot v_z=\Delta a_T,\qquad \dot\phi=\omega
\end{equation}
이고, Euler 이산화 $A_d=I+\Delta tA$, $B_d=\Delta tB$에 대해 이산 LQR을 설계한다.
\begin{equation}
e_{k+1}=A_de_k+B_d\Delta u_k,\qquad \Delta u_k=-Ke_k,\qquad K=(R_u+B_d^\top PB_d)^{-1}B_d^\top PA_d,\qquad V(e)=e^\top Pe.
\label{eq:lqr}
\end{equation}
$P$는 이산 Riccati 방정식의 해이다. 가중치는 $Q=\mathrm{diag}(2,2,8,\ 0.01,0.01,0.5,\ 8,8,2)$, $R_u=\mathrm{diag}(0.06,0.15,0.15,0.12)$이다.

\subsection{정지점의 재설정}
속도가 있는 기체는 제자리에 멈출 수 없으므로, $V$를 최소로 만드는 목표점을 해석적으로 구한다.
$r=(v,\phi)$로 두고 $P$를 위치와 나머지로 분할하면
\begin{equation}
V=(p-a)^\top P_{pp}(p-a)+2(p-a)^\top P_{pr}r+r^\top P_{rr}r .
\end{equation}
$\partial V/\partial a=-2P_{pp}(p-a)-2P_{pr}r=0$에서
\begin{equation}
a^\star(x)=p+Dr,\quad D=P_{pp}^{-1}P_{pr};\qquad
V_{\min}(r)=\min_aV=r^\top Sr,\quad S=P_{rr}-P_{rp}P_{pp}^{-1}P_{pr}
\label{eq:anchor}
\end{equation}
이다($S$는 Schur 여행렬). 이때 $e=(-Dr,r)$이므로 재설정 LQR은 $u=\ustar-\tilde Kr$, $\tilde K=K_r-K_pD$로, 어디서든 목표점 없이 정지하는 제어기가 된다.
설계값에서 $D$의 속도 성분은 약 0.385\,s로, 수평 속도 1\,m/s마다 정지점이 0.385\,m 앞에 놓인다.

\subsection{호버 타원체와 기준 집합}
$\{e:e^\top Pe\le c\}$에서 위치 성분의 방향 $n$으로의 최대 변위는 Cauchy--Schwarz 부등식으로 $\sqrt{c\,n^\top[P^{-1}]_{pp}n}$이다. 따라서
\begin{equation}
\rho_{xy}=\sqrt{c_B\,\lambda_{\max}\big([P^{-1}]_{xy}\big)}=1.22\ \mathrm m,\qquad
z_{\mathrm{ext}}=\sqrt{c_B\,[P^{-1}]_{zz}}=0.63\ \mathrm m,\qquad c_B=50 .
\label{eq:extent}
\end{equation}
필터가 쓰는 기준 집합은 정적 환경의 호버 회복 가능 상태로, 추격자 조건을 포함하지 않는다(추격자는 회피 반사와 종단 조건으로 다룬다, \S\ref{sec:family}, \ref{sec:terminal}).
\begin{equation}
\Base=\left\{x\ \middle|\
\begin{aligned}
&h_f,h_o,h_w,h_z\ge0,\qquad r^\top Sr\le c_B,\\
&\sdf_O(a^\star_{xy})\ge\rho_{xy}+r_d,\qquad d_{\mathrm{wall}}(a^\star_{xy})\ge\rho_{xy}+r_d,\\
&a^\star_z-z_{\mathrm{ext}}\ge z_f,\qquad z_{\max}-a^\star_z\ge\rho_{xy}+r_d
\end{aligned}\right\}.
\label{eq:base}
\end{equation}
즉 충분히 느리고($V_{\min}\le c_B$), LQR이 정착하는 동안 움직일 수 있는 타원체 전체가 장애물·바닥·벽에서 떨어져 있다.
천장 조건은 구현대로 $\rho_{xy}$를 쓰며, $\rho_{xy}\ge z_{\mathrm{ext}}$이므로 보수적이다.
Phase I의 기준 집합 $\BaseP$는 벽과 천장 대신 추격자 구 조건 $\|a^\star-c_j\|\ge R_{\mathrm{ball}}+\rho$($\rho=1.22$\,m, 3차원 범위)를 쓴다.

\subsection{$c_B$의 인증}
목표점을 $a_0$에 고정한 LQR $u=\ustar-K(p-a_0,r)$이 $\{V\le c_B\}$를 유지하려면 $c_B$가 다음 세 한계보다 작아야 한다.
\begin{itemize}
\item 입력 한계: 타원체 위에서 $\max|K_ie|=\sqrt{c\,K_iP^{-1}K_i^\top}\le m_i$이므로 $c_U=\min_i m_i^2/(K_iP^{-1}K_i^\top)=56.1$ ($m=(g,18,18,18)$).
\item 자세 좌표: $|\phi|<2$이면 좌표가 일대일이므로 $c_{\mathrm{chart}}=4/\lambda_{\max}([P^{-1}]_{\phi\phi})=102.1$.
\item 비선형 감소: 비선형 Euler 모델에서 $\{V=c\}$ 위 4000개 표본의 한 스텝 감소 $V^+/V<1$을 확인한다. $c=\min(c_U,c_{\mathrm{chart}})=56.1$에서 이미 모든 표본이 감소하므로(최악 방향 0.3\,m/s$^2$ 바람 포함) $c_{\mathrm{Lyap}}\ge56.1$이다. 감소하지 않으면 이분법으로 줄인다.
\end{itemize}
$c_B=50$에서 표본 최대 $V^+/V$는 0.980(바람 포함 0.986)이다. $\Base$ 안에서 수평 속력은 2.67\,m/s, 수직 속력은 2.12\,m/s, 기울기는 $52.7^\circ$ 이하이다.
표본 검사이므로 연속 상태공간 전체의 엄밀한 증명은 아니다.

\subsection{호버 latch}
계획이 $\Base$에 들어가고 회피 반사가 꺼져 있으면($\sigma=0$, \S\ref{sec:family}) 그때의 목표점 $a_0$을 고정하고 고정 목표점 LQR로 호버한다.
인증이 다루는 제어기가 바로 이 고정 목표점 LQR이기 때문이다. 매 스텝 재설정하면 목표점이 밀리미터 단위로 흔들려, 경계에서 들어온 상태가 다시 $\Base$ 밖으로 나가는 일이 생긴다.
고정된 호버는 $V(p-a_0,r)\le c_B$인 동안 종단 조건으로 인정된다.

%======================================================================
\section{Phase I: 안전 도착 회복 정책}\label{sec:phase1}

\paragraph{정책} LQR 입력에 잔차를 더한다.
\begin{equation}
\pi_\theta(x,\zeta)=\sat\big(u_{\mathrm{LQR}}(x)+\Delta u_\theta(o)\big),\qquad u_{\mathrm{LQR}}(x)=\ustar-\tilde Kr .
\label{eq:phase1-policy}
\end{equation}
관측 $o\in\R^{98}$은 $z-z_f$, $v$, $q$(8개), 가까운 순서의 장애물 슬롯 8개(최근접점 벡터 2, 중심 오프셋 2, 반폭 2, $\cos2\psi$, $\sin2\psi$, 모서리 반경, 존재 여부: 각 10개), 추격자 구 슬롯 2개(중심 오프셋 3, 반경 $R_{\mathrm{ball}}$, 존재 여부: 각 5개)로 구성된다.
슬롯 순서는 회복 전개를 시작할 때의 순서로 고정하여 전개 중 입력이 불연속으로 바뀌지 않게 한다. $\Safe$와 $\BaseP$의 소속은 $o$의 정확한 함수이다.

\paragraph{목적함수와 Bellman 관계} 에피소드는 표본 (상태, 장면)에서 시작한 회복 전개 하나이며, $\BaseP$ 도달(성공), $\Safe$ 이탈(실패), 112스텝(2.24\,s) 중 먼저 일어나는 사건에서 끝난다.
\begin{equation}
J(\theta)=\mathbb E\big[\beta^{\tau_B}\,\mathbf 1\{\tau_B<\tau_{\mathrm{out}}\}\big],\qquad \beta=0.99 .
\label{eq:sa}
\end{equation}
장면을 포함한 상태 $y=(x,\zeta)$와 결합 전이 $y^+=G(y,u)$에 대해 최적값은
\begin{equation}
W^\star(y)=1\ (y\in\BaseP\cap\Safe),\qquad W^\star(y)=0\ (y\notin\Safe),\qquad W^\star(y)=\beta\max_{u\in\mathcal U}W^\star(G(y,u))\ \ (\text{그 밖})
\end{equation}
를 만족하고, $W^\star(y)=\beta^{d(y)}$($d$는 안전하게 도착하는 최소 스텝 수)이다. TD3의 critic 목표는 이 관계를 그대로 쓴다.
\begin{equation}
y_{\mathrm{TD}}=b(o)+\beta\,c(o)\,\min_{i=1,2}Q'_i\big(o',\mu'(o')\big),\qquad b=\mathbf 1\{\BaseP\},\ c=\mathbf 1\{\Safe\setminus\BaseP\},
\end{equation}
여기서 $Q'(o',\cdot)$는 $o'\in\BaseP$이면 1, $o'\notin\Safe$이면 0으로 둔다. 엔트로피 항은 critic 목표에 넣지 않는다(도착 순서 ``빨리 $>$ 늦게 $>$ 실패''를 보존하기 위함).
actor 손실은 $-\mathbb E[\min_iQ_i(o,\Delta u_\theta(o))]+0.2\,\mathbb E\big[\tfrac14\|\Delta u/u_{\mathrm{scale}}\|^2\big]$이고, 둘째 항이 잔차를 LQR 쪽으로 붙잡는다. actor는 LQR 모방($\Delta u\approx0$)으로 초기화한다.

\paragraph{학습 분포} 일반(속력 $\le10$\,m/s, 기울기 $\le45^\circ$, 장애물 0--6개, 추격자 구 0--2개), 정지 경로 막힘(속력 $\ge3$\,m/s, LQR 정지점 부근에 장애물), 추격자 근접(속력 $\ge3$\,m/s, LQR 정지점 부근에 추격자 구), 호버 근방(느리고 정지점이 비어 있음) 영역을 가중치 $(0.4,0.2,0.2,0.2)$로 섞는다.
평가에는 저속 상태 근처에 추격자 구가 있는 저속 탈출 영역을 추가로 쓴다(가중치 0). 영역마다 5만 개 시작점을 미리 뽑고, 검증 1000개와 시험 2000개는 따로 뽑는다.

\paragraph{평가 지표와 결과} 시험 상태가 $C_N$에 속하는 비율 $\hat m$을 쓴다. $C_N$은 $N=T/\Delta t=75$스텝(1.5\,s) 안에 $\Safe$를 벗어나지 않고 $\BaseP$에 처음 도달하는 상태의 집합이다.
학습 회복 정책은 가중 평균 $70.1\pm0.3\,\%$(세 시드)로 호버 LQR의 $43.8\,\%$보다 높다(표~\ref{tab:table4_phase1}, 그림~\ref{fig:phase1}).
이 평가는 추격자를 정지한 구로 근사한 Phase I 조건이며, 최종 필터의 이동 추격자 성능과는 구분해야 한다.

%======================================================================
\section{회복 정책군}\label{sec:family}

\paragraph{기본 회복 입력} 기준 집합 안이거나 호버가 고정되어 있으면 고정 목표점 LQR, 밖이면 학습 회복 정책(또는 호버 LQR 변형에서는 재설정 LQR)을 쓴다.
\begin{equation}
\ubase^{(j)}(x)=\begin{cases}\ustar-K(p-a_0,r) & x\in\Base\ \text{또는 호버 고정},\\ \pi_\theta(x,\zeta_{\mathrm{view}}) & \text{그 밖(학습 변형)},\\ u_{\mathrm{LQR}}(x) & \text{그 밖(LQR 변형)}.\end{cases}
\end{equation}
호버가 고정되기 전에는 $a_0=a^\star(x)$이므로 첫 줄은 재설정 LQR과 같다. 학습 회복 정책의 장면 입력 $\zeta_{\mathrm{view}}$는 계획을 만든 시점의 장면(가까운 장애물 8개와 추격자 구)으로 고정한다. Phase I에서 학습한 범위(장애물 10\,m, 추격자 12\,m 이내)만 넣는다.

\paragraph{회피 반사} 최소 추격자 여유도 $h_{\min}=\min_jh_{c,j}$에 따라 회피 입력을 섞는다.
\begin{equation}
\sigma(x)=\clip\!\Big(\frac{h_{\mathrm{off}}-h_{\min}(x)}{h_{\mathrm{off}}-h_{\mathrm{on}}},0,1\Big),\qquad
\pi_b^{(j)}(x)=\sat\big((1-\sigma)\,\ubase^{(j)}(x)+\sigma\,\uesc^{(j)}(x)\big),\qquad h_{\mathrm{on}}=0.5,\ h_{\mathrm{off}}=2.0\ \mathrm m.
\label{eq:reflex}
\end{equation}
$h_{\min}\ge2$\,m이면 기본 회복 입력만, $h_{\min}=1.25$\,m이면 절반씩, $h_{\min}\le0.5$\,m이면 회피 입력만 쓴다. 반사가 작동하면($\sigma>0$) 호버 고정을 푼다.

\paragraph{회피 입력} 가장 위험한 추격자 $j^\star=\arg\min_jh_{c,j}$의 반대 방향 $\hat e_{xy}$를 변형 각 $\vartheta_j$만큼 돌리고, 정지점 $a^\star$ 주변의 장애물과 벽의 바깥 법선으로 민다.
\begin{align}
d&=\nrm\Big(R(\vartheta_j)\,\hat e_{xy}+2\varpi_o\,n_o+2\textstyle\sum_{i=1}^4\varpi_{w,i}\,n_{w,i}\Big),\qquad
\varpi=\clip\!\Big(\frac{d_{\mathrm{push}}-(\mathrm{dist}(a^\star_{xy})-\rho_{xy}-r_d)}{d_{\mathrm{push}}},0,1\Big),\\
v_{\mathrm{des}}&=\big(v_{\mathrm{esc},j}\,d,\ 1.5\,(z_{\mathrm{ref}}-z)\big),\qquad z_{\mathrm{ref}}=\clip(z,\ z_f+1,\ z_{\max}-3),
\end{align}
여기서 $n_o=\nabla\sdf_O(a^\star_{xy})/\|\nabla\sdf_O\|$, $n_{w,i}$는 벽의 안쪽 법선, $\mathrm{dist}$는 각 장애물·벽까지의 거리, $d_{\mathrm{push}}=1.5$\,m이다.
속도 루프는 $a=k_v(v_{\mathrm{des}}-v)$(수평 성분은 $\|a_{xy}\|\le g\tan45^\circ$로 제한), $a_{\mathrm{des}}=(a_{xy},a_z+g)$에서
$a_T=a_{\mathrm{des}}^\top z_b$, $\omega=k_{\mathrm{att}}R^\top(z_b\times\hat a_{\mathrm{des}})-k_\psi\psi\,e_3$로 입력을 만들고 포화시킨다($k_v,k_{\mathrm{att}},k_\psi)=(2,10,3)$, $z_b=R(q)e_3$, $\psi$는 요각).
추격자 여유가 충분할 때는 학습 회복 정책이 장애물을 처리하고, $h_c$가 줄어들수록 회복 정책은 ``$v_{\mathrm{esc}}$로 도피''로 바뀐다. 이는 열린 공간에서 $h_c$를 줄이지 않는 행동이다(\S\ref{sec:hc}).

\paragraph{일곱 변형} 회피 속력비 $v_{\mathrm{esc}}/v_c$와 회피 각 $\vartheta$의 조합으로 회복 정책을 늘린다(표~\ref{tab:variants}).
필터는 어느 하나라도 인증하는 후보를 받아들이므로, 회복 가능 집합은 각 변형 회복 가능 집합의 합집합이 된다.

\begin{table}[h]\centering\small
\caption{필터의 회복 정책 일곱 변형. 번호는 동률일 때의 우선순위이다.}\label{tab:variants}
\begin{tabular}{cccc}\toprule
번호 $j$ & 기본 회복 입력(기준 집합 밖) & $v_{\mathrm{esc}}/v_c$ & $\vartheta$ \\\midrule
0 & 학습 회복 정책 $\pi_\theta$ & 2 & $0^\circ$ \\
1 & 재설정 호버 LQR & 2 & $0^\circ$ \\
2, 3 & 학습 회복 정책 & 2 & $+60^\circ$, $-60^\circ$ \\
4 & 학습 회복 정책 & 3.5 & $0^\circ$ \\
5, 6 & 학습 회복 정책 & 3.5 & $+60^\circ$, $-60^\circ$ \\\bottomrule
\end{tabular}
\end{table}

%======================================================================
\section{이산 다중 회복 안전 필터}\label{sec:filter}



\subsection{후보 입력}
임무 정책의 명목 입력 $\unom$을 각 회복 정책의 현재 입력과 섞어 후보를 만든다.
\begin{equation}
u^{(j)}_\lambda=\sat\big(\lambda\,\unom+(1-\lambda)\,\pi_b^{(j)}(x)\big),\qquad \lambda\in\Lambda=\{1,0.75,0.5,0.25,0\},\quad j=0,\dots,6 .
\label{eq:cand}
\end{equation}
$7\times5=35$개 후보를 검사한다($\lambda=1$에서는 일곱 개가 같은 입력이다).
이는 입력 집합 전체를 최적화하는 것이 아니라 명목 입력과 각 회복 입력을 잇는 선분 위의 유한한 점을 검사하는 것이다.

\subsection{회복 궤적의 전개와 추격자 예측}
후보를 한 스텝 실행한 다음 상태에서 시작해 그 변형의 회복 정책을 $N=150$스텝(3\,s) 전개한다. 추격자는 예측 기체를 향해 $v_c$로 직진한다고 가정한다(거리 감소에 가장 불리한 운동).
\begin{align}
x_1&=f_d\big(x,u^{(j)}_\lambda\big),\qquad x_{i+1}=f_d\big(x_i,\ \pi_b^{(j)}(x_i;c_i)\big),\qquad i=1,\dots,N, \label{eq:rollout}\\
c_{\ell,i+1}&=c_{\ell,i}+\min\!\big(v_c\Delta t,\ \|p_i-c_{\ell,i}\|\big)\,\frac{p_i-c_{\ell,i}}{\|p_i-c_{\ell,i}\|},\qquad c_{\ell,0}=c_\ell(\text{현재}),\qquad c_{\ell,i,z}\ge z_f+r_c . \label{eq:pursue}
\end{align}
$\ell$은 추격자 번호이다. 회피 반사는 예측 추격자 위치 $c_i$를 보고, 학습 회복 정책의 장면 입력은 계획 시점의 것으로 고정된다. 종단 첨자는 $H=N+1$이다.

\subsection{궤적 조건(신규 계획 여유)}
새 계획은 모든 예측 노드 $i=1,\dots,H$에서 여유를 더한 안전 집합에 있어야 한다.
\begin{equation}
\min\{h_f,h_o,h_w,h_z\}(x_i)\ge m_o=0.15\ \mathrm m,\qquad \min_\ell h_{c,\ell}(x_i;c_{\ell,i})\ge m_c=0.2\ \mathrm m .
\label{eq:fresh}
\end{equation}
여유는 반사 방향의 재조준과 실제 추격자가 예측과 다르게 움직이는 것에 대비한 것이다. 계획 유지(\S\ref{sec:commit})의 재검사는 여유 없이($m_o=m_c=0$) 안전 집합만 요구한다.

\subsection{종단 집합}\label{sec:terminal}
종단 상태 $x_H$는 호버하거나 직선 제동으로 멈출 수 있어야 한다.
\begin{align}
\Term_{\mathrm{hover}}&=\big\{x_H:\ \text{호버 고정이면 } V(p_H-a_0,r_H)\le c_B,\ \text{아니면 } x_H\in\Base\big\},\\
\Term_{\mathrm{flee}}&=\big\{x_H:\ \sdf_O(p_{\mathrm{brake}})\ge r_d,\ d_{\mathrm{wall}}(p_{\mathrm{brake}})\ge r_d\big\},\qquad
p_{\mathrm{brake}}=p_{H,xy}+\frac{v_{H,xy}\|v_{H,xy}\|}{2a_e},\\
\Term&=\big(\Term_{\mathrm{hover}}\cup\Term_{\mathrm{flee}}\big)\cap\big\{\min_\ell h_{c,\ell}(x_H;c_{\ell,H})\ge m_c\big\}.
\label{eq:terminal}
\end{align}
$p_{\mathrm{brake}}$는 수평 속도 방향으로 $a_e$로 제동할 때의 정지점이다. 따라서 종단이 항상 $\Base$인 것은 아니며, 도피 중인 상태도 제동 공간이 남아 있으면 허용한다.
회복 정책 $j$의 유한시간 회복 가능 집합은
\begin{equation}
C^{(j)}=\big\{(x,\zeta):\ x_i\in\Safe^{+}\ (i=1,\dots,H),\ x_H\in\Term\big\},\qquad C=\bigcup_{j=0}^{6}C^{(j)}
\end{equation}
이며($\Safe^+$는 식~\eqref{eq:fresh}의 여유를 더한 안전 집합), 후보 실행 직후 상태가 $C^{(j)}$에 속하는지를 검사한다.
합집합은 각 집합을 포함하므로 같은 검사 조건에서 한 회복 정책만 쓰는 경우보다 선택지가 줄지 않는다. 다만 모든 회피 가능 상태를 찾는다는 뜻은 아니다.

\subsection{선택 규칙}
인증된 조합의 집합을 $\mathcal A$라 할 때 명목 입력의 비중이 가장 큰 후보를 실행한다.
\begin{equation}
(\lambda^\star,j^\star)\in\arg\max_{(\lambda,j)\in\mathcal A}\lambda,\qquad u=u^{(j^\star)}_{\lambda^\star}\quad(\text{동률이면 작은 }j).
\label{eq:select}
\end{equation}
변형마다 회복 입력이 다르므로 이는 $\|u-\unom\|$의 최소화와 일반적으로 같지 않다. ``명목 입력이 회복 가능성을 유지하면 그대로 통과시키고, 필요할 때만 회복 쪽으로 움직인다''는 최소 수정의 이산 형태이다.
그림~\ref{fig:snapshot}은 실제 에피소드에서 필터가 한 스텝에 검사한 내용을 다시 계산한 것이다. 명목 입력($\lambda=1$)은 앞의 장애물 때문에 모든 회복 정책에서 거부되고, 회복 정책 0과 $\lambda=0.75$의 조합이 인증되어 실행되었다(평가 기록과 일치).
\IfFileExists{\snapdir/snapshot_ep845_step091.png}{%
\begin{figure}[H]\centering
\includegraphics[width=\linewidth]{\snapdir/snapshot_ep845_step091.png}
\caption{필터의 한 스텝 결정(시드 0, 에피소드 845, $t=1.8$\,s). 왼쪽: 일곱 회복 정책의 3\,s 계획($\lambda=0$)과 끝나는 방식(점선 원: 호버 타원체, $\times$: 제동 정지점), 선택한 후보의 계획(굵은 파랑)과 추격자 예측(초록 점선), 추격자의 $h_c=0$ 거리(옅은 원). 오른쪽: 후보 35개의 인증 결과(초록: 인증, 빨강: 궤적 조건 위반, 주황: 종단 조건 위반)와 선택한 칸.}\label{fig:snapshot}
\end{figure}}{}

\subsection{계획 유지와 추격자 예측 재생}\label{sec:commit}
학습 회복 정책의 입력(가까운 장애물 순서, 범위 안의 장애물)은 스텝마다 바뀔 수 있어 직전에 인증한 계획이 다음 스텝에 달라질 수 있다.
그래서 필터는 기억 $m=(\zeta_{\mathrm{view}},a_0,\text{호버 고정 여부},j_c,\text{추격자 슬롯 순서},\hat c)$를 저장한다. $\hat c$는 인증한 계획이 예측한 다음 스텝의 추격자 위치이다.
$\mathcal A=\varnothing$이면 저장한 변형 $j_c$의 회복 입력을 \emph{저장한 장면 입력과 예측 추격자 위치} $\hat c$로 계산해 실행한다.
모델이 정확하면 이 입력은 인증된 회복 궤적의 다음 입력과 정확히 같다(committed continuation). 이 입력도 실제 추격자 위치에 대해 여유 없이 다시 검사하고, 실패하면 그대로 실행하되 ``인증 없음''으로 기록한다.
예측 추격자 위치는 식~\eqref{eq:pursue}로 한 스텝씩 진행한다.

\subsection{알고리즘}
\begin{algorithm}[H]
\caption{이산 다중 회복 안전 필터 $F$ (한 제어 스텝)}\label{alg:filter}
\KwIn{상태 $x$, 장면(장애물, 추격자 $c$), 명목 입력 $\unom$}
\KwMem{$m=(\zeta_{\mathrm{view}},a_0,\text{고정 여부},j_c,\text{슬롯 순서},\hat c)$}
\KwOut{실행 입력 $u$, 갱신된 기억 $m'$, 정보 $(\lambda^\star,\text{인증 여부})$}
현재 장면에서 새 장면 입력 $\zeta_{\mathrm{new}}$를 만들고, 추격자 예측 $c_1\gets$ 식~\eqref{eq:pursue}를 한 스텝 적용\;
\For{$j=0,\dots,6$}{
  $\pi_b^{(j)}(x)$ 계산(식~\eqref{eq:reflex})\;
  \For{$\lambda\in\{1,0.75,0.5,0.25,0\}$}{
    $u^{(j)}_\lambda\gets\sat(\lambda\unom+(1-\lambda)\pi_b^{(j)}(x))$, \ $x_1\gets f_d(x,u^{(j)}_\lambda)$\;
    회복 정책 $j$를 $N=150$스텝 전개(식~\eqref{eq:rollout}--\eqref{eq:pursue})\;
    모든 노드가 식~\eqref{eq:fresh}를 만족하고 $x_H\in\Term$이면 $(\lambda,j)$를 $\mathcal A$에 추가\;
  }
}
\eIf{$\mathcal A\neq\varnothing$}{
  $(\lambda^\star,j^\star)\gets\arg\max_{\mathcal A}\lambda$ (동률이면 작은 $j$), \ $u\gets u^{(j^\star)}_{\lambda^\star}$\;
  $m'\gets(\zeta_{\mathrm{new}},a_0',\text{고정 여부}',j^\star,\text{새 슬롯 순서},c_1)$\;
}{
  $u\gets\pi_b^{(j_c)}(x;\zeta_{\mathrm{view}},\hat c)$ \tcp*{직전 인증 계획의 다음 입력}
  여유 없이 실제 추격자에 대해 재검사하고, 실패하면 ``인증 없음''으로 기록\;
  $m'\gets(\zeta_{\mathrm{view}},a_0',\text{고정 여부}',j_c,\text{슬롯 순서},\Pi(\hat c))$\;
}
\KwRet{$u,\ m'$}
\end{algorithm}

\subsection{보장되는 것과 보장되지 않는 것}\label{sec:guarantee}
\begin{keybox}{명제 (가정 아래의 진술)}
가정: (A1) 기체 모델이 정확하고 외란이 없다. (A2) 각 추격자의 속력은 $v_c$ 이하이다. (A3) 3\,s 예측 구간에서 속력 $v_c$의 순수 추적이 $h_c$에 대해 최악의 경우이다(거리에 대해서는 성립, 시선 방향에 대해서는 근사).

이때 실행되는 모든 입력은 인증된 후보이거나 직전 인증 계획의 이어짐이다. 인증된 계획은 3\,s 동안 상태를 여유를 둔 $\Safe$ 안에 두고 $\Term$에서 끝난다.
추격자 예측 재생 덕분에 이어짐은 인증된 회복 궤적의 꼬리를 정확히 재현하며, 실제 추격자 위치에 대해 다시 검사된다.
\end{keybox}
남는 간극은 $\Term$이 추격자에 대해 불변 집합이 아니라는 점이다. 호버는 결국 추격자에게 따라잡히고, 도피는 열린 공간이 남아 있는 동안만 안전하다.
또한 이어짐을 재검사할 때 예측 구간이 한 스텝 늘어나므로, 그 마지막 상태가 $\Term$을 만족하지 못해 인증을 잃을 수 있다.
따라서 이 필터는 3\,s 앞을 인증하고 20\,ms마다 다시 계획하는 장치이며, 이동 추격자에 대한 무한시간 불변성의 증명은 아니다.
실제 측정에서는 인증 상실이 PS2 정책 에피소드의 1.5\,\%에서 나타났다(표~\ref{tab:table2_final}).

\subsection{계산량}
한 제어 스텝에서 후보 35개와 이어짐 7개, 즉 42개의 150스텝 전개(제어 스텝당 약 6300회의 회복 입력 계산)를 수행하며, 학습 변형은 그때마다 98차원 입력의 신경망을 평가한다.
GPU(RTX 4050 노트북)에서 1000개 에피소드를 묶어 평가하면 15\,s 에피소드 전체가 약 350\,s 걸린다(묶음 기준 약 0.46\,ms/기체-스텝).
단일 기체의 실시간 최악 실행 시간은 측정하지 않았다.

%======================================================================
\section{비교 필터: bCBF-QP}\label{sec:bcbf}

PS2-RL의 backup CBF QP를 같은 집합에 적용한 비교 구성이다. 기체와 추격자를 포함한 상태를 $y=(p,v,q,c_1,\dots,c_C)$, 결합 동역학을 $\dot y=f(y)+g(y)u$,
고정된 회복 정책(반사 포함)의 흐름을 $\varphi(\tau,y)$, 그 벡터장을 $f_\pi$라 하자.
민감도 $Q_i=\partial\varphi(\tau_i,y)/\partial y$는 Euler 전개의 정확한 Jacobian으로 전파한다.
절대 예측 시각 $t_i$를 고정하면 $\tau_i=t_i-t$가 시간에 따라 줄어들므로, 매끄러운 구간에서
\begin{equation}
\frac{d}{dt}h\big(\varphi(t_i-t,y(t))\big)=\nabla h(\varphi_i)\big[Q_i\big(f(y)+g(y)u\big)-f_\pi(\varphi_i)\big]
\end{equation}
이다. 둘째 항이 남은 예측 시간이 줄어드는 데서 나온다. 이를 $-\alpha h(\varphi_i)$ 이상으로 제한한 선형 제약을 노드 $\tau_i$(1.5\,s 구간, 3스텝 간격)와 모든 안전 함수 $h\in\{h_o,h_f,h_w,h_z,h_c\}$에 대해 둔다($\alpha=4$).
종단 노드에서는 기준 집합 함수(반사가 작동 중이면 제동점 함수)와 $h_c$에 같은 형태의 제약을 $\alpha_B=2$로 둔다.
\begin{equation}
\min_{u,\ s\ge0}\ \tfrac12\sum_kw_k(u_k-u_{\mathrm{nom},k})^2+\tfrac12w_ss^2\quad\text{s.t. 완화된 CBF 제약, }u\in\mathcal U,\qquad w=(1,0.1,0.1,1),\ w_s=10^6 .
\end{equation}
roll/pitch 각속도의 가중치가 작아 추력보다 기울기로 경로를 바꾸는 해를 선호한다.
회복 입력 $u=\pi_b(x)$가 제약을 $10^{-2}$보다 크게 위반하면(정규화된 행 기준) QP를 건너뛰고, QP 출력은 다시 이산 필터(최종 회복 정책군)를 통과시킨다.
따라서 이 구성의 접촉 억제는 QP 단독의 결과가 아니며, 선형화·이산화·slack 때문에 표준 연속시간 CBF 불변성 정리를 그대로 적용할 수 없다.

%======================================================================
\section{임무 정책}\label{sec:striker}

\paragraph{기준 접근 제어} 표적 위치 $p_t$와 속도 $v_t$에 대해
\begin{equation}
v_{\mathrm{des}}=v_{\mathrm{cr}}\frac{p_t-p}{\|p_t-p\|}+v_t,\qquad v_{\mathrm{cr}}=6\ \mathrm{m/s}
\end{equation}
를 회피 반사와 같은 속도·자세 루프($k_v=2.5$, 기울기 $50^\circ$, $k_{\mathrm{att}}=10$, $k_\psi=3$)로 추종해 $\upur$를 만든다. 장애물과 추격자를 고려하지 않는다.

\paragraph{잔차 정책}
\begin{equation}
\unom=\sat\big(\upur(x)+\Delta u_\xi(o)\big),\qquad |\Delta a_T|\le0.1\,a_{\max}=3.92\ \mathrm{m/s^2},\quad |\Delta\omega_i|\le0.25\,\omega_{\max}=4.5\ \mathrm{rad/s}.
\label{eq:residual}
\end{equation}
관측은 135차원이다: 회복 정책 형식의 장면 관측 98개, 표적 오프셋/10, 표적 속도, 추격자 속도, 경기장 내 위치(14개),
유도 특징 19개(시선 방향 3, $d/10$, $\log(d+0.1)$, 상대 속도/5 3, 접근 속도/5, 시선각 변화율/2 3, 요 좌표계의 상대 위치/10 3과 상대 속도/5 3, 표적 높이/2),
그리고 보정 대상인 기준 입력 $(\upur-\ustar)/u_{\mathrm{scale}}$(4개).

\paragraph{보상} 표적 거리 $d_k=\|p_k-p_{t,k}\|$에 대해
\begin{equation}
r_k=w_p(d_{k-1}-d_k)-w_t-\frac{w_a}{4}\Big\|\frac{\Delta u_k}{u_{\mathrm{scale}}}\Big\|^2+100\cdot\mathbf 1_{\text{성공}}-150\cdot\mathbf 1_{\text{접촉 또는 이탈}},\qquad w_p=2,\ w_t=w_a=0.05 .
\label{eq:reward}
\end{equation}
거리 항은 $w_p(d_0-d_{\mathrm{end}})$로 합쳐지고, 시간 비용은 빠른 성공을 장려한다. 시간초과는 critic 목표에서 종료로 보지 않는다(부트스트랩 유지).

\paragraph{학습} SAC, $\gamma=0.995$, 보상 척도 0.1, 병렬 환경 64개, 배치 256, 재생 버퍼 $10^6$, 학습률 $3\times10^{-4}$, $\tau=0.005$, 온도 초기값 0.1(하한 $10^{-3}$, 목표 엔트로피 $-4$), 은닉층 256$\times$2(ELU).
처음 20만 스텝은 actor를 고정하고 critic만 학습한다(critic warm-up). actor 손실에 고정점 벌점 $0.2\,\mathbb E\big[\tfrac14\|\Delta u/b\|^2\big]$($b$는 잔차 한계)을 더한다.
1000만 스텝 학습하며, 25만 스텝마다 128개 에피소드로 평가해 반환값이 가장 큰 체크포인트를 쓴다. 처음부터 학습한 정책은 실패했고(성공 0, 이후 0.41--0.49), 기준 접근 제어의 잔차로 두어야 안정적으로 학습되었다.

%======================================================================
\section{PS2: 필터를 포함한 임무 학습}\label{sec:ps2}

필터를 환경의 일부로 두고 임무 정책을 다시 학습한다.
\begin{align}
u_{\mathrm{nom},k}&=\sat\big(\upur(x_k)+\Delta u_\xi(o_k)\big),\qquad (u_k,m_{k+1})=F\big(x_k,\zeta_k,m_k,u_{\mathrm{nom},k}\big),\qquad x_{k+1}=f_d(x_k,u_k), \label{eq:ps2}\\
r^{\mathrm{PS2}}_k&=r_k-w_f\|\eta_k\|^2,\qquad \eta_k=\frac{u_k-u_{\mathrm{nom},k}}{u_{\mathrm{scale}}},\qquad w_f=0.5 . \label{eq:ps2-reward}
\end{align}
정책은 자신이 제안한 입력이 그대로 실행된다고 가정하지 않고, 필터 개입 이후의 상태와 보상으로 학습한다. 필터가 자주 막는 행동보다 실제로 진전이 생기고 수정 비용이 작은 행동이 유리해진다.
필터 기억 $m$까지 포함한 전이 $(x,\zeta,m)\mapsto(x^+,\zeta^+,m^+)$는 정책에서 보면 환경의 일부이므로, 이산 필터를 미분하지 않고도 off-policy SAC로 학습할 수 있다.
다만 $m$은 관측되지 않으므로 관측이 완전한 상태라고 주장하지 않는다. 필터 출력을 복사하는 모방학습이 아니라, 수정된 전이에서 같은 강화학습 목표를 다시 최적화하는 것이다.

\paragraph{설정} 같은 시드의 잔차 정책 최고 체크포인트에서 시작하여, 처음 10만 스텝은 actor를 고정하고 critic을 재적응시키며, 온도는 시작 값으로 고정한다(학습률 0).
배치할 필터(최종 회복 정책군, 추격자 예측 재생)와 같은 필터로 300만 스텝 학습하고(학습 = 배치), 25만 스텝마다 64개 에피소드로 평가해 반환값이 가장 큰 체크포인트를 쓴다.
PS2 정책은 필터와 함께 쓰도록 학습한 정책이며, 필터 없이 실행하면 성공률이 잔차 정책보다 낮다(62.1\,\% 대 74.0\,\%). 필터가 장애물에서 비켜 주는 것에 의존하기 때문이다.

\paragraph{학습 시간} RTX 4050 노트북 GPU에서 시드당 Phase I 약 10분(1000만 스텝), 잔차 정책 28--34분(1000만 스텝), PS2 72--76분(300만 스텝)이 걸렸다.

%======================================================================
\section{평가 방법}\label{sec:protocol}

\paragraph{절차} 시드 $s\in\{0,1,2\}$마다 전체 파이프라인(Phase I $\to$ 잔차 정책 $\to$ PS2)을 그 시드로 학습한다.
모든 구성을 같은 1000개 평가 장면(장면 시드 9001, 학습 장면 풀과 별도)으로 평가하며, 장면마다 출발 교란과 표적·추격자 초기 조건까지 같다.
에피소드는 15\,s이고, 10\,s 결과는 같은 에피소드를 10\,s에서 절단한 것이다(10\,s까지 끝나지 않은 에피소드는 시간초과로 센다).
각 시드는 그 시드의 임무 정책을 학습한 추격자 설정(원래 추격자)으로 평가한다.

\paragraph{지표} 성공률, 접촉률(장애물·바닥·추격자), 시간초과율, 영역 이탈률, 성공 에피소드의 평균 도달 시간,
평균 명목 유지율 $\bar\lambda$(에피소드 스텝 평균 $\lambda^\star$, 새 인증이 없는 스텝은 0),
인증 상실률(한 스텝 이상 유효한 인증 계획이 없었던 에피소드의 비율). 인증 상실은 접촉과 다른 지표이다.

\paragraph{통계} 시드별 비율의 평균 $\pm$ 표준편차를 보고하고, 시드를 합친 횟수 $k/n$에 대해 양측 95\,\% Clopper--Pearson 구간을 쓴다.
\begin{equation}
p_L=\mathrm{B}^{-1}\!\big(\tfrac{\alpha}{2};\,k,\,n-k+1\big),\qquad p_U=\mathrm{B}^{-1}\!\big(1-\tfrac{\alpha}{2};\,k+1,\,n-k\big),\qquad \alpha=0.05,
\end{equation}
여기서 $\mathrm B^{-1}$는 베타 분포의 분위수이다. $k=0$이면 $p_U=1-(\alpha/2)^{1/n}$이고, 단측 95\,\% 상한은 $1-0.05^{1/n}\approx3/n$이다.
$n=1000$에서 0.30\,\%(단측)와 0.37\,\%(양측), $n=3000$에서 0.10\,\%와 0.12\,\%이다.
모든 구성이 같은 장면을 쓰므로 두 구성의 비교는 \emph{쌍 비교}이다. 구성 A만 성공한 장면 수 $b$와 B만 성공한 장면 수 $c$로 정확 McNemar 검정을 한다.
\begin{equation}
p_{\mathrm{McNemar}}=\min\Big\{1,\ 2\sum_{i=0}^{\min(b,c)}\binom{b+c}{i}2^{-(b+c)}\Big\}.
\end{equation}
구성끼리 비교하는 표(표~\ref{tab:table2_final}, \ref{tab:table3_ablation})는 모든 행이 평가된 공통 시드만 쓴다.

%======================================================================
\section{결과}\label{sec:results}

\begin{warnbox}{수치의 기준}
아래 표와 그림은 업로드된 평가 결과(시드 0 전체, 시드 1 일부, 시드 2 미평가)에서 만든 것이다. 각 행의 ``시드 $\times$ 장면'' 열이 사용한 에피소드 수를 나타낸다.
세 시드 평가를 마친 뒤 \file{STEPS=tables bash docker\_batch/run\_strike\_paper.sh}로 다시 만들고 이 문서를 다시 컴파일하면 표와 그림이 갱신된다.
\end{warnbox}

\subsection{주요 비교}
\tabin{table1_10s}
\tabin{table1_15s}

\begin{figure}[H]\centering
\figin{\figpath{fig_outcomes.pdf}}
\caption{구성별 결과 비율(왼쪽 10\,s, 오른쪽 15\,s 제한). 막대 안의 수는 성공률(\%)이다.}\label{fig:outcomes}
\end{figure}

\paragraph{필터의 효과} 같은 임무 정책에서 최종 필터는 접촉을 거의 없앤다. 잔차 정책은 26.1\,\%(521/2000)에서 0/2000으로, 기준 접근 제어는 57.5\,\%에서 1/2000으로 줄었다.
대신 우회와 정지 때문에 10\,s 시간초과가 늘어난다(잔차 정책 0\,\%$\to$15.5\,\%). 15\,s 제한에서는 시간초과 대부분이 늦은 성공으로 바뀐다.

\paragraph{PS2의 효과} 같은 최종 필터를 쓸 때 PS2 정책은 잔차 정책보다 성공률이 높다(10\,s 84.5$\to$90.4\,\%, 15\,s 94.4$\to$97.6\,\%).
같은 장면 쌍 비교에서 PS2만 성공한 장면이 172개, 잔차 정책만 성공한 장면이 54개로(10\,s) 차이는 유의하다($p<10^{-6}$, 표~\ref{tab:table5_paired}).
평균 도달 시간도 6.53\,s에서 6.11\,s로 줄었다. 필터를 고정하고 임무 정책만 바꾼 비교이므로 필터 포함 재학습의 효과를 보여 주는 핵심 근거이다.

\begin{figure}[H]\centering
\begin{subfigure}{0.49\linewidth}\centering\figin{\figpath{fig_strike_time.pdf}}\caption{누적 성공률}\label{fig:strike-time}\end{subfigure}\hfill
\begin{subfigure}{0.49\linewidth}\centering\figin{\figpath{fig_safety.pdf}}\caption{추격자 최소 여유 거리의 누적 분포}\label{fig:safety}\end{subfigure}
\caption{(a) 시간에 따른 누적 성공률(음영은 시드를 합친 95\,\% 구간). 점선은 10\,s. (b) 에피소드별 추격자 최소 여유 거리 $\min_t(\|p-c\|-r_c-r_d)$의 누적 분포. 필터를 쓴 구성에서 0.2\,m의 계단은 추격자가 바짝 따라오는 동안 필터가 $h_c$를 신규 계획 여유 $m_c=0.2$\,m에 유지한 에피소드들이다(멀어지는 중에는 $h_c$가 거리 여유와 같다).}
\end{figure}

\subsection{최종 필터 선택과 구성 요소}
\tabin{table2_final}
\tabin{table3_ablation}
\tabin{table5_paired}

\begin{figure}[H]\centering
\figin{\figpath{fig_ablation.pdf}}
\caption{구성 요소 제거 실험(PS2 정책, 15\,s). 오차 막대는 95\,\% Clopper--Pearson 구간이다.}\label{fig:ablation}
\end{figure}

\begin{itemize}
\item \textbf{회피 변형}: 단일 회피 필터는 접촉 1.7\,\%(17/1000)로, 최종 필터(1/1000)와의 쌍 비교에서 차이가 유의하다($p=3.1\times10^{-5}$). 성공률도 낮다(15\,s 95.7\,\% 대 98.4\,\%).
\item \textbf{bCBF-QP}: 성공률은 비슷하지만 최종 필터가 10\,s와 15\,s 모두 높고($p=0.0015$, $0.043$), 평균 도달 시간이 짧다(5.93\,s 대 6.33\,s, 공통 시드 기준). bCBF는 명목 유지율이 더 높지만(0.97) 이후 이산 검사에서 인증을 더 자주 잃는다(2.5\,\% 대 1.5\,\%).
\item \textbf{학습 회복 정책}: 회복 정책을 호버 LQR로 바꾸면 15\,s 성공률이 82--84\,\%로 크게 떨어지고 시간초과가 15--17\,\%로 늘어난다(시드 2개와 공통 시드 기준). Phase I 회복 정책이 필터의 허용 범위를 넓히는 것이 성능의 주요 원인이다.
\item \textbf{구형 근사}: 추격자를 정지한 구로 보면 접촉이 26.3\,\%로, 상대 접근 속도를 반영한 $h_c$가 필수적이다.
\item \textbf{계획 유지의 추격자 예측 재생과 호버 latch}: 1000개 장면에서 성공·접촉의 차이는 관측되지 않았다(98.6\,\%, 98.3\,\% 대 98.4\,\%). 두 요소는 성능보다 이어짐이 인증된 계획을 정확히 재현하고 인증 대상 제어기와 일치하게 하는 일관성을 위한 것이다.
\end{itemize}

\begin{figure}[H]\centering
\figin{\figpath{fig_filter.pdf}}
\caption{필터 개입. 왼쪽: 에피소드 평균 명목 유지율 $\bar\lambda$의 누적 분포(1에서의 계단은 필터가 한 번도 수정하지 않은 에피소드). 오른쪽: 인증 상실 에피소드 비율.}\label{fig:filter}
\end{figure}

\subsection{학습 과정과 Phase I}
\begin{figure}[H]\centering
\figin{\figpath{fig_training.pdf}}
\caption{학습 중 평가 성공률(시드별). 원은 평가에 사용한 체크포인트. 필터 없이 학습한 잔차 정책은 시드 1, 2에서 500--650만 스텝 이후 붕괴했으며, 평가에는 최고 반환값 체크포인트(시드 0: 850만, 시드 1: 125만, 시드 2: 175만 스텝)를 썼다. PS2 학습은 필터 안에서 안정적이다.}\label{fig:training}
\end{figure}
\tabin{table4_phase1}
\begin{figure}[H]\centering
\figin[0.55\linewidth]{\figpath{fig_phase1.pdf}}
\caption{Phase I 회복 성공률(1.5\,s 안에 안전하게 기준 집합 도달). 학습 회복 정책은 세 시드 평균과 표준편차.}\label{fig:phase1}
\end{figure}

\subsection{대표 에피소드}\label{sec:results-episodes}
\file{run\_strike\_paper.sh}는 시드 0의 평가에서 필터가 장애물 충돌을 막은 장면, 추격자 접촉을 막은 장면, 추격자 근접 회피, 유일한 접촉 장면,
기준 제어가 실패한 장면 등을 자동으로 고르고(\file{episode\_picks.json}), 각 장면을 여러 구성으로 다시 실행해 다음을 만든다.
\begin{itemize}
\item 2차원 궤적 비교(구성별 패널), 3차원 궤적(속력 색), 3차원 비교 그림;
\item 상세 신호 그림: 표적 거리, 장애물·벽·바닥·천장 여유, $h_c$와 반사 가중치 $\sigma$, 속력과 고도, $\lambda^\star$와 인증 상태(신규/이어짐/인증 없음), 사용한 회복 변형, 입력과 명목 입력;
\item 2차원 GIF와 회전하는 3차원 GIF(추격자 $h_c$ 요구 거리 구 표시);
\item 필터 스냅샷(\file{scripts/strike\_filter\_snapshot.py}): 처음 명목 입력을 줄인 스텝, 이전 계획을 이어 간 스텝, 인증을 잃은 스텝에서 필터의 검사를 다시 계산한 그림(그림~\ref{fig:snapshot}).
\end{itemize}
\IfFileExists{\episodedir/overview_main.png}{%
\begin{figure}[H]\centering
\includegraphics[width=\linewidth]{\episodedir/overview_main.png}
\caption{대표 에피소드(시드 0)의 수평 궤적. 실선은 최종 필터를 쓴 PS2 정책과 잔차 정책, 점선은 필터 없는 PS2 정책이다. 별은 성공, X는 접촉, 회색 원은 시간초과.}\label{fig:episodes-overview}
\end{figure}}{\IfFileExists{\episodedir/overview.png}{%
\begin{figure}[H]\centering
\includegraphics[width=\linewidth]{\episodedir/overview.png}
\caption{대표 에피소드(시드 0)의 수평 궤적.}\label{fig:episodes-overview}
\end{figure}}{}}

\paragraph{유일한 접촉 사례(시드 0, 에피소드 845)} 기체는 큰 원기둥과 위쪽 벽 사이의 좁은 공간($y\approx9$\,m)으로 들어간 뒤 약 6\,s부터 새 인증을 찾지 못해
이전 계획(호버 LQR + 회피 변형)을 이어 갔고, 약 7.5\,s부터는 그 계획도 재검사를 통과하지 못했다(인증 없음).
추격자가 다가오면서 회피 반사가 작동했지만 벽과 장애물이 도피 방향을 막아 9.2\,s에 접촉했다. 같은 장면에서 bCBF 구성도 같은 방식으로 실패했다.
이 사례는 \S\ref{sec:guarantee}의 남는 간극, 즉 호버 종단이 다가오는 추격자에 대해 불변 집합이 아니고 막힌 공간에서는 도피도 보장되지 않는다는 점이 실제로 나타난 경우이다.
인증을 잃은 7.5\,s의 검사를 다시 계산하면(그림~\ref{fig:snapshot-lost}) 35개 후보 모두 궤적 조건을 위반했다. 위쪽 벽($y=10$\,m) 쪽으로 움직이던 기체에서 일곱 회복 정책의 3\,s 계획이 모두 벽을 넘었기 때문이다.
Phase I 회복 정책은 벽 없이 학습되었고, 회피 반사는 추격자 여유가 $h_{\mathrm{off}}=2$\,m 이상이면 섞이지 않는다. 벽을 포함한 Phase I 학습이 남은 개선 방향이다.
\IfFileExists{\snapdir/snapshot_ep845_step374.png}{%
\begin{figure}[H]\centering
\includegraphics[width=\linewidth]{\snapdir/snapshot_ep845_step374.png}
\caption{인증을 잃은 시점의 검사(에피소드 845, $t=7.5$\,s). 일곱 회복 정책의 계획이 모두 위쪽 벽을 넘어 새 인증이 없고, 필터는 이전 계획을 이어 실행했다.}\label{fig:snapshot-lost}
\end{figure}}{}
\IfFileExists{\episodedir/ep845/compare2d.png}{%
\begin{figure}[H]\centering
\includegraphics[width=\linewidth]{\episodedir/ep845/compare2d.png}
\caption{에피소드 845를 여러 구성으로 다시 실행한 수평 궤적(CPU 재실행).}\label{fig:ep845-2d}
\end{figure}}{}
\IfFileExists{\episodedir/ep845/detail_ps2_ours.png}{%
\begin{figure}[H]\centering
\includegraphics[width=0.86\linewidth]{\episodedir/ep845/detail_ps2_ours.png}
\caption{에피소드 845, PS2 정책 + 최종 필터의 상세 신호. 위에서부터 표적 거리, 기하 여유, $h_c$와 반사 가중치 $\sigma$, 속력과 고도, $\lambda^\star$와 인증 상태(노랑: 이전 계획 유지, 빨강: 인증 없음), 사용한 회복 변형, 입력과 명목 입력.}\label{fig:ep845-detail}
\end{figure}}{}

%======================================================================
\section{한계와 적용 범위}\label{sec:limits}
\begin{enumerate}
\item \textbf{종단 집합의 불변성}: 호버 기준 집합은 정적 장애물에 대한 회복 수단이다. 추격자가 계속 다가오면 호버는 영구히 안전하지 않고, 도피 종단도 유한시간 예측의 조건이다. 정적 기준 집합에 기대는 원래 이론의 보장을 그대로 옮길 수 없다.
\item \textbf{추격자 모델}: 예측은 현재 시선 방향의 순수 추적이며, 장애물 우회와 시선각 변화를 포함한 모든 추격 전략의 최악값이라는 증명은 없다. 반응 지연과 이탈 가속도에 기반한 $h_c$도 근사이다.
\item \textbf{모델 오차}: 입력 지연, 공력, 센서 오차, 바람이 있으면 3\,s 검사와 실제 경로가 달라질 수 있다. 모든 오차를 포함하는 강건 경계는 제시하지 않았다.
\item \textbf{시간 이산 검사}: 검사는 20\,ms 노드에서 한다. 노드 사이 연속시간 충돌까지 배제하려면 이동량이나 보간 오차를 이용한 추가 조건이 필요하다.
\item \textbf{인증 상실}: 계획 유지 중에도 인증을 잃을 수 있다(PS2 에피소드의 1.5\,\%). 이 시점은 명제의 전제를 벗어난다.
\item \textbf{통계}: 1000개 장면과 2--3개 시드의 결과이다. 접촉 1/2000은 접촉확률이 0이라는 증명이 아니며 95\,\% 상한은 0.28\,\%이다.
\item \textbf{실기체}: 결과는 시뮬레이션의 성공과 접촉이다. 실시간 최악 실행 시간, PX4 추력·각속도 추종 오차, 실기체 운용 성능은 검증하지 않았다.
\end{enumerate}

%======================================================================
\appendix
\section{하이퍼파라미터}\label{app:hparams}
\begin{center}\small
\begin{tabularx}{\linewidth}{>{\raggedright\arraybackslash}p{3.0cm}>{\raggedright\arraybackslash}X}\toprule
구분 & 값 \\\midrule
모델 & $\Delta t=0.02$\,s, $a_T\in[0,4g]$, $|\omega_i|\le18$\,rad/s, $r_d=0.3$\,m, $z_f=1$\,m, $z_{\max}=10$\,m \\
경기장·장면 & 40$\times$20\,m, 장애물 3개, 추격자 1개($r_c=0.4$\,m, $v_c=2$\,m/s, $a_c=2$\,m/s$^2$, 반발 거리 2\,m), 표적 0.2--0.8\,m/s, 성공 반경 0.6\,m, 15\,s \\
$h_c$ & $\delta=0.2$\,s, $a_e=5$\,m/s$^2$ \\
호버 LQR & $Q=\mathrm{diag}(2,2,8,0.01,0.01,0.5,8,8,2)$, $R_u=\mathrm{diag}(0.06,0.15,0.15,0.12)$, $c_B=50$, $\rho_{xy}=1.22$\,m, $z_{\mathrm{ext}}=0.63$\,m \\
Phase I & TD3, 1000만 스텝, 병렬 64, 배치 256, 버퍼 $10^6$, $\beta=0.99$, 학습률 actor $10^{-4}$/critic $3\times10^{-4}$, $\tau=0.0025$, 정책 지연 2, 탐색 잡음 0.1, 잔차 벌점 0.2, 에피소드 112스텝, 구형 근사 $T=1.5$\,s, 영역 가중치 $(0.4,0.2,0.2,0.2)$ \\
회피 반사 & $h_{\mathrm{on}}=0.5$\,m, $h_{\mathrm{off}}=2.0$\,m, $(k_v,k_{\mathrm{att}},k_\psi)=(2,10,3)$, 기울기 $45^\circ$, $d_{\mathrm{push}}=1.5$\,m, $z_{\mathrm{ref}}\in[z_f+1,z_{\max}-3]$ \\
회복 정책군 & $(v_{\mathrm{esc}}/v_c,\vartheta)\in\{(2,0^\circ)\text{ 학습},(2,0^\circ)\text{ LQR},(2,\pm60^\circ),(3.5,0^\circ),(3.5,\pm60^\circ)\}$ \\
이산 필터 & $\Lambda=\{1,0.75,0.5,0.25,0\}$, 예측 3\,s(150스텝), 신규 계획 여유 $m_o=0.15$\,m·$m_c=0.2$\,m, 종단 $\Term_{\mathrm{hover}}\cup\Term_{\mathrm{flee}}$, 계획 유지(추격자 예측 재생), 호버 latch \\
bCBF-QP & $\alpha=4$, $\alpha_B=2$, 1.5\,s, 3스텝 간격, $w=(1,0.1,0.1,1)$, $w_s=10^6$, 건너뛰기 허용치 $10^{-2}$, 보조 필터 = 최종 이산 필터 \\
기준 접근 제어 & $v_{\mathrm{cr}}=6$\,m/s, $k_v=2.5$, 기울기 $50^\circ$ \\
잔차 정책 & 잔차 한계 추력 10\,\%, 각속도 25\,\%; 관측 135; SAC $\gamma=0.995$, 보상 척도 0.1, 병렬 64, 배치 256, 학습률 $3\times10^{-4}$, $\tau=0.005$, 온도 0.1(하한 $10^{-3}$), critic warm-up 20만, 고정점 벌점 0.2, 1000만 스텝 \\
보상 & $w_p=2$, $w_t=0.05$, $w_a=0.05$, 성공 $+100$, 접촉·이탈 $-150$ \\
PS2 & $w_f=0.5$, 잔차 정책에서 시작, critic 재적응 10만, 온도 고정, 300만 스텝 \\
평가 & 시드 3개, 장면 1000개(시드 9001), 15\,s와 10\,s 절단, Clopper--Pearson 95\,\%, 정확 McNemar \\\bottomrule
\end{tabularx}
\end{center}

\section{기호표}\label{app:symbols}
\begin{center}\small
\begin{tabularx}{\linewidth}{>{\raggedright\arraybackslash}p{3.4cm}X}\toprule
기호 & 의미 \\\midrule
$x=(p,v,q)$, $u=(a_T,\omega)$ & 기체 상태, 입력(추력 가속도, 기체 각속도) \\
$f_d$, $\Delta t$ & 이산 동역학, 시간 간격 \\
$\zeta$, $c_\ell$, $\nu_\ell$ & 장면 정보, 추격자 위치, 추격자 속도 \\
$r_d$, $r_c$, $v_c$, $a_c$ & 기체 반경, 추격자 반경, 추격자 속력·가속도 상한 \\
$h_f,h_o,h_w,h_z$ & 바닥, 장애물, 벽, 천장 여유 \\
$h_{c,j}$, $s_j^+$, $\delta$, $a_e$ & 추격자 여유도, 양의 접근 속도, 반응 지연, 유효 이탈 가속도 \\
$\Safe$, $\Base$, $\BaseP$ & 안전 집합, 필터의 기준 집합, Phase I 기준 집합 \\
$e$, $P$, $K$, $V$ & 호버 오차, Riccati 해, LQR 이득, $V=e^\top Pe$ \\
$a^\star$, $D$, $S$ & 최적 정지점, 정지점 행렬, Schur 여행렬 ($V_{\min}=r^\top Sr$) \\
$c_B$, $\rho_{xy}$, $z_{\mathrm{ext}}$ & 기준 집합 수준, 타원체의 수평·수직 범위 \\
$\pi_\theta$, $\Delta u_\theta$ & Phase I 학습 회복 정책과 그 잔차 \\
$\pi_b^{(j)}$, $\uesc^{(j)}$, $\sigma$ & 회복 정책 변형 $j$, 회피 입력, 반사 가중치 \\
$\Lambda$, $\lambda^\star$, $\bar\lambda$ & 혼합률 집합, 선택한 혼합률, 에피소드 평균 \\
$N$, $H$, $m_o$, $m_c$ & 예측 스텝(150), 종단 첨자($N+1$), 신규 계획 여유 \\
$\Term$, $p_{\mathrm{brake}}$ & 종단 집합, 제동 정지점 \\
$F$, $m$ & 안전 필터, 필터 기억(계획 유지) \\
$\upur$, $\Delta u_\xi$, $\unom$ & 기준 접근 입력, 임무 정책 잔차, 명목 입력 \\
$r_k$, $r^{\mathrm{PS2}}_k$, $\eta_k$, $w_f$ & 보상, PS2 보상, 정규화한 필터 수정량, 수정 비용 가중치 \\\bottomrule
\end{tabularx}
\end{center}

\section{원고 초안과 다른 점}\label{app:draft}
초안을 이 버전의 결과로 바꿀 때 확인할 항목이다.
\begin{enumerate}
\item \textbf{평가 규모}: 초안의 ``단일 시드, 128개 장면'' 수치(44.5, 82.8, 93.8, 100\,\% 등)는 이전 실험이다. 현재 표는 1000개 장면, 시드 2--3개이며 15\,s 에피소드와 10\,s 절단을 함께 보고한다.
\item \textbf{접촉 0건 표현}: PS2 정책 + 최종 필터는 접촉 1/2000(시드 0의 한 장면)이다. ``접촉이 관측되지 않았다'' 대신 ``2000회 중 1회(95\,\% 상한 0.28\,\%)''로 쓴다. 잔차 정책 + 최종 필터는 0/2000이다.
\item \textbf{15\,s 성공률}: 초안의 ``128개 장면 모두 성공(100\,\%)''은 현재 97.6\,\%이다(시간초과 2.4\,\%).
\item \textbf{PS2 정책의 학습 필터}: 초안 표 1은 단일 회피 필터로 학습한 PS2 정책을 썼다. 최종 파이프라인의 PS2 정책은 배치하는 최종 필터(회피 변형, 추격자 예측 재생)로 학습했다.
\item \textbf{bCBF slack 가중치}: 초안 식 (34) 뒤의 $w_s=10$은 $w_s=10^6$이다.
\item \textbf{신규 계획 여유}: 0.15\,m는 장애물과 벽뿐 아니라 바닥과 천장에도 적용된다.
\item \textbf{기준 집합의 추격자 조건}: 필터의 기준 집합 $\Base$에는 추격자 조건이 없고, 추격자는 종단 조건 $h_c\ge0.2$\,m와 회피 반사로 다룬다. Phase I의 기준 집합만 추격자 구 조건을 가진다.
\item \textbf{천장 조건}: 식 (20)의 $z_{\max}-a^\star_z\ge\rho_{\mathrm{top}}+r_d$에서 구현은 $\rho_{\mathrm{top}}=\rho_{xy}$이다(초안과 같음, 보수적).
\item \textbf{제동 정지점}: 식 (30)은 구현에서 수평 성분만 쓴다: $p_{\mathrm{brake}}=p_{H,xy}+v_{H,xy}\|v_{H,xy}\|/(2a_e)$. 이 점이 장애물과 벽에서 $r_d$ 이상 떨어져야 한다.
\item \textbf{추격자 운동}: 환경의 추격자는 속력 2\,m/s와 함께 가속도 2\,m/s$^2$ 상한을 가지며, 장애물 2\,m 이내에서 밀려난다(\S\ref{sec:arena}). 필터의 예측(가속도 제한 없는 순수 추적)과 다르다는 점을 밝혀 둔다.
\item \textbf{Phase I 수치}: 가중 성공률은 세 시드 평균 70.1$\pm$0.3\,\%(초안 69.9\,\%), LQR 43.8\,\%이다.
\item \textbf{이동 추격자 회복 학습 변형}: 초안 표 2의 해당 행은 최종 파이프라인 평가에 포함되지 않은 128개 장면 실험이다. 쓰려면 이전 실험임을 밝히거나 표에서 뺀다.
\item \textbf{수식 표기}: 초안 식 (2)·(6)의 $|\omega_i)$, $|\ell_i)$는 절댓값 $|\omega_i|$, $|\ell_i|$이고, 식 (4)의 괄호는 $v_{k+1}=v_k+\Delta t\,(a_{T,k}R(q_k)e_3-ge_3)$이다. 회피 각 $\vartheta$와 임무 정책 매개변수(초안 $\Delta u_\vartheta$)의 기호가 겹치므로 이 문서는 임무 정책 잔차를 $\Delta u_\xi$로 썼다.
\end{enumerate}

\section{한글(HWP) 수식 편집기 입력}\label{app:hwp}
본문의 주요 수식을 한글 수식 편집기 문법으로 옮긴 것이다(같은 내용: \file{results/strike\_paper/hwp\_equations.txt}). 붙여넣은 뒤 모양을 확인한다.
\IfFileExists{hwp_equations.txt}{\VerbatimInput[fontsize=\scriptsize,breaklines,breakanywhere,firstline=4]{hwp_equations.txt}}{\fbox{hwp\_equations.txt 없음}}

\end{document}
